 

%\documentclass[english,aps,prx,twocolumn,nofootinbib,notitlepage,longbibliography]{revtex4-1}
%\usepackage[T1]{fontenc}
%\usepackage[latin9]{inputenc}
%\setcounter{secnumdepth}{3}
%\usepackage{array}
%\usepackage{multirow}
%\usepackage{amsmath}
%\usepackage{amssymb}
%%-------------------Below is revTex setting-----------------------
\documentclass[11pt,nofootinbib,preprint]{revtex4-1}
\pdfoutput=1
\usepackage[T1]{fontenc}
\usepackage[latin9]{inputenc}
\synctex=-1
%\usepackage[sort&compress]{natbib}
\usepackage{times}
\usepackage{graphicx}
\usepackage{tabularx}
\usepackage{graphicx}
\usepackage{dcolumn}
\usepackage{bm}
\usepackage{booktabs}
\usepackage{makecell}
\usepackage{makebox}
\usepackage{multirow}
%\usepackage{natbib}
\usepackage[figuresright]{rotating}
\usepackage{amsmath}
\usepackage{amstext}
\usepackage{amssymb}
\usepackage{xfrac}
\usepackage[colorlinks=true,citecolor={blue},linktocpage=true]{hyperref}\usepackage{graphicx}

\usepackage{amsmath}
\usepackage{amstext}
\usepackage{amssymb}
\usepackage{amsfonts}
\usepackage{longtable,booktabs}
\usepackage{hyperref}
\usepackage{url}
\usepackage{subfigure}
\usepackage{dsfont}
\usepackage{booktabs}
\usepackage{amsbsy}
%\usepackage{dcolumn}
\usepackage{amsthm}
\usepackage{bm}
\usepackage{esint}
\usepackage{multirow}
\usepackage{hyperref}
\usepackage{cleveref}
\usepackage{mathrsfs}
\usepackage{amsfonts}
\usepackage{amsbsy}
\usepackage{dcolumn}
\usepackage{bm}
\usepackage{multirow}
\usepackage{color}
\usepackage{extarrows}
\usepackage{datetime}
\usepackage{comment}
\usepackage[super]{nth}
\usepackage{tikz-cd}
\usepackage{float}

\makeatletter

%%%%%%%%%%%%%%%%%%%%%%%%%%%%%% LyX specific LaTeX commands.
%% Because html converters don't know tabularnewline
\providecommand{\tabularnewline}{\\}

%%%%%%%%%%%%%%%%%%%%%%%%%%%%%% User specified LaTeX commands.
%%\usepackage{jheppub} % for details on the use of the package, please
                     % see the JHEP-author-manual

%%%\usepackage[T1]{fontenc} % if needed


\def\Z{\mathbb{Z}}
\newcommand{\red}[1]{{\textcolor{red}{#1}}}
\def\epsi{\epsilon^{\mu\nu\mathscr{D}} }
\def\ep{\epsilon^{\mu\nu\mathscr{D}\rho} }

\makeatother


\begin{document}
\title{Fusion rules and shrinking rules of topological orders in five dimensions}
%\title{Non-Abelian Fusion Rules and Hierarchical Shrinking Rules in 5D Topological Orders}
\author{Yizhou Huang}
\thanks{These authors contributed equally to this work.}
 \author{Zhi-Feng Zhang}
\thanks{These authors contributed equally to this work.}
 \author{Peng Ye}
\email{yepeng5@mail.sysu.edu.cn}
\affiliation{School of Physics, Sun Yat-sen University, Guangzhou, 510275, China}
\date{\today}
\begin{abstract}
As a series of work about 5D (spacetime) topological orders, here we employ the path-integral formalism of 5D  topological quantum field theory  (TQFT) established in \href{https://doi.org/10.1007/JHEP04(2022)138}{Zhang and Ye, JHEP 04 (2022) 138}  to explore non-Abelian fusion rules,  hierarchical shrinking rules and quantum dimensions of particle-like, loop-like and membrane-like topological excitations in 5D topological orders.  To illustrate, we focus on a prototypical example of twisted $BF$ theories that comprise the twisted topological terms of the $BBA$ type. First, we classify topological excitations by establishing equivalence classes among all gauge-invariant Wilson operators. Then, we compute fusion rules from the path-integral and find that fusion rules may be of non-Abelian nature, i.e., the fusion outcome can be a direct sum of distinct excitations. We further compute shrinking rules. Especially, we discover exotic hierarchical structures hidden in shrinking processes of 5D or higher:  a membrane is shrunk into particles and loops, and the loops are subsequently shrunk into a direct sum of particles. We obtain the algebraic structure of shrinking coefficients and fusion coefficients.  We compute the quantum dimensions of all excitations and find that sphere-like membranes and torus-like membranes differ not only by their shapes but also by their quantum dimensions. We further study the algebraic structure that determines anomaly-free conditions on fusion coefficients and shrinking coefficients. Besides $BBA$, we explore general properties of all twisted terms in $5$D. Together with braiding statistics reported before, the theoretical progress here paves the way toward characterizing and classifying topological orders in higher dimensions where topological excitations consist of both particles and spatially extended objects. 
\end{abstract}
\maketitle
\newpage
\tableofcontents{}

\section{Introduction\label{s1}}
Originally discovered in the context of strongly correlated electron systems, such as the fractional quantum Hall effect and chiral spin liquids, \textit{topological order} has been widely explored in the literature of quantum many-body physics, high energy physics, mathematical physics, and quantum information~\cite{wenZootopoRMP}. While it is well-known that Ginzburg-Landau order parameters cannot characterize topological orders and their phase transitions, for the past decades, topological data, such as the number of topological excitations (i.e., anyons), ground state degeneracy, edge chiral central charge, braiding statistics, and fusion rules of anyons, have been extensively studied and utilized to characterize  3D (spacetime) topological orders.  As a matter of fact, the practical goal of characterizing and classifying topological orders can be regarded as the establishment of a set of such topological data or physical observables for every type of topological order, assuming that the ground states are realized by Hamiltonians with short-range interactions and a bulk gap. To achieve this goal, various theoretical approaches have been developed, including but not limited to, the construction of exactly solvable models (e.g., string-net models~\cite{string1}, Dijkgraaf-Witten models~\cite{dwitten}, and stabilizer codes~\cite{Kitaev2003faulttolerant}), quantum-informative perspectives (e.g.,  long-range entanglement nature of topological order~\cite{CGW2010LRE},   the correspondence between zero modes of entanglement spectrum and the physical edge state~\cite{Li_Haldane}, and topological entanglement entropy~\cite{KitaevTopoEE,Levin_Wen_TEE}), general algebraic structure based on unitary modular category~\cite{Kitaev2006,string1,Kong2020}, and topological quantum  field theoretical (TQFT) analysis~\cite{Witten1989,Turaev2016}.  


As one of the most prominent   TQFTs, the Chern-Simons theory governs the low-energy physics of $3$D topological orders and provides a quantitative tool for computing topological properties of anyons which are a key ingredient of topological quantum computation~\cite{sarma_08_TQC}. Meanwhile, by placing the Chern-Simons theory on an open disc, the 2D boundary theory of 3D topological orders is shown to be dominated by gapless modes and mathematically described by conformal field theory (CFT)~\cite{moore1989classical} with  gravitational anomaly~\cite{Wen_grav_prd2013,string8}.  Furthermore, the Gutzwiller projective construction of both  bulk ground state wavefunctions and the associated edge CFT algebraic construction can be performed in a systematic manner~\cite{Wen99,Barkeshli2010,barkeshli_wen}.  It should be noticed that the Chern-Simons theory mentioned in the present paper belongs to  hydrodynamical field-theoretical approach, which is different from the flux-attachment   approach of composite bosons and fermions~\cite{composite_Jain_2,composite_Fradkin,PhysRevLett.62.82}. The Chern-Simons theory of 3D topological orders with global symmetry can be applied to characterizing and classifying 3D symmetry-enriched topological phases (SET) in which anyons may carry a fractional quantum number of global symmetry group~\cite{PhysRevB.93.155121,PhysRevB.87.195103}. Recently, the Chern-Simons theory at the trivial level~\cite{PhysRevB.86.125119,YW12,PhysRevB.93.115136,Ye14b}  has been   applied to describe the bulk effective field theory of bosonic symmetry-protected ``topological'' phases (SPT) as long as global symmetry is imposed nontrivially. Furthermore, the Chern-Simons theory at the trivial level, after weakly gauging, leads to   topological response theory of bosonic SPTs~\cite{PhysRevLett.112.141602,Hung_Wen_gauge}.  

 
While the notion of topological order was born in $3$D,  lifting spatial dimensions of topological orders to $4$D and higher is an active branch of frontier research. In higher dimensions, the self-statistics of particle excitations is no longer anyonic~\cite{wu84} but the presence of spatially extended excitations (e.g., loops in $4$D and higher and membranes in $5$D and higher) may significantly enrich the topological data. Specifically, in $4$D topological orders, topological excitations include both  point-like particle excitations and loop-like excitations.  Consider a discrete Abelian gauge group $G=\prod_{i=1}^n{\mathbb{Z} _{N_i}}$. 
The underlying TQFTs are called ``twisted $BF$ theories'' where the topological action consists of  a multi-component $BF$ term \cite{Horowitz1990,hansson2004superconductors}  and some ``twisted (topological) terms''. The $BF$ term in $4$D is written as a wedge product of $B$ and $dA$: $BdA$  ($B$ and $A$ are respectively $2$- and $1$-form gauge fields), which can be physically understood as a result of exotic boson condensation and duality  transformations~\cite{hansson2004superconductors,YW13a,ye16a,Moy_Fradkin2023}. Twisted terms (or ``twists'' for short) consist of  $AAdA$, $AAAA$~\cite{2016arXiv161209298P,PhysRevB.99.235137,YeGu2015,ypdw}, and $AAB$~\cite{yp18prl}.  From these topological terms, braiding statistics among topological excitations  can be systematically computed, such as 
   particle-loop braiding~\cite{hansson2004superconductors,abeffect,PRESKILL199050,PhysRevLett.62.1071,PhysRevLett.62.1221,ALFORD1992251}, multi-loop braiding~\cite{wang_levin1}, and particle-loop-loop braiding (i.e., Borromean-Rings braiding) \cite{yp18prl}, which are tightly connected to link invariants, e.g., Hopf linking number of Hopf links and Milner's triple linking number of general Brunnian links~\cite{kauffman2006formal,mellor2003geometric}.  By excluding all actions that are not gauge-invariant, ref.~\cite{zhang2021compatible} exhausted all legitimate combinations of these topological terms to obtain a complete list of topological orders in $4$D.  
   
   Recently, fusion rules, shrinking rules, and quantum dimensions of all topological excitations in $4$D  twisted $BF$ theories are computed systematically~\cite{PhysRevB.107.165117}, where shrinking rules of loops are studied concretely and  non-Abelian fusion rules are found despite the Abelian nature of the gauge group.   By \textit{shrinking}, we mean that spatially extended topological excitations, i.e., loops, are adiabatically shrunk into particles, which is a unique topological property for topological orders in $4$D and higher.  
   In addition, twisted $BF$ theories  of $4$D topological orders with various types of  symmetry can be applied to $4$D SET phases, leading to a field-theoretical classification and characterization of a large class of $4$D SET phases, see, e.g., refs.~\cite{ye16_set,2016arXiv161008645Y,Ye:2017aa,ye16a} and ref.~\cite{Ning2018prb} on the general theory. Again, twisted $BF$ theories at the trivial level have been applied to describe the bulk effective field theory of $4$D SPT phases~\cite{YeGu2015,bti2}, which, by weakly gauging, leads to the topological response theory of $4$D SPT phases~\cite{PhysRevLett.114.031601,PhysRevB.99.205120,bti6}. In addition, $BB$ and $dAdA$ type terms were also discussed respectively for fermionic topological phases and topological response theory of bosonic topological insulators~\cite{Kapustin:2014gua,PhysRevB.99.235137,ye2023unp,bti2,YW13a,lapa17}.
   
   
   By lifting spatial dimensions further, topological orders in $5$D admit particle-like, loop-like, and membrane-like topological excitations, which exhibit highly unexplored features of both the formalism of TQFTs and braiding statistics~\cite{Zhang:2021ycl}.   If we still consider the gauge group $G=\prod_{i=1}^n\Z_{N_i}$,  TQFTs in $5$D may contain two types of (generalized) $BF$ terms, i.e., $CdA$ and $\tilde{B}dB$, where $C$ is $3$-form, $B$ and $\tilde{B}$ are two different $2$-form, $A$ is $1$-form.  Therefore, for each $\Z_{N_i}$ subgroup, there are two choices for the assignment of gauge charges and corresponding $BF$ terms, which significantly enriches topological orders. On top of two types of $BF$ terms, there are many different twisted topological terms, e.g., $AAAAA$, $AAAdA$, $AdAdA$, $AAC$, $AAAB$, $BBA$, $AdAB$, $AAdB$, which lead to many different twisted $BF$ theories after proper combinations in $5$D,  as long as gauge invariance is guaranteed~\cite{zhang2021compatible}.  In ref.~\cite{Zhang:2021ycl}, braiding statistics encoded in all topological excitations has been investigated  via computing correlation functions of Wilson operators whose   expectation values  rely on  counting intersections of sub-manifolds in $5$D. The physical theory presented in ref.~\cite{Zhang:2021ycl}  provides an alternative route to understanding the link theory of higher-dimensional compact manifolds.  

While braiding statistics has shown how diverse the topological data of $5$D topological orders are~\cite{Zhang:2021ycl}, the fusion rules and shrinking rules of $5$D topological orders are still yet to be uncovered. Since  there are  membrane topological excitations that are spatially extended excitations of two dimensions, it is natural to expect more exotic fusion and shrinking rules than that of $4$D studied in ref.~\cite{PhysRevB.107.165117}, which motivates the present paper. In this paper, we adopt the path-integral formalism of $5$D topological orders constructed in ref.~\cite{Zhang:2021ycl}. We focus   on two prototypical twisted topological terms, namely, $BBA$ and $AAAB$ both of which exhibit not only non-Abelian fusion rules but also   hierarchical structure in shrinking rules.  By \textit{hierarchical}, we mean that a membrane is  shrunk into  particles and loops, and  the loops are subsequently shrunk into a direct sum of particles. Such phenomena can only exist in $5$D and higher.  We  realize that membrane excitations potentially possess very different topological properties, e.g., quantum dimensions if the membranes have different topologies, e.g., sphere and torus. We concretely compute all fusion rules and verify the consistency between fusion rules and shrinking rules, which leads to the algebraic structure hidden in the fusion coefficients and shrinking coefficients. Any violation of the algebraic structure may be associated with a quantum anomaly of some kind in the topological order. 
   
     This paper is organized as follows. In section~\ref{s2}, we first review the twisted $BF$ theory with the $BBA$ twisted term introduced in ref.~\cite{Zhang:2021ycl}. Then we construct Wilson operators for all topological excitations (see table~\ref{tab_excitation}) and classify them into different equivalence classes. In section~\ref{s3}, we calculate the fusion rules (see table~\ref{table_fusionBBA}) for the $BBA$ twisted term explicitly from the path integral. We also compute quantum dimensions (see table~\ref{table_quantumdBBA}) for all topological excitations. In section~\ref{s4}, we calculate the shrinking rules (see table~\ref{table_shrinking}) for the $BBA$ twisted term. We concretely discuss the hierarchical structure in the shrinking rules. In section~\ref{section_algebraic}, we study the general algebraic structure that determines anomaly-free topological orders. In section~\ref{section_general_discuss_twisted5D}, we present several general properties (see table~\ref{table_comparison_5dtwisted}) of all twisted terms in $5$D. 
A brief summary is given in section~\ref{section_summary}. Several appendices are provided.

 
\section{Topological excitations and  Wilson operators for the $BBA$ twisted term\label{s2}}

\subsection{Twisted $BF$ theory with the $BBA$ twisted term}
Any gauge group $G=\prod_{i=1}^n{\mathbb{Z} _{N_i}}$ with $n\geqslant3$ can support $BBA$ term in TQFT action. For simplicity, we consider $G=\left(\mathbb{Z} _{2}\right)^{3}$ and take the TQFT action to be~\cite{Zhang:2021ycl}
\begin{equation}
	S=\int{\frac{N_1}{2\pi}\tilde{B}^1dB^1+}\frac{N_2}{2\pi}\tilde{B}^2dB^2+\frac{N_3}{2\pi}C^3dA^3+qB^1B^2A^3\,,
	\label{eq_action_BBA}
\end{equation}
where $q$ is a periodic and quantized coefficient. $A^3$ is a 1-form gauge field. $B^1$, $B^2$, $\tilde{B}^1$, and $\tilde{B}^2$ are 2-form gauge fields.  $C^3$ is a 3-form gauge field.  We note that there are two kinds of $BF$ terms in the action~(\ref{eq_action_BBA}), i.e., $BdB$ and $CdA$, so the action represents a mixed $BF$ theory following the definition in ref.~\cite{Zhang:2021ycl}.  The coefficient $q$ is quantized and periodically identified: $q=\frac{pN_1N_2N_3}{\left( 2\pi \right) ^2N_{123}}$, $p\in \mathbb{Z} _{N_{123}}$, where $N_{123}$ is the greatest common divisor of $N_{1}$, $N_{2}$ and $N_{3}$, and  $N_{1}=N_{2}=N_{3}=N_{123}=2$ for $G=\left(\mathbb{Z} _{2}\right)^{3}$ considered here.  The gauge fields $\tilde{B}^1$, $\tilde{B}^2$ and $C^3$ serve as the Lagrange multipliers in action~(\ref{eq_action_BBA}). $B^1$, $B^2$ and $A^3$ satisfy the flat connection conditions: $dB^1=0$, $dB^2=0$ and $dA^3=0$. The gauge transformations are given by:
\begin{align}
	&B^1\rightarrow B^1+dV^1, \tilde{B}^1\rightarrow \tilde{B}^1+d\tilde{V}^1+\frac{2\pi q}{N_1}\left( V^2A^3+B^2\chi ^3+V^2d\chi ^3 \right)\,,   
	\\
	&B^2\rightarrow B^2+dV^2, \tilde{B}^2\rightarrow \tilde{B}^2+d\tilde{V}^2+\frac{2\pi q}{N_2}\left( V^1A^3+B^1\chi ^3+V^1d\chi ^3 \right)\,,
	\\
	&A^3\rightarrow A^3+d\chi ^3, C^3\rightarrow C^3+dT^3+\frac{2\pi q}{N_3}\left( -V^1B^2-B^1V^2-\frac{1}{2}V^1dV^2-\frac{1}{2}V^2dV^1 \right) \,,\label{eq_gauge_C3}
\end{align}
where $\chi^3$, $V^i$, $\tilde{V}^i$ and $T^3$ are $0$-form, $1$-form, $1$-form and $2$-form gauge parameters respectively. They satisfy the following compactness conditions: $\int d\chi^3\in2\pi\mathbb{Z}$, $\int dV^i\in2\pi\mathbb{Z}$, $\int d\tilde{V}^i\in2\pi\mathbb{Z}$ and $\int dT^3\in2\pi\mathbb{Z}$, where the integrals are performed over compact submanifolds.

\subsection{Wilson operators for topological excitations}
Topological excitations in $5$D spacetime consist of point-like particle excitations, one-dimensional closed loop excitations, and two-dimensional closed membrane excitations. We will simply call them particles, loops, and membranes unless otherwise specified. Since loops and membranes are spatially extended excitations, they can have different shapes in space such as self-linking loops, knotted loops, or membranes in the shape of a torus with multiple handles. We   consider some simplest shapes, which are $S^1$ loop, $S^2$ sphere and $T^2$ torus (with genus$=1$), as shown in figure~\ref{fig_excitation}. 
\begin{figure}
	\centering
	\includegraphics[scale=0.9,keepaspectratio]{fig_excitation}
	\caption{Illustration of particle excitations and spatially extended topological excitations with different shapes in the $5$D twisted $BF$ theory with the $BBA$ twisted term.}
	\label{fig_excitation}
\end{figure}


In TQFTs, topological excitations  are represented by Wilson operators. In other words, by respectively using gauge fields of $1$-, $2$-, and $3$-form, we can construct Wilson loop operators for particle excitations, Wilson surface operators for loop excitations, and  Wilson volume operators for membrane excitations. Thus,   the equivalence classes of Wilson operators directly render the equivalence classes of topological excitations, which eventually leads to a finite set of topological excitations despite the infinite number of gauge-invariant Wilson operators that can be written.  Below, we will demonstrate the details of some examples of  Wilson operators, and leave all nonequivalent Wilson operators in table~\ref{tab_excitation}. In this table, there are $14$ topologically distinct excitations and one of them is the topologically trivial excitation denoted as $\mathsf{1}$.  It is important to point out that trivial particles, trivial loops, and trivial membranes are all identical to each other and are eventually denoted as $\mathsf{1}$. 
\begin{table}
	\caption{\label{tab_excitation}Operators for nonequivalent excitations in $5$D topological order with action $S=\int{\frac{N_1}{2\pi}\tilde{B}^1dB^1+}\frac{N_2}{2\pi}\tilde{B}^2dB^2+\frac{C_1}{2\pi}C^3dA^3+qB^1B^2A^3$. }
	\centering
	\begin{tabular*}{\textwidth}{ccc}
		\toprule
		\makebox[0.15\textwidth][c]{excitation} & \makebox[0.55\textwidth][c]{operator for excitation} & \makebox[0.3\textwidth][c]{equivalent excitations} \\
		\midrule
		$\mathsf{1}$ & $\exp \left( 0 \right)$ & $-$ \\
		\midrule
		$\mathsf{P}_{001}$ & $\exp \left( i\int_{\gamma}{A^3} \right)$ & $-$ \\
		\midrule
		$\mathsf{L}_{100}$ & $\exp \left( i\int_{\sigma}{B^1} \right)$ & $-$ \\
		\midrule
		$\mathsf{L}_{010}$ & $\exp \left( i\int_{\sigma}{B^2} \right)$ & $-$ \\
		\midrule
		$\mathsf{L}_{110}$ & $\exp \left( i\int_{\sigma}{B^1}+i\int_{\sigma}{B^2} \right)$ & $-$ \\
		\midrule
		$\mathsf{L}_{100}^{001}$ & $\exp \left( i\int_{\sigma}{B^1}+i\int_{\gamma}{A^3} \right)$ & $-$ \\
		\midrule
		$\mathsf{L}_{010}^{001}$ & $\exp \left( i\int_{\sigma}{B^2}+i\int_{\gamma}{A^3} \right)$ & $-$ \\
		\midrule
		$\mathsf{L}_{110}^{001}$ & $\exp \left( i\int_{\sigma}{B^1}+i\int_{\sigma}{B^2}+i\int_{\gamma}{A^3} \right)$ & $-$ \\
		\midrule
		$\mathsf{M^S}_{001}$ & \thead{$2\exp \left[ i\int_{\omega}{C^3+\frac{1}{2}\frac{2\pi q}{N_3}\left( d^{-1}B^1B^2+d^{-1}B^2B^1 \right)} \right]$ \\ $\times\delta \left( \int_{\sigma}{B^1} \right) \delta \left( \int_{\sigma}{B^2} \right)$} & \thead{$\mathsf{M^S}_{001}=\mathsf{M^S}_{001}^{,ij0}$ \\  $i,j=0,1$} \\
		\midrule
		$\mathsf{M^S}_{001}^{001,}$ & \thead{$2\exp \left[ i\int_{\gamma}{A^3}+i\int_{\omega}{C^3+\frac{1}{2}\frac{2\pi q}{N_3}\left( d^{-1}B^1B^2+d^{-1}B^2B^1 \right)} \right]$ \\ $\times\delta \left( \int_{\sigma}{B^1} \right) \delta \left( \int_{\sigma}{B^2} \right)$} & \thead{$\mathsf{M^S}_{001}^{001,}=\mathsf{M^S}_{001}^{001,ij0}$ \\ $i,j=0,1$} \\
		\midrule
		$\mathsf{M^T}_{001}$ & \thead{$4\exp \left[ i\int_{\tilde{\omega}}{C^3+\frac{1}{2}\frac{2\pi q}{N_3}\left( d^{-1}B^1B^2+d^{-1}B^2B^1 \right)} \right]$ \\ $\times\delta \left( \int_{\sigma}{B^1} \right) \delta \left( \int_{\sigma}{B^2} \right)$} & \thead{$\mathsf{M^T}_{001}=\mathsf{M^T}_{001}^{,ij0}$ \\ $i,j=0,1$} \\
		\midrule
		$\mathsf{M^T}_{001}^{001,}$ & \thead{$4\exp \left[ i\int_{\gamma}{A^3}+i\int_{\tilde{\omega}}{C^3+\frac{1}{2}\frac{2\pi q}{N_3}\left( d^{-1}B^1B^2+d^{-1}B^2B^1 \right)} \right]$ \\ $\times\delta \left( \int_{\sigma}{B^1} \right) \delta \left( \int_{\sigma}{B^2} \right)$} & \thead{$\mathsf{M^T}_{001}^{001,}=\mathsf{M^T}_{001}^{001,ij0}$ \\ $i,j=0,1$} \\
		\midrule
		$\mathsf{M^{ST}}$ & \thead{$8\exp \left[ i\int_{\omega}{C^3+\frac{1}{2}\frac{2\pi q}{N_3}\left( d^{-1}B^1B^2+d^{-1}B^2B^1 \right)} \right]$ \\$\times\exp \left[ i\int_{\tilde{\omega}}{C^3+\frac{1}{2}\frac{2\pi q}{N_3}\left( d^{-1}B^1B^2+d^{-1}B^2B^1 \right)} \right]$\\ $\times\delta \left( \int_{\sigma}{B^1} \right) \delta \left( \int_{\sigma}{B^2} \right)$} & \thead{$\mathsf{M^{ST}}=\mathsf{M^{ST}}^{,ij0}$ \\ $i,j=0,1$} \\
		\midrule
		$\mathsf{M^{ST}}^{001,}$ & \thead{$8\exp \left[ i\int_{\gamma}{A^3}+i\int_{\omega}{C^3+\frac{1}{2}\frac{2\pi q}{N_3}\left( d^{-1}B^1B^2+d^{-1}B^2B^1 \right)} \right]$ \\$\times\exp \left[ i\int_{\tilde{\omega}}{C^3+\frac{1}{2}\frac{2\pi q}{N_3}\left( d^{-1}B^1B^2+d^{-1}B^2B^1 \right)} \right]$\\ $\times\delta \left( \int_{\sigma}{B^1} \right) \delta \left( \int_{\sigma}{B^2} \right)$} & \thead{$\mathsf{M^{ST}}^{001,}=\mathsf{M^{ST}}^{001,ij0}$ \\ $i,j=0,1$} \\
		\bottomrule
	\end{tabular*}
\end{table}



\textbf{\textit{Wilson operators for particles.}} 
 A particle carrying  unit gauge charge of $1$-form gauge field $A^3$ is called a \textit{$\mathbb{Z} _{N_3}$ particle}, which can be represented by the following gauge invariant Wilson operator:
\begin{align}
	\mathsf{P}_{001}=\exp \left( i\int_{\gamma}{A^3} \right) \,.
\end{align}
Here we use $\mathsf{P}_{n_{1}n_{2}n_{3}}$ to represent a particle simultaneously carrying $n_{1}$, $n_{2}$, $n_{3}$ units of $\mathbb{Z} _{N_1}$, $\mathbb{Z} _{N_2}$, $\mathbb{Z} _{N_3}$ gauge charge. $\gamma$ is the closed world-line of the particle. Since  the action~(\ref{eq_action_BBA}) contains only one $1$-form gauge field, i.e., $A^3$, there is only one kind of particle with unit gauge charge, which is $\mathbb{Z} _{N_3}$ particle. The antiparticle of $\mathsf{P}_{n_{1}n_{2}n_{3}}$ is represented by $\bar{\mathsf{P}}_{n_{1}n_{2}n_{3}}$. For instance, the antiparticle of $\mathsf{P}_{001}$ is represented by $\bar{\mathsf{P}}_{001}$. 
In section~\ref{s3}, we will show that for $G=\left(\mathbb{Z} _{2}\right)^{3}$, the antiparticle of any excitation is itself because the fusion of two identical excitations can be written as $c\cdot\left(\mathsf{1}\oplus\cdots\right)$, where $c$ is just a constant. Also, $n_{i}$ only takes $0$ or $1$ mod $2$ due to the $\mathbb{Z} _{2}$ cyclic group structure, which largely simplifies our analysis.

\textbf{\textit{Wilson operators for loops.}}
Similarly,  a $\mathbb{Z} _{N_1}$ loop carries unit gauge charge of the $2$-form gauge field $B^1$, and can be represented by
\begin{align}
	\mathsf{L}_{100}=\exp \left( i\int_{\sigma}{B^1} \right)\,,
\end{align}
where $\sigma$ is the closed world sheet of the loop excitation.  Such a loop is a \textit{pure loop} that only carries gauge charges of  either one or more $2$-form gauge fields.  In contrast, we can  further  consider a \textit{decorated loop} that simultaneously carries gauge charges of both $1$-form gauge fields and $2$-form gauge fields. For instance, a $\left( \mathbb{Z} _{N_1},\mathbb{Z} _{N_2} \right)$ loop decorated by a $\mathbb{Z} _{N_3}$ particle is represented by
\begin{align}
	\mathsf{L}_{110}^{001}=\exp \left( i\int_{\sigma}{B^1}+i\int_{\sigma}{B^2}+i\int_{\gamma}{A^3} \right)\,.
\end{align}
Here, the superscript ``$001$'' denotes a $\mathbb{Z} _{N_3}$ particle decoration with $\gamma\in \sigma$, and the subscript ``$110$'' means that the loop simultaneously carries unit gauge charges of $B^1$ gauge field and  $B^2$ gauge field. If there is no decoration at all, $\mathsf{L}_{110}^{000}$ is a pure loop and can be rewritten as $\mathsf{L}_{110}$ for the notational simplicity.


\textbf{\textit{Wilson operators for membranes.}}
As mentioned above, we will consider   membrane excitations of two different shapes: the membrane denoted by $\mathsf{M^S}$ is in the shape of $S^2$ and the membrane denoted by $\mathsf{M^T}$ is in the shape of $T^2$. For instance, one may write a Wilson (volume) operator as $
	\mathsf{M^S}_{001}=\exp \left( i\int_{\omega}{C^3} \right)$,  where the world volume $\omega$ is an $S^2\times S^1$ manifold in $5$D spacetime. However, the Wilson operator we construct above is not gauge invariant because $C^3$ is transformed with nontrivial shifts as shown in eq.~(\ref{eq_gauge_C3}). The correct form should be:
\begin{align}
	\mathsf{M^S}_{001}=N\exp \left[ i\int_{\omega}{C^3+\frac{1}{2}\frac{2\pi q}{N_3}\left( d^{-1}B^1B^2+d^{-1}B^2B^1 \right)} \right] \delta \left( \int_{\sigma}{B^1} \right) \delta \left( \int_{\sigma}{B^2} \right) \,,
	\label{eq3}
\end{align}
where $N$ is a normalization factor and in section~\ref{s3} we will show that $N=2$. We define $d^{-1}B^1$ and $d^{-1}B^2$ as
$	d^{-1}B^1=\int_{\mathcal{A} \in \sigma}{B^1} \quad\text{and}\quad d^{-1}B^2=\int_{\mathcal{A} \in \sigma}{B^2}
 $ ($\mathcal{A}$ is an open area on $\sigma$) whose legitimacy relies on  extra constraints:  $\int_{\sigma}{B^1}=0$ mod $2\pi$, $\int_{\sigma}{B^1}=0$ mod $2\pi$. To enforce these constraints, we introduce two Dirac's delta-functional in eq.~(\ref{eq3}):
\begin{align}
	\delta \left( \int_{\sigma}{B^1} \right) =\begin{cases}
		1, \quad \int_{\sigma}{B^1}=0   \,\, \mathrm{mod}\,\, 2\pi\\
		0, \quad \mathrm{else}\\
	\end{cases}\,,\,\,
	\delta \left( \int_{\sigma}{B^2} \right) =\begin{cases}
		1, \quad \int_{\sigma}{B^2}=0 \,\,  \mathrm{mod} \,\, 2\pi\\
		0, \quad \mathrm{else}\\
	\end{cases}\,.
\end{align}
Similarly, the operator for  a $\mathbb{Z} _{N_3}$ membrane of $T^2$ shape can be written as:
\begin{equation}
	\mathsf{M^T}_{001}=\tilde{N}\exp \left[ i\int_{\tilde{\omega}}{C^3+\frac{1}{2}\frac{2\pi q}{N_3}\left( d^{-1}B^1B^2+d^{-1}B^2B^1 \right)} \right] \delta \left( \int_{\sigma}{B^1} \right) \delta \left( \int_{\sigma}{B^2} \right) \,,
\end{equation}
where $\tilde{\omega}$ is a $T^2\times S^1$ manifold in $5$D spacetime, and in section~\ref{s4} we will show that  $\tilde{N}=4$. Symbolically, once we write down the Wilson operator for an $S^2$ membrane, we can formally express  the Wilson operator for a $T^2$ membrane  simply through replacing $\mathsf{M^S}$, $N$ and $\omega$ by $\mathsf{M^T}$, $\tilde{N}$ and $\tilde{\omega}$. 

Just like loop excitations that can be decorated by particles, membranes can   be decorated by particles and loops simultaneously. For instance, a $\mathbb{Z} _{N_3}$ membrane decorated by a $\mathbb{Z} _{N_3}$ particle and a $\mathbb{Z} _{N_1}$ loop is
\begin{multline}
	\mathsf{M^S}_{001}^{001,100}=2\exp \left[ i\int_{\gamma}{A^3}+i\int_{\sigma}{B^1}+i\int_{\omega}{C^3+\frac{1}{2}\frac{2\pi q}{N_3}\left( d^{-1}B^1B^2+d^{-1}B^2B^1 \right)} \right]\\ \times\delta \left( \int_{\sigma}{B^1} \right) \delta \left( \int_{\sigma}{B^2} \right) \,.
	\label{eq2}
\end{multline}
The number  $001$ and $100$ in the superscript respectively denote a $\mathbb{Z} _{N_3}$ particle  decoration and a $\mathbb{Z} _{N_1}$ loop decoration respectively. The subscript $001$ means that the membrane carries a unit gauge charge of $3$-form gauge field $C^3$. $\delta \left( \int_{\sigma}{B^1} \right)$ in eq.~(\ref{eq2}) plays an important role because it induces an equivalence relation between $\mathsf{M^S}_{001}^{001,100}$ and $\mathsf{M^S}_{001}^{001,000}$, which will be  discussed in section~\ref{ss23}. If there is no particle decoration at all, $\mathsf{M^S}_{001}^{000,100}$ can be  rewritten as $\mathsf{M^S}_{001}^{,100}$ for the notational simplicity.  If there is no loop decoration at all, $\mathsf{M^S}_{001}^{001,000}$ can be  rewritten as $\mathsf{M^S}_{001}^{001,}$ for the notational simplicity. If all decorations are absent, the pure membrane $\mathsf{M^S}_{001}^{000,000}$ is directly replaced by $\mathsf{M^S}_{001}$.




\textbf{\textit{Wilson operators for $\mathsf{M^{ST}}$-excitations.}}
We have already introduced two kinds of membranes: $S^2$ membranes and $T^2$ membranes. For the completeness of fusion rules to be discussed shortly, it is vital to incorporate another type of membrane formed by fusing an $S^2$ membrane and a $T^2$ membrane. The resulting object is no longer a manifold.  The corresponding Wilson operator is denoted as $\mathsf{M^{ST}}$:
\begin{align}
	\mathsf{M^{ST}}=&8\exp \left[ i\int_{\omega}{C^3+\frac{1}{2}\frac{2\pi q}{N_3}\left( d^{-1}B^1B^2+d^{-1}B^2B^1 \right)} \right] \nonumber
	\\
	&\times \exp \left[ i\int_{\tilde{\omega}}{C^3+\frac{1}{2}\frac{2\pi q}{N_3}\left( d^{-1}B^1B^2+d^{-1}B^2B^1 \right)} \right] \nonumber
	\\
	&\times \delta \left( \int_{\sigma}{B^1} \right) \delta \left( \int_{\sigma}{B^2} \right)\,.
\end{align}
 

\subsection{Equivalence classes of Wilson operators\label{ss23}}
In the above discussion, we mentioned that a   $\mathbb{Z} _{N_3}$ membrane that is decorated by a $\mathbb{Z} _{N_3}$ particle and a $\mathbb{Z} _{N_1}$ loop can be represented by eq.~(\ref{eq2}), and it is equivalent to a   $\mathbb{Z} _{N_3}$ membrane  that is only decorated by a $\mathbb{Z} _{N_3}$ particle,  because of the constraint $\delta \left( \int_{\sigma}{B^1} \right)$: 
\begin{multline}
\!\!	\!\!\!\!\mathsf{M^S}_{001}^{001,000}\!\!=\!\!2\exp \left[ i\int_{\gamma}{A^3}+i\int_{\omega}{C^3+\frac{1}{2}\frac{2\pi q}{N_3}\!\left( d^{-1}B^1B^2\!+\!d^{-1}B^2B^1 \right)} \right] \!\!\delta \left( \int_{\sigma}{B^1} \right) \!\!\delta \left( \int_{\sigma}{B^2} \right) .\!\!
\end{multline}
For simplicity, we omit $000$ and use $\mathsf{M^S}_{001}^{001,}$ to represent $\mathsf{M^S}_{001}^{001,000}$. We say $\mathsf{M^S}_{001}^{001,100}$ and $\mathsf{M^S}_{001}^{001,}$ are equivalent to each other because they always give the same answers when we calculate any gauge-invariant correlation functions in the path integral. More concretely, the correlation function of $\mathsf{M^S}_{001}^{001,100}$ and any operator $\mathcal{O}$ can be written as
\begin{align}
	\langle \mathcal{O} \mathsf{M^S}_{001}^{001,100} \rangle =\frac{1}{\mathcal{Z}}\int\mathcal{D} \left[ ABC \right] \exp \left( iS \right) \times \mathcal{O} \times \mathsf{M^S}_{001}^{001,100}\,,
\end{align}
where $\mathcal{Z} =\int\mathcal{D} \left[ABC\right] \exp \left( iS \right) $ is the partition function and $\mathcal{D}[ABC]$ denotes   configurations of all gauge fields $A^3,B^1,\tilde{B}^1,B^2,\tilde{B}^2,C^3$. The constraint $\delta \left( \int_{\sigma}{B^1} \right) $ means that in the functional integration, the integration over $B^1$ is constrained in the subspace in which $B^1$ is flat such that $\exp \left( i\int_{\sigma}{{B}^1} \right)=1 $. Immediately, one can realize that $\langle \mathcal{O} \mathsf{M^S}_{001}^{001,100} \rangle=\langle \mathcal{O} \mathsf{M^S}_{001}^{001,} \rangle$. 
% We use $\hat{\mathcal{O}}$ to represent the configuration of $\mathcal{O}$ after integrating out $\tilde{B}^1$, $\tilde{B}^2$ and $C^3$. Notice that $\delta \left( \int_{\sigma}{\hat{B}^1} \right) $ enforces $\exp \left( i\int_{\sigma}{\hat{B}^1} \right)=1 $, thus we have [\textbf{The following derivation should be edited!}]:
%\begin{align}
%	\langle \mathcal{O} \mathsf{M^S}_{001}^{001,100} \rangle =&\frac{1}{\mathcal{Z}}\int{\mathcal{D} \left[ A^i \right]}\mathcal{D} \left[ B^i \right] \mathcal{D} \left[ C^i \right] \exp \left( iS \right) \times \mathcal{O} \times \mathsf{M^S}_{001}^{001,100}\nonumber
%	\\
%	=&\hat{\mathcal{O}}\times 2\exp \left[ i\int_{\gamma}{\hat{A}^3}+i\int_{\sigma}{\hat{B}^1}+i\int_{\omega}{\hat{C}^3+\frac{1}{2}\frac{2\pi q}{N_3}\left( d^{-1}\hat{B}^1\hat{B}^2+d^{-1}\hat{B}^2\hat{B}^1 \right)} \right]\nonumber
%	\\
%	&\times\delta \left( \int_{\sigma}{\hat{B}^1} \right) \delta \left( \int_{\sigma}{\hat{B}^2} \right) \nonumber
%	\\
%	=&\hat{\mathcal{O}}\times 2\exp \left[ i\int_{\gamma}{\hat{A}^3}+i\int_{\omega}{\hat{C}^3+\frac{1}{2}\frac{2\pi q}{N_3}\left( d^{-1}\hat{B}^1\hat{B}^2+d^{-1}\hat{B}^2\hat{B}^1 \right)} \right] \nonumber
%	\\
%	&\times\delta \left( \int_{\sigma}{\hat{B}^1} \right) \delta \left( \int_{\sigma}{\hat{B}^2} \right) \nonumber
%	\\
%	=&\langle \mathcal{O} \mathsf{M^S}_{001}^{001,} \rangle 
%\end{align}
Since $\mathcal{O}$ is an arbitrary operator, $\mathsf{M^S}_{001}^{001,100}$ and $\mathsf{M^S}_{001}^{001,}$ are indistinguishable in all gauge-invariant observables. In this sense, we conclude that these two operators as well as the associated two topological excitations  are said to  belong to the same equivalence class. Similarly, noticing that $\delta \left( \int_{\sigma}{{B}^2} \right) $ enforces $\exp \left( i\int_{\sigma}{{B}^2} \right)=1$ in the functional integration,   we can obtain another equivalence class between  $\mathsf{M^S}_{001}^{001,010}$ and $\mathsf{M^S}_{001}^{001,110}$. In conclusion, we have
$	\mathsf{M^S}_{001}^{001,}=\mathsf{M^S}_{001}^{001,100}=\mathsf{M^S}_{001}^{001,010}=\mathsf{M^S}_{001}^{001,110}\,$,
where ``$=$'' simply means that they are indistinguishable in all topological data and thus belong to the same equivalence class. 
By exhausting all such equivalence classes,  we obtain $14$ topologically distinct excitations belonging to $14$ topologically distinct equivalence classes, which  are listed in table~\ref{tab_excitation}.


\section{Fusion rules for the $BBA$ twisted term\label{s3}}
\subsection{Fusion rules and the path-integral representation}
Fusion rules are important properties in topological orders. Several prototypical fusion processes are pictorially illustrated in figure~\ref{fig_fusion} and detailed calculations of these fusion rules are shown in section~\ref{ss31}.
\begin{figure}
	\centering
	\includegraphics[scale=0.9,keepaspectratio]{fig_fusion}
	\caption{Some typical fusion rules for the $5$D twisted $BF$ theory with the $BBA$ twisted term. Fusing two $\mathsf{P}_{001}$'s leads to a  vacuum (trivial excitation).  Fusing a $\mathsf{L}_{100}$ and a $\mathsf{P}_{001}$ leads to a  decorated loop $\mathsf{L}_{100}^{001}$. Fusing two $\mathsf{M^S}_{001}$'s  leads to a direct sum of a  vacuum and three different loops denoted by $\mathsf{L}_{100}$, $\mathsf{L}_{010}$ and $\mathsf{L}_{110}$. The first two fusion rules shown in this figure are Abelian fusion rules since each of them only contains a single fusion channel. The  last fusion rule is non-Abelian fusion  since it has multiple fusion channels.}
	\label{fig_fusion}
\end{figure}
A general fusion process can be symbolically written as:
\begin{align}
	\mathsf{a}\otimes \mathsf{b}=\oplus _{\mathsf{c}}N_{\mathsf{c}}^{\mathsf{a}\mathsf{b}}\mathsf{c}
	\label{eq_general_fusion}\,,
\end{align}
where $\mathsf{a}$, $\mathsf{b}$ and $\mathsf{c}$ denote topological excitations. $N_{\mathsf{c}}^{\mathsf{a}\mathsf{b}}\in\Z$ is  fusion coefficient.  In the condensed matter literature, a fusion process is said to be Abelian if there is only one $\mathsf{c}$ such that $N_{\mathsf{c}}^{\mathsf{a}\mathsf{b}}$ is nonzero, e.g., $\mathsf{a}\otimes \mathsf{b}=\mathsf{c}$. Otherwise, the fusion process is said to be non-Abelian, e.g., $\mathsf{a}\otimes \mathsf{b}=\mathsf{c}_1\oplus\mathsf{c}_2$. We follow this nomenclature in this paper.  For an excitation $\mathsf{a}$, if $\mathsf{a}\otimes \mathsf{b}$ is always Abelian for any $\mathsf{b}$, we call $\mathsf{a}$ an Abelian excitation. Otherwise, we call $\mathsf{a}$ a non-Abelian excitation.
In this paper, we use Wilson operators to represent topological excitations, the fusion process (\ref{eq_general_fusion}) can be interpreted as the equivalence of expectation values: 
\begin{align}
	\langle \mathcal{O}_\mathsf{a}\otimes \mathcal{O}_\mathsf{b} \rangle=\langle \oplus _{\mathsf{c}}N_{\mathsf{c}}^{\mathsf{a}\mathsf{b}}\mathcal{O}_\mathsf{c} \rangle\,.
\end{align}
In the path-integral representation, the above identity is represented by:   
\begin{align}
	 &\frac{1}{\mathcal{Z}}\int\mathcal{D} \left[ ABC\right] \exp \left( iS \right) \times \left( \mathcal{O} _{\mathsf{a}}\times \mathcal{O} _{\mathsf{b}} \right) =\frac{1}{\mathcal{Z}}\int \mathcal{D} \left[ ABC\right] \exp \left( iS \right) \times \left( \sum_\mathsf{c}{N_{\mathsf{c}}^{\mathsf{a}\mathsf{b}}}\mathcal{O} _{\mathsf{c}} \right)  \,.
	 	\label{eq_fusion_decomposition}
\end{align}
If the  decomposition (\ref{eq_fusion_decomposition}) is successfully performed,  all fusion coefficients can be obtained.  From the above formula, we can also verify that: 
\begin{align}
	&\mathsf{a}\otimes \mathsf{b}=\mathsf{b}\otimes \mathsf{a}\,, \left(\mathsf{a}\otimes \mathsf{b}\right)\otimes \mathsf{c}=\mathsf{a}\otimes\left(\mathsf{b}\otimes \mathsf{c}\right)\,,\mathsf{a}\oplus \mathsf{b}=\mathsf{b}\oplus \mathsf{c}\,,\left(\mathsf{a}\oplus \mathsf{b}\right)\oplus \mathsf{c}=\mathsf{a}\oplus\left(\mathsf{b}\oplus \mathsf{c}\right)\,.
	\label{eq_fusion_algebra_rule}
\end{align} 
\subsection{Examples\label{ss31}}
In the following, through     several examples, we will  concretely compute the fusion rules of the twisted $BF$ theory with the $BBA$ twisted term whose action is given in eq.~(\ref{eq_action_BBA}).  All fusion rules for the $BBA$ twisted term are summarized   in table~\ref{table_fusionBBA}.




\textbf{\textit{Fusing a $\mathbb{Z} _{N_3}$ particle and a $\mathbb{Z} _{N_3}$ particle.}}
Using eq.~(\ref{eq_fusion_decomposition}), we can calculate that
\begin{align}
	\langle \mathsf{P}_{001}\otimes \mathsf{P}_{001} \rangle 
	=&\frac{1}{\mathcal{Z}}\int{\mathcal{D} \left[ ABC   \right] \exp \left( iS \right) \exp \left( i\int_{\gamma}{A^3}+i\int_{\gamma}{A^3} \right) } \nonumber
	\\
	=&\frac{1}{\mathcal{Z}}\int{\mathcal{D} \left[ ABC   \right]  \exp \left( iS \right) \exp \left( i2\int_{\gamma}{A^3} \right) }  \,.
\end{align}
Integrating out Lagrange's multipliers, we obtain constraints for $A^3$, $B^1$, and $B^2$ respectively:
$\int{A^3}=\frac{2\pi m_3}{N_3}  \,,
	\int{B^1}=\frac{2\pi m_1}{N_1}  \,,
	\int{B^2}=\frac{2\pi m_2}{N_2} $, 
where $m_{1,2,3}\in\mathbb{Z}$. Since $N_1=N_2=N_3=2$, we have $\exp \left( i2\int_{\gamma}{A^3} \right) =1$, thus
 $	\langle\mathsf{P}_{001}\otimes \mathsf{P}_{001} \rangle =1=\langle\mathsf{1} \rangle$,
i.e., 
$	\mathsf{P}_{001}\otimes \mathsf{P}_{001}=\mathsf{1}\,$.  This result indicates that fusing two $\mathsf{P}_{001}$'s leads to the trivial excitation. It also   means the antiparticle of $\mathsf{P}_{001}$ is itself, which is   consistent with the group structure of $\mathbb{Z}_2$ cyclic group.


\textbf{\textit{Fusing a $\mathbb{Z} _{N_3}$ particle and a $\mathbb{Z} _{N_1}$ loop.}}
In the path-integral representation, fusion of a $\mathbb{Z} _{N_3}$ particle and a $\mathbb{Z} _{N_1}$ loop is represented by:
\begin{align}
&	\langle\mathsf{P}_{001}\otimes \mathsf{L}_{100} \rangle  =\frac{1}{\mathcal{Z}}\int{\mathcal{D} \left[ ABC \right] \exp \left( iS \right) \exp \left( i\int_{\gamma}{A^3}+i\int_{\sigma}{B^1} \right)} =\langle \mathsf{L}_{100}^{001} \rangle\,.
\end{align}
This result tells us that a $\mathbb{Z} _{N_3}$ decorated $\mathbb{Z} _{N_1}$ loop can be constructed by fusing a $\mathbb{Z} _{N_3}$ particle and a $\mathbb{Z} _{N_1}$ loop.


\textbf{\textit{Fusing  a $\mathbb{Z} _{N_1}$ loop and a $\mathbb{Z} _{N_3}$ membrane.}}
Fusion of a $\mathbb{Z} _{N_1}$ loop and a $\mathbb{Z} _{N_3}$ membrane is
\begin{align}
&	\langle \mathsf{L}_{100}\otimes \mathsf{M^S}_{001} \rangle \nonumber\\
	=&\frac{1}{\mathcal{Z}}\int\!\! \mathcal{D} \!\left[ ABC \right]\! \exp \left( iS \right) \times2\exp \left( i\int_{\sigma}{B^1}+i\int_{\omega}{C^3+f} \right)\delta \left( \int_{\sigma}{B^1} \right) \delta \left( \int_{\sigma}{B^2} \right) \nonumber
	\\
	=&\frac{1}{\mathcal{Z}}\int\!\! \mathcal{D} \!\left[ ABC \right]\!   \exp \left( iS \right) \times 2\exp \left( i\int_{\omega}{C^3+f} \right)\delta \left( \int_{\sigma}{B^1} \right) \delta \left( \int_{\sigma}{B^2} \right) =\langle \mathsf{M^S}_{001} \rangle \,,
\end{align}
where $f=\frac{1}{2}\frac{2\pi q}{N_3}\left( d^{-1}B^1B^2+d^{-1}B^2B^1 \right) $.
This is reasonable because $\mathsf{M^S}_{001}$ is equivalent to $\mathsf{M^S}_{001}^{,100}$ due to the delta functional $\delta \left( \int_{\sigma}{B^1} \right) $. Thus, $\mathsf{L}_{100}$ behaves trivially on $\mathsf{M^S}_{001}$.



\textbf{\textit{Fusing a $\mathbb{Z} _{N_3}$ membrane and a $\mathbb{Z} _{N_3}$ membrane.}}
Fusion of a $\mathbb{Z} _{N_3}$ membrane and a $\mathbb{Z} _{N_3}$ membrane is
\begin{align}
	&\langle \mathsf{M^S}_{001}\otimes \mathsf{M^S}_{001} \rangle  \nonumber\\
	=& \frac{1}{\mathcal{Z}}\!\!\int\!\! \mathcal{D} \!\left[ ABC \right]\!   \exp \left( iS \right)\nonumber
	\\
	&\times 2\exp \left[ i\int_{\omega}{C^3+\frac{1}{2}\frac{2\pi q}{N_3}\left( d^{-1}B^1B^2+d^{-1}B^2B^1 \right)} \right] \delta \left( \int_{\sigma}{B^1} \right) \delta \left( \int_{\sigma}{B^2} \right) \nonumber
	\\
	&\times 2\exp \left[ i\!\int_{\omega}{C^3+\frac{1}{2}\frac{2\pi q}{N_3}\left( d^{-1}B^1B^2+d^{-1}B^2B^1 \right)} \right] \!\delta \left( \int_{\sigma}{B^1} \right)\! \delta \left( \int_{\sigma}{B^2} \right) \,.
\end{align}
Integrating out Lagrange's multipliers, i.e., $\tilde{B}^1$, $\tilde{B}^2$, and $A^3$, we obtain $\int_{\sigma}{\hat{B}^1}=\frac{2\pi m_1}{N_1}$, $\int_{\sigma}{\hat{B}^2}=\frac{2\pi m_2}{N_2}$, $\int_{\omega}{\hat{C}^3}=\frac{2\pi m_3}{N_3}$ with $m_1,m_2,m_3\in\Z$, where we add a hat for each gauge field in order to specify the configurations after the Lagrange's multipliers are integrated. As a result, 
\begin{align}
\!\!\!\!\!\!	\!\!\langle \mathsf{M^S}_{001}\!\otimes \!\mathsf{M^S}_{001} \rangle \!=&\frac{1}{\mathcal{Z}_{\rm eff}}   \!\!\sum_{\hat{B}^1,\hat{B}^2,\hat{A}^3}\!\! \bigg\{e^{iS_{\rm eff}}  \!\!
	\times \!\!4\exp\! \left[ i2\int_{\omega}{\frac{1}{2}\frac{2\pi}{N_3}\frac{pN_1N_2N_3}{\left( 2\pi \right) ^2N_{123}}\left( d^{-1}\hat{B}^1\hat{B}^2+d^{-1}\hat{B}^2\hat{B}^1 \!\right)} \!\right] \nonumber
	\\
	&\times \delta \left( \int_{\sigma}{\hat{B}^1} \right) \delta \left( \int_{\sigma}{\hat{B}^2} \right)     \bigg\}\,,
	\label{eq_illustrate_nonAbelianfusion}
\end{align}
where two redundant delta functionals have been dropped as one delta functional  is sufficient for each constraint. The summation $\sum_{\hat{B}^1,\hat{B}^2,\hat{A}^3}$ denotes the functional integration over  the configurations after integrating Lagrange's multipliers. The effective partition function $\mathcal{Z}_{\rm eff}=\sum_{\hat{B}^1,\hat{B}^2,\hat{A}^3}\exp(iS_{\rm eff})$ with the effective action: 
\begin{align}
S_{\rm eff}= \int{\frac{pN_1N_2N_3}{\left( 2\pi \right) ^2N_{123}}\hat{B}^1\hat{B}^2\hat{A}^3}\,.
\end{align}
In deriving eq.~(\ref{eq_illustrate_nonAbelianfusion}), the exponential term ``$\exp\left(i2\times\int C^3\right)$'' has been dropped since it equals to one. Since $\int_{\sigma}{\hat{B}^1}=\frac{2\pi m_1}{N_1}$, $\int_{S^1}{d^{-1}\hat{B}^1}=\int_{S^1}{\int_{\mathcal{A} \in \sigma}{\hat{B}^1}}=\frac{2\pi k_1}{N_1}$ with $k_1$ being an integer and there exists $k_{1}^{\prime}$ such that $k_1+k_{1}^{\prime}=m_1$, we have 
$	\int_{\omega}{d^{-1}\hat{B}^1\hat{B}^2}=\int_{S^1}{d^{-1}\hat{B}^1}\times \int_{\sigma}{\hat{B}^2}=\frac{2\pi k_1}{N_1}\frac{2\pi m_2}{N_2}
 $. Likewise, for $\int_{\omega}{d^{-1}\hat{B}^2\hat{B}^1}$, we have similar discussion and we conclude that:
 $
	\int_{\omega}{d^{-1}\hat{B}^2\hat{B}^1}=\frac{2\pi m_1}{N_1}\frac{2\pi k_2}{N_2}
 $. According to the discussion in Appendix \ref{ap1},     $\exp \left[ i2\times \int \frac{1}{2}\frac{2\pi q}{N_3}\left( d^{-1}B^1B^2+d^{-1}B^2B^1 \right) \right] =1$ can also be dropped.
 By further noting that the delta functionals in eq.~(\ref{eq_illustrate_nonAbelianfusion}) can be expanded via Fourier-like transformations $(i=1,2)$:
\begin{align}
	\delta \left( \int_{\sigma}{\hat{B}^i} \right) &=\delta \left( \frac{2\pi m_i}{N_i} \right) =\frac{1}{2}\left[ 1+\exp \left( \frac{i2\pi m_i}{2} \right) \right] \,,
\end{align}
%One can also verify that
%\begin{align}
%	\delta \left( \int_{\sigma}{\hat{B}^1} \right) \delta \left( \int_{\sigma}{\hat{B}^2} \right) =\delta \left( \int_{\sigma}{\hat{B}^1} \right) \delta \left( \int_{\sigma}{\hat{B}^2} \right) \delta \left( \int_{\sigma}{\hat{B}^1} \right) \delta \left( \int_{\sigma}{\hat{B}^2} \right) 
%	\label{eq13}
%\end{align}
we finally    obtain the following relation between fusion input and outcome:
\begin{align}
	&\langle \mathsf{M^S}_{001}\otimes \mathsf{M^S}_{001} \rangle \nonumber\\
	=& \frac{1}{\mathcal{Z}_{\rm eff}}  \!  \sum_{\hat{B}^1,\hat{B}^2,\hat{A}^3} \! \exp \! \left( iS_{\rm eff}\right) \! 
\left[\!  1+\exp \left( i\int_{\sigma}{\hat{B}^1} \right) +\exp \left( i\int_{\sigma}{\hat{B}^2} \right) +\exp \left( i\int_{\sigma}{\hat{B}^1+\hat{B}^2} \right) \! \right] \nonumber
	\\
\!\!=& \frac{1}{\mathcal{Z}}\!\!\int\!\! \mathcal{D} \!\left[ ABC \right]\!\exp \left( iS \right) \!\!
\left[ 1\!+\!\exp \left( i\int_{\sigma}\!{ {B}^1} \right) \!+\!\exp \left(\! i\int_{\sigma}\!{ {B}^2} \!\right) \!+\!\exp \left(\! i\int_{\sigma}\!{ {B}^1\!+\! {B}^2} \!\right) \right] \nonumber
	\\
	=&\langle \mathsf{1}\oplus \mathsf{L}_{100}\oplus \mathsf{L}_{010}\oplus \mathsf{L}_{110} \rangle	\label{equation_Ms001Ms001}\,.
\end{align}
It is apparent that this is a non-Abelian fusion with multiple fusion outcomes. Fusing two identical $\mathbb{Z} _{N_3}$ membranes renders   one  vacuum (i.e., trivial excitation), one $\mathbb{Z} _{N_1}$ loop, one $\mathbb{Z} _{N_2}$ loop, and one $ \left( \mathbb{Z} _{N_1},\mathbb{Z} _{N_2} \right) $ loop. 




\textbf{\textit{Fusing an $\mathsf{M^{ST}}$-excitation and a $\mathbb{Z} _{N_3}$ membrane.}} 
Fusion of an $\mathsf{M^{ST}}$ and a $\mathbb{Z} _{N_3}$ membrane is
\begin{align}
	&\langle \mathsf{M^{ST}}\otimes \mathsf{M^S}_{001} \rangle \nonumber\\
	=&\frac{1}{\mathcal{Z}}\int{\mathcal{D} \left[ ABC\right]  \exp \left( iS \right)} 
	\times 8\exp \left[ i\int_{\omega}{C^3+\frac{1}{2}\frac{2\pi q}{N_3}\left( d^{-1}B^1B^2+d^{-1}B^2B^1 \right)} \right] \nonumber
	\\
	&\times \exp \left[ i\int_{\tilde{\omega}}{C^3+\frac{1}{2}\frac{2\pi q}{N_3}\left( d^{-1}B^1B^2+d^{-1}B^2B^1 \right)} \right] \delta \left( \int_{\sigma}{B^1} \right) \delta \left( \int_{\sigma}{B^2} \right) \nonumber
	\\
	&\times 2\exp \left[ i\int_{\omega}{C^3+\frac{1}{2}\frac{2\pi q}{N_3}\left( d^{-1}B^1B^2+d^{-1}B^2B^1 \right)} \right] \delta \left( \int_{\sigma}{B^1} \right) \delta \left( \int_{\sigma}{B^2} \right) \,.
\end{align}
Similar to the derivation of  eq.~(\ref{equation_Ms001Ms001}),  we can obtain the fusion rule below:
 $	\langle \mathsf{M^{ST}}\otimes \mathsf{M^S}_{001} \rangle =\langle 4\cdot \mathsf{M^T}_{001} \rangle 
$.  As we mentioned above, $\mathsf{M^{ST}}$ is  constructed by fusing an $\mathsf{M^S}_{001}$ and an $\mathsf{M^T}_{001}$. So we can also calculate $\mathsf{M^{ST}}\otimes \mathsf{M^S}_{001}$ in an alternative way:
\begin{align}
	\mathsf{M^{ST}}\otimes \mathsf{M^S}_{001}=&\mathsf{M^S}_{001}\otimes \mathsf{M^T}_{001}\otimes \mathsf{M^S}_{001}=\mathsf{M^T}_{001}\otimes \mathsf{M^S}_{001}\otimes \mathsf{M^S}_{001} \nonumber
	\\
	=&\mathsf{M^T}_{001}\otimes \left(\mathsf{1}\oplus \mathsf{L}_{100}\oplus \mathsf{L}_{010}\oplus \mathsf{L}_{110}\right) 
	=\mathsf{M^T}_{001}\oplus \mathsf{M^T}_{001}\oplus \mathsf{M^T}_{001}\oplus \mathsf{M^T}_{001} \nonumber
	\\
	=&4\cdot \mathsf{M^T}_{001}\,,
\end{align}
where we have applied the fusion rule in eq.~(\ref{equation_Ms001Ms001}) and the properties of fusion rules listed in eq.~(\ref{eq_fusion_algebra_rule}). 
 
\subsection{Fusion table and quantum dimensions}
For the $BBA$ twisted term, we can exhaust all combinations of two topological excitations and calculate their fusion rules. We list all possible fusion processes and the corresponding fusion rules in table~\ref{table_fusionBBA}. There are $8$ Abelian excitations and $6$ non-Abelian excitations. We can see that this fusion table is symmetric, which is consistent with the condition $\mathsf{a}\otimes \mathsf{b}=\mathsf{b}\otimes \mathsf{a}$. 
\begin{sidewaystable*}
	\caption{\label{table_fusionBBA}Fusion table for the TQFT action~(\ref{eq_action_BBA}). We call an excitation an Abelian excitation if it always has a single fusion channel with any other excitations; otherwise, we call it a non-Abelian excitation. This table shows that there are $8$ Abelian excitations and $6$ non-Abelian excitations. We can directly obtain quantum dimension of each excitation via computing  $N_{\mathsf{a}}$ from this table (see table~\ref{table_quantumdBBA}).}
	\centering
	\tiny
	\begin{tabular*}{\textwidth}{@{\extracolsep{\fill}}|c|c|c|c|c|c|c|c|c|c|c|c|c|c|c|}
		 \hline
		 
		 \hline
		 & \multicolumn{8}{c|}{\textbf{ABELIAN EXCITATIONS}} & \multicolumn{6}{c|}{\textbf{NON-ABELIAN EXCITATIONS}} \\
		 \hline
		 $\otimes$ & $\mathsf{1}$ & $\mathsf{P}_{001}$ & $\mathsf{L}_{100}$ & $\mathsf{L}_{010}$ & $\mathsf{L}_{110}$ & $\mathsf{L}_{100}^{001}$ & $\mathsf{L}_{010}^{001}$ & $\mathsf{L}_{110}^{001}$ & $\mathsf{M^S}_{001}$ & $\mathsf{M^S}_{001}^{001,}$ & $\mathsf{M^T}_{001}$ & $\mathsf{M^T}_{001}^{001,}$ & $\mathsf{M^{ST}}$ & $\mathsf{M^{ST}}^{001,}$ \\
		 \hline
		 $\mathsf{1}$ & $\mathsf{1}$ & $\mathsf{P}_{001}$ & $\mathsf{L}_{100}$ & $\mathsf{L}_{010}$ & $\mathsf{L}_{110}$ & $\mathsf{L}_{100}^{001}$ & $\mathsf{L}_{010}^{001}$ & $\mathsf{L}_{110}^{001}$ & $\mathsf{M^S}_{001}$ & $\mathsf{M^S}_{001}^{001,}$ & $\mathsf{M^T}_{001}$ & $\mathsf{M^T}_{001}^{001,}$ &  $\mathsf{M^{ST}}$ & $\mathsf{M^{ST}}^{001,}$\\
		 \hline
		 $\mathsf{P}_{001}$ & $\mathsf{P}_{001}$ & $\mathsf{1}$ & $\mathsf{L}_{100}^{001}$ & $\mathsf{L}_{010}^{001}$ & $\mathsf{L}_{110}^{001}$ &  $\mathsf{L}_{100}$ & $\mathsf{L}_{010}$ & $\mathsf{L}_{110}$ & $\mathsf{M^S}_{001}^{001,}$ & $\mathsf{M^S}_{001}$ & $\mathsf{M^T}_{001}^{001,}$ & $\mathsf{M^T}_{001}$  & $\mathsf{M^{ST}}^{001,}$ & $\mathsf{M^{ST}}$ \\
		 \hline
		 $\mathsf{L}_{100}$ & $\mathsf{L}_{100}$ & $\mathsf{L}_{100}^{001}$ & $\mathsf{1}$ & $\mathsf{L}_{110}$ & $\mathsf{L}_{010}$ & $\mathsf{P}_{001}$ & $\mathsf{L}_{110}^{001}$ & $\mathsf{L}_{010}^{001}$ & $\mathsf{M^S}_{001}$ & $\mathsf{M^S}_{001}^{001,}$ & $\mathsf{M^T}_{001}$ & $\mathsf{M^T}_{001}^{001,}$ & $\mathsf{M^{ST}}$ & $\mathsf{M^{ST}}^{001,}$\\
		 \hline
		 $\mathsf{L}_{010}$ & $\mathsf{L}_{010}$ & $\mathsf{L}_{010}^{001}$ & $\mathsf{L}_{110}$ & $\mathsf{1}$ & $\mathsf{L}_{100}$ & $\mathsf{L}_{110}^{001}$ & $\mathsf{P}_{001}$ & $\mathsf{L}_{100}^{001}$ & $\mathsf{M^S}_{001}$ & $\mathsf{M^S}_{001}^{001,}$ & $\mathsf{M^T}_{001}$ & $\mathsf{M^T}_{001}^{001,}$ & $\mathsf{M^{ST}}$ & $\mathsf{M^{ST}}^{001,}$ \\
		 \hline
		 $\mathsf{L}_{110}$ & $\mathsf{L}_{110}$ & $\mathsf{L}_{110}^{001}$ & $\mathsf{L}_{010}$ & $\mathsf{L}_{100}$ & $\mathsf{1}$ & $\mathsf{L}_{010}^{001}$ & $\mathsf{L}_{100}^{001}$ & $\mathsf{P}_{001}$ & $\mathsf{M^S}_{001}$ &  $\mathsf{M^S}_{001}^{001,}$& $\mathsf{M^T}_{001}$ & $\mathsf{M^T}_{001}^{001,}$ & $\mathsf{M^{ST}}$ & $\mathsf{M^{ST}}^{001,}$\\
		 \hline
		 $\mathsf{L}_{100}^{001}$ & $\mathsf{L}_{100}^{001}$ & $\mathsf{L}_{100}$ & $\mathsf{P}_{001}$ & $\mathsf{L}_{110}^{001}$ & $\mathsf{L}_{010}^{001}$ & $\mathsf{1}$ & $\mathsf{L}_{110}$ & $\mathsf{L}_{010}$ & $\mathsf{M^S}_{001}^{001,}$ & $\mathsf{M^S}_{001}$ & $\mathsf{M^T}_{001}^{001,}$ & $\mathsf{M^T}_{001}$ & $\mathsf{M^{ST}}^{001,}$ & $\mathsf{M^{ST}}$ \\
		 \hline
		 $\mathsf{L}_{010}^{001}$ & $\mathsf{L}_{010}^{001}$ & $\mathsf{L}_{010}$ & $\mathsf{L}_{110}^{001}$ & $\mathsf{P}_{001}$ & $\mathsf{L}_{100}^{001}$ & $\mathsf{L}_{110}$ & $\mathsf{1}$ & $\mathsf{L}_{100}$ & $\mathsf{M^S}_{001}^{001,}$ & $\mathsf{M^S}_{001}$ & $\mathsf{M^T}_{001}^{001,}$ & $\mathsf{M^T}_{001}$ & $\mathsf{M^{ST}}^{001,}$ & $\mathsf{M^{ST}}$ \\
		 \hline
		 $\mathsf{L}_{110}^{001}$ & $\mathsf{L}_{110}^{001}$ & $\mathsf{L}_{110}$ & $\mathsf{L}_{010}^{001}$ & $\mathsf{L}_{100}^{001}$ & $\mathsf{P}_{001}$ & $\mathsf{L}_{010}$ & $\mathsf{L}_{100}$ & $\mathsf{1}$ & $\mathsf{M^S}_{001}^{001,}$ & $\mathsf{M^S}_{001}$ & $\mathsf{M^T}_{001}^{001,}$ & $\mathsf{M^T}_{001}$ & $\mathsf{M^{ST}}^{001,}$ & $\mathsf{M^{ST}}$ \\
		 \hline
		 $\mathsf{M^S}_{001}$ & $\mathsf{M^S}_{001}$ & $\mathsf{M^S}_{001}^{001,}$ & $\mathsf{M^S}_{001}$ & $\mathsf{M^S}_{001}$ & $\mathsf{M^S}_{001}$ & $\mathsf{M^S}_{001}^{001,}$ & $\mathsf{M^S}_{001}^{001,}$ & $\mathsf{M^S}_{001}^{001,}$ & \makecell{ $\mathsf{1}$\\  $\oplus \mathsf{L}_{100}$\\  $\oplus \mathsf{L}_{010}$\\  $\oplus \mathsf{L}_{110}$} & \makecell{ $\mathsf{P}_{001}$\\  $\oplus \mathsf{L}_{100}^{001}$\\  $\oplus \mathsf{L}_{010}^{001}$\\  $\oplus \mathsf{L}_{110}^{001}$} & $\mathsf{M^{ST}}$ & $\mathsf{M^{ST}}^{001,}$ & $4\cdot\mathsf{M^T}_{001}$ & $4\cdot\mathsf{M^T}_{001}^{001,}$ \\
		 \hline
		 $\mathsf{M^S}_{001}^{001,}$ & $\mathsf{M^S}_{001}^{001,}$ & $\mathsf{M^S}_{001}$ & $\mathsf{M^S}_{001}^{001,}$ & $\mathsf{M^S}_{001}^{001,}$ & $\mathsf{M^S}_{001}^{001,}$ & $\mathsf{M^S}_{001}$ & $\mathsf{M^S}_{001}$ & $\mathsf{M^S}_{001}$ & \makecell{ $\mathsf{P}_{001}$\\  $\oplus \mathsf{L}_{100}^{001}$\\  $\oplus \mathsf{L}_{010}^{001}$\\  $\oplus \mathsf{L}_{110}^{001}$} & \makecell{ $\mathsf{1}$\\  $\oplus \mathsf{L}_{100}$\\  $\oplus \mathsf{L}_{010}$\\  $\oplus \mathsf{L}_{110}$} & $\mathsf{M^{ST}}^{001,}$ & $\mathsf{M^{ST}}$ & 4$\cdot$$\mathsf{M^T}_{001}^{001,}$ & $4\cdot\mathsf{M^T}_{001}$ \\
		 \hline
		 $\mathsf{M^T}_{001}$ & $\mathsf{M^T}_{001}$ & $\mathsf{M^T}_{001}^{001,}$ & $\mathsf{M^T}_{001}$ & $\mathsf{M^T}_{001}$ & $\mathsf{M^T}_{001}$ & $\mathsf{M^T}_{001}^{001,}$ & $\mathsf{M^T}_{001}^{001,}$ & $\mathsf{M^T}_{001}^{001,}$ & $\mathsf{M^{ST}}$ & $\mathsf{M^{ST}}^{001,}$ & \makecell{ $4\cdot\mathsf{1}$\\  $\oplus4\cdot \mathsf{L}_{100}$\\  $\oplus4\cdot \mathsf{L}_{010}$\\  $\oplus4\cdot \mathsf{L}_{110}$} & \makecell{ $4\cdot \mathsf{P}_{001}$\\  $\oplus4\cdot \mathsf{L}_{100}^{001}$\\  $\oplus4\cdot \mathsf{L}_{010}^{001}$\\  $\oplus4\cdot \mathsf{L}_{110}^{001}$} & $16\cdot \mathsf{M^S}_{001}$ & $16\cdot \mathsf{M^S}_{001}^{001,}$ \\
		 \hline
		 $\mathsf{M^T}_{001}^{001,}$ & $\mathsf{M^T}_{001}^{001,}$ & $\mathsf{M^T}_{001}$ & $\mathsf{M^T}_{001}^{001,}$ & $\mathsf{M^T}_{001}^{001,}$ & $\mathsf{M^T}_{001}^{001,}$ & $\mathsf{M^T}_{001}$ & $\mathsf{M^T}_{001}$ & $\mathsf{M^T}_{001}$ & $\mathsf{M^{ST}}^{001,}$ & $\mathsf{M^{ST}}$ & \makecell{ $4\cdot \mathsf{P}_{001}$\\  $\oplus4\cdot \mathsf{L}_{100}^{001}$\\  $\oplus4\cdot \mathsf{L}_{010}^{001}$\\  $\oplus4\cdot \mathsf{L}_{110}^{001}$} & \makecell{ $4\cdot\mathsf{1}$\\  $\oplus4\cdot \mathsf{L}_{100}$\\  $\oplus4\cdot \mathsf{L}_{010}$\\  $\oplus4\cdot \mathsf{L}_{110}$} & $16\cdot \mathsf{M^S}_{001}^{001,}$ & $16\cdot \mathsf{M^S}_{001}$ \\
		 \hline
		 $\mathsf{M^{ST}}$ & $\mathsf{M^{ST}}$ & $\mathsf{M^{ST}}^{001,}$ & $\mathsf{M^{ST}}$ & $\mathsf{M^{ST}}$ & $\mathsf{M^{ST}}$ & $\mathsf{M^{ST}}^{001,}$ & $\mathsf{M^{ST}}^{001,}$ & $\mathsf{M^{ST}}^{001,}$ & $4\cdot\mathsf{M^T}_{001}$ & $4\cdot\mathsf{M^T}_{001}^{001,}$ & $16\cdot \mathsf{M^S}_{001}$ & $16\cdot \mathsf{M^S}_{001}^{001,}$ & \makecell{ $16\cdot\mathsf{1}$\\  $\oplus16\cdot \mathsf{L}_{100}$\\  $\oplus16\cdot \mathsf{L}_{010}$\\  $\oplus16\cdot \mathsf{L}_{110}$} & \makecell{ $16\cdot \mathsf{P}_{001}$\\  $\oplus16\cdot \mathsf{L}_{100}^{001}$\\  $\oplus16\cdot \mathsf{L}_{010}^{001}$\\  $\oplus16\cdot \mathsf{L}_{110}^{001}$} \\
		 \hline
		 $\mathsf{M^{ST}}^{001,}$ & $\mathsf{M^{ST}}^{001,}$ & $\mathsf{M^{ST}}$ & $\mathsf{M^{ST}}^{001,}$ & $\mathsf{M^{ST}}^{001,}$ & $\mathsf{M^{ST}}^{001,}$ & $\mathsf{M^{ST}}$ & $\mathsf{M^{ST}}$ & $\mathsf{M^{ST}}$ & $4\cdot\mathsf{M^T}_{001}^{001,}$ & $4\cdot\mathsf{M^T}_{001}$ & $16\cdot \mathsf{M^S}_{001}^{001,}$ & $16\cdot \mathsf{M^S}_{001}$ & \makecell{ $16\cdot \mathsf{P}_{001}$\\  $\oplus16\cdot \mathsf{L}_{100}^{001}$\\  $\oplus16\cdot \mathsf{L}_{010}^{001}$\\  $\oplus16\cdot \mathsf{L}_{110}^{001}$} & \makecell{ $16\cdot\mathsf{1}$\\  $\oplus16\cdot \mathsf{L}_{100}$\\  $\oplus16\cdot \mathsf{L}_{010}$\\  $\oplus16\cdot \mathsf{L}_{110}$} \\
		 \hline
	\end{tabular*}
\end{sidewaystable*}

For a topological excitation $\mathsf{a}$, we can define a fusion coefficient matrix $N_{\mathsf{a}}$, whose element is $ \left( N_{\mathsf{a}} \right) _{\mathsf{b}\mathsf{c}}=N_{\mathsf{c}}^{\mathsf{a}\mathsf{b}}$. Quantum dimension is defined by the greatest eigenvalue of the matrix $N_{\mathsf{a}}$. For example, we can obtain the fusion coefficient matrix for $\mathsf{M^S}_{001}^{001,}$:
\begin{align}
	N_{\mathsf{M^S}_{001}^{001,}}=\left[\begin{array}{cccccccccccccc}
		0&0&0&0&0&0&0&0&0&1&0&0&0&0\\
		0&0&0&0&0&0&0&0&1&0&0&0&0&0\\
		0&0&0&0&0&0&0&0&0&1&0&0&0&0\\
		0&0&0&0&0&0&0&0&0&1&0&0&0&0\\
		0&0&0&0&0&0&0&0&0&1&0&0&0&0\\
		0&0&0&0&0&0&0&0&1&0&0&0&0&0\\
		0&0&0&0&0&0&0&0&1&0&0&0&0&0\\
		0&0&0&0&0&0&0&0&1&0&0&0&0&0\\
		0&1&0&0&0&1&1&1&0&0&0&0&0&0\\
		1&0&1&1&1&0&0&0&0&0&0&0&0&0\\
		0&0&0&0&0&0&0&0&0&0&0&0&0&1\\
		0&0&0&0&0&0&0&0&0&0&0&0&1&0\\
		0&0&0&0&0&0&0&0&0&0&0&4&0&0\\
		0&0&0&0&0&0&0&0&0&0&4&0&0&0\\
	\end{array}\,.
	\right]
\end{align}
We can calculate that its eigenvalues are 
\begin{align}
	\text{eig}\left(N_{\mathsf{M^S}_{001}^{001,}}\right)=\left(2,2,2,2,0,0,0,0,0,0,-2,-2,-2,-2\right)\,.
\end{align}
Thus the quantum dimension of $\mathsf{M^S}_{001}^{001,}$ is $2$.
We can obtain all matrices $N_{\mathsf{a}}$ from table~\ref{table_fusionBBA} and then calculate the quantum dimensions of all excitations, as shown in table~\ref{table_quantumdBBA}. We conclude that for a non-Abelian excitation, the quantum dimension is larger than $1$. For an Abelian excitation, the quantum dimension is always $1$.
\begin{table}
	\caption{\label{table_quantumdBBA}Quantum dimension of excitations in table~\ref{table_fusionBBA}. For a non-Abelian excitation, the quantum dimension is larger than $1$ and equal to the normalization factor of the Wilson operator. For an Abelian excitation, the quantum dimension is always $1$. }
	\centering
	\footnotesize
	\begin{tabular*}{\textwidth}{@{\extracolsep{\fill}}ccccccccccccccc}
		\toprule
		Excitation & $\mathsf{1}$ & $\mathsf{P}_{001}$ & $\mathsf{L}_{100}$ & $\mathsf{L}_{010}$ & $\mathsf{L}_{110}$ & $\mathsf{L}_{100}^{001}$ & $\mathsf{L}_{010}^{001}$ & $\mathsf{L}_{110}^{001}$ & $\mathsf{M^S}_{001}$ & $\mathsf{M^S}_{001}^{001,}$ & $\mathsf{M^T}_{001}$ & $\mathsf{M^T}_{001}^{001,}$ & $\mathsf{M^{ST}}$ & $\mathsf{M^{ST}}^{001,}$ \\
		\hline
		Quantum dimension & $1$ & $1$ & $1$ & $1$ & $1$ & $1$ & $1$ & $1$ & $2$ & $2$ & $4$ & $4$ & $8$ & $8$ \\
		\bottomrule
	\end{tabular*}
\end{table}


 
\section{Shrinking rules for the $BBA$ twisted term\label{s4}}



\subsection{General discussion on shrinking in $5$D and its hierarchical structure}
In $3$D topological orders, shrinking rules are absent because all topological excitations are  point-like particles, i.e., anyons. However, in $4$D and higher, spatially extended excitations exist, such that the inclusion of shrinking processes makes sense. While loop excitations in $4$D brings us  exotic braiding statistics, as reviewed in section~\ref{s1}, one may wonder  what are the resulting outcomes  by adiabatically shrinking loop excitations. Indeed, a recent study in $4$D indicates that  the topological data given by shrinking rules   significantly enrich  the physics of  higher-dimensional topological orders~\cite{PhysRevB.107.165117}. 
More concretely, the world-sheet of a loop excitation is an $S^1\times S^1=T^2$ spacetime manifold. If we shrink the loop to a point, then the world-sheet is simultaneously shrunk  to a world-line, i.e., $T^2\rightarrow S^1$, which can be either the world-line of a particle or many distinct particles. The former case is an Abelian shrinking process while the latter is a non-Abelian shrinking process.  
\begin{figure}
	\centering
	\includegraphics[scale=0.8,keepaspectratio]{fig_shrinking}
	\caption{Typical shrinking rules for loops and $S^2$ membranes in the $5$D twisted $BF$ theory with the $BBA$ twisted term. Shrinking an $\mathsf{L}_{100}$ and shrinking  an $\mathsf{L}_{100}^{001}$ respectively lead to a vacuum and a $\mathsf{P}_{001}$. Shrinking an $\mathsf{M^S}_{001}$ and shrinking an $\mathsf{M^S}_{001}^{001,}$ respectively lead to a direct sum of $2$ vacua and a direct sum  of $2$ $\mathsf{P}_{001}$'s.}
	\label{fig_shrinking}
\end{figure}
\begin{figure}
	\centering
	\includegraphics[scale=0.8,keepaspectratio]{fig_hierarchy_nontrivial}
	\caption{A typical example of hierarchical shrinking rules for $T^2$ membranes in the $5$D twisted $BF$ theory with the $BBA$ twisted term. We can shrink a $T^2$ membrane $\mathsf{M^T}_{001}$ to a superposition of a vacuum, a $\mathsf{L}_{100}$, a $\mathsf{L}_{010}$ and a $\mathsf{L}_{110}$ first. Then we can continue to shrink them to $4$ vacua.}
	\label{fig_hierarchy_nontrivial}
\end{figure}
\begin{figure}
	\centering
	\includegraphics[scale=0.9,keepaspectratio]{fig_hierarchy_trivial}
	\caption{A typical example of shrinking rules without hierarchical structure for $T^2$ membranes in the $5$D twisted $BF$ theory with the $AAAdA$ twisted term.}
	\label{fig_hierarchy_trivial}
\end{figure}

In $5$D topological orders, spatially extended excitations include both loops and membranes, and the latter can have different shapes, e.g., $S^2$ and $T^2$. Therefore, we expect the underlying shrinking rules will be significantly diversified, as shown in figures \ref{fig_shrinking}, \ref{fig_hierarchy_nontrivial},  and \ref{fig_hierarchy_trivial}. Intuitively,  an $S^2$ membrane can be directly shrunk to particles while a $T^2$ membrane can be shrunk to loops and particles first, and then we can further  shrink these loops to particles. As membranes exist in topological orders of $5$D and higher, this \textit{hierarchical structure} of shrinking rules can only appear in $5$D and higher. A typical example of hierarchical shrinking is shown in figure \ref{fig_hierarchy_nontrivial}. 

It seems that   two successive steps are required to shrink $T^2$ membranes completely into particles. But it is not always the case. As shown in figure \ref{fig_hierarchy_trivial}, some $T^2$ membranes in the twisted $BF$ theory with the $AAAdA$ twisted term, upon being shrunk in the first step, only lead to trivial loops and some particles. As trivial loops and trivial particles are equivalent to each other and all of them are identical to a vacuum, these $T^2$ membranes are completely shrunk to particles including vacua in the first step. If we need two steps to shrink a $T^2$ membrane to particles, we say this $T^2$ membrane has hierarchical shrinking rules. 
 
 
 
To quantitatively calculate shrinking rules, we  define a shrinking operator $\mathcal{S}$ and the associated shrinking coefficient $\mathrm{S}_{\mathsf{b}}^{\mathsf{a}}$ as below:
\begin{align}
\mathcal{S}\left(\mathsf{a}\right)=\underset{X_{1}\rightarrow X_{2}}{\lim}\mathsf{a}=\oplus_{\mathsf{b}} \mathrm{S}^{\mathsf{a}}_{\mathsf{b}} \mathsf{b}\,,
\end{align}
where the shrinking coefficient $S^{\mathsf{a}}_{\mathsf{b}}\in\Z$. $X_1$ and $X_2$ are respectively  spacetime trajectories of $\mathsf{a}$ before and after shrinking, which respects $X_2\subset X_1$.   The summation   $\oplus_{\mathsf{b}}$ exhausts all $14$ topologically distinct excitations.  We leave details of $X_1$ and $X_2$ to section~\ref{section_shrinking_nohier} and \ref{section_shrinking_hier} where concrete calculations are given.  
 
In the path-integral representation, we may compute the following expectation value in the ground state:
\begin{align}
	\langle \mathcal{S} \left( \mathsf{a} \right) \rangle =&\langle \underset{X_1\rightarrow X_2}{\lim}\mathsf{a} \rangle =\underset{X_1\rightarrow X_2}{\lim}\frac{1}{\mathcal{Z}}\int{\mathcal{D} \left[ ABC \right] \exp \left( iS \right) \mathcal{O} _{\mathsf{a}}} \nonumber
	\\
	=&\sum_{\mathsf{b}}\frac{1}{\mathcal{Z}}\int{\mathcal{D} \left[ ABC \right] \exp \left( iS \right) \mathrm{S}_{\mathsf{b}}^{\mathsf{a}}\mathcal{O} _{\mathsf{b}}} =\langle \oplus_{\mathsf{b}}\mathrm{S}_{\mathsf{b}}^{\mathsf{a}}\mathsf{b} \rangle \,.
	\label{eq_shrinking_pathintegral}
\end{align}
 If the integer $\mathrm{S}_{\mathsf{b}}^{\mathsf{a}}=1$ for some topological excitation $\mathsf{b}$ and vanishes for any other excitations,  the shrinking is said to be Abelian. If $\mathrm{S}_{\mathsf{b}}^{\mathsf{a}}$ is nonzero for more than one topological excitations or if there exists $\mathsf{b}$ such that  $\mathrm{S}_{\mathsf{b}}^{\mathsf{a}}\textgreater 1$,  the shrinking is said to be non-Abelian. We collect all shrinking rules in the  shrinking table~\ref{table_shrinking}.
 
 
 
 
\subsection{Shrinking rules without hierarchical structure\label{section_shrinking_nohier}}
Next, we present several typical examples in table~\ref{table_shrinking}.  In this subsection, we study shrinking rules without hierarchical structure. $X_1$ and $X_2$ meet the following conditions: (i) if $\mathsf{a}$ is a particle ($\mathsf{1}$ and $\mathsf{P}_{001}$), then the world-line of the particle is unchanged, i.e., $X_1=X_2=\gamma$, which means that particles are unshrinkable; (ii) if $\mathsf{a}$ is a loop, then the world-sheet of the loop is shrunk into a world-line, i.e., $X_1=\sigma, X_2=\gamma$; (iii) if $\mathsf{a}$ is an $S^2$ membrane ($\mathsf{M^S}_{001}$ and $\mathsf{M^S}_{001}^{001,}$), then the world-volume  of the membrane is shrunk into a world-line and the manifold $\sigma$ in the delta functionals is also shrunk into a world-line, i.e., $X_1=(\omega,\sigma), X_2=(\gamma,\gamma)$.



\textbf{\textit{Shrinking a $\mathbb{Z}_{N_1}$ loop.}}
Shrinking a $\mathbb{Z}_{N_1}$ loop results a collapse of  its world sheet to a world-line, which is $\sigma\rightarrow \gamma$. From the path integral, we obtain
\begin{align}
	\langle \mathcal{S} \left( \mathsf{L}_{100} \right) \rangle &=\langle \underset{\sigma \rightarrow \gamma}{\lim}\mathsf{L}_{100} \rangle =\underset{\sigma \rightarrow \gamma}{\lim}\frac{1}{\mathcal{Z}}\int{\mathcal{D} \left[ ABC \right] \exp \left( iS \right) \exp \left( i\int_{\sigma}{B^1} \right)}\nonumber
	\\
	&=\frac{1}{\mathcal{Z}}\int{\mathcal{D} \left[ ABC \right] \exp \left( iS \right) \exp \left( i0 \right)}=\langle \mathsf{1} \rangle\,. 
\end{align}
So we conclude that
 $	\mathcal{S} \left( \mathsf{L}_{100} \right)=\mathsf{1}
$, which means that a $\mathbb{Z}_{N_1}$ loop denoted as $\mathsf{L}_{100}$ is shrunk into a vacuum.




\textbf{\textit{Shrinking a $\mathbb{Z}_{N_1}$ loop decorated by $\mathbb{Z}_{N_3}$ particle.}}
Similarly, the shrinking rule for an $\mathsf{L}_{100}^{001}$ loop can be obtained
\begin{align}
	\langle \mathcal{S} \left( \mathsf{L}_{100}^{001} \right) \rangle &=\langle \underset{\sigma \rightarrow \gamma}{\lim}\mathsf{L}_{100}^{001} \rangle =\underset{\sigma \rightarrow \gamma}{\lim}\frac{1}{\mathcal{Z}}\int{\mathcal{D} \left[ ABC\right] \exp \left( iS \right) \exp \left( i\int_{\sigma}{B^1}+i\int_{\gamma}{A^3} \right)}\nonumber
	\\
	&=\frac{1}{\mathcal{Z}}\int{\mathcal{D} \left[ ABC \right] \exp \left( iS \right) \exp \left( i0 \right) \exp \left( i\int_{\gamma}{A^3} \right)}=\langle \mathsf{P}_{001} \rangle \,.
\end{align}
Therefore, particle decoration is unaffected by the shrinking process. Also notice that 
\begin{align}
	\mathcal{S} \left( \mathsf{L}_{100}^{001} \right)=\mathcal{S} \left( \mathsf{L}_{100}\otimes \mathsf{P}_{001} \right)=\mathcal{S} \left( \mathsf{L}_{100} \right)\otimes\mathcal{S} \left( \mathsf{P}_{001} \right)=\mathsf{P}_{001}
\end{align}
implying that shrinking rules respect the fusion rules.


\textbf{\textit{Shrinking a $\mathbb{Z}_{N_3}$ membrane in the shape of $S^2$.}}
From path integral we obtain 
\begin{align}
&	\langle \mathcal{S} \left( \mathsf{M^S}_{001} \right) \rangle 	=\langle \underset{\omega \rightarrow \gamma, \sigma \rightarrow \gamma}{\lim}\mathsf{M^S}_{001} \rangle \nonumber
	\\
	=&\underset{\omega \rightarrow \gamma, \sigma \rightarrow \gamma}{\lim}\frac{1}{\mathcal{Z}}\int{\mathcal{D} \left[ ABC \right] \exp \left( iS \right)}
	\times 2\exp \left[ i\int_{\omega}{C^3+\frac{1}{2}\frac{2\pi q}{N_3}\left( d^{-1}B^1B^2+d^{-1}B^2B^1 \right)} \right]  \nonumber
	\\
	&\times \delta \left( \int_{\sigma}{B^1} \right) \delta \left( \int_{\sigma }{B^2} \right) =\frac{1}{\mathcal{Z}}\int{\mathcal{D} \left[ ABC\right] \exp \left( iS \right) 2\cdot\exp \left( i0 \right) \delta \left( 0 \right) \delta \left( 0 \right)}=\langle 2\cdot\mathsf{1} \rangle \,.
\end{align}
In other words, the shrinking process renders a superposition of $2$ vacua when shrinking a $\mathbb{Z}_{N_3}$ membrane in the shape of $S^2$. This is reasonable because this membrane is a non-Abelian excitation and has a quantum dimension of $2$. If we consider $2\cdot\mathsf{1}$ as an excitation and derive its fusion coefficient matrix $N_{2\cdot\mathsf{1}}$, we will find that $N_{2\cdot\mathsf{1}}$ is $2$ times identity matrix since fusion coefficient matrix of $\mathsf{1}$ is an identity matrix. Thus the greatest eigenvalue of $N_{2\cdot\mathsf{1}}$ is $2$ and we conclude that $2\cdot\mathsf{1}$ has quantum dimension $2$. It may be seen as preserving the quantum dimension. As a matter of fact, by exhausting all shrinking processes, we can verify that \textit{all shrinking processes preserve quantum dimension.}



\subsection{Hierarchical shrinking process (i): shrinking a   $T^2$ membrane\label{section_shrinking_hier}}
Spatially, when shrinking a $T^2$ torus, we will first get an $S^1$ closed line and then get a point. Thus for a $T^2$ membrane, we can first shrink it into loops and particles and then shrink them to particles. If a $T^2$ membrane gives nontrivial loops after the first shrinking process, we will further shrink these loops into  particles. This is  aforementioned hierarchical shrinking. We define that $\mathcal{S}$ will shrink $\tilde{\omega}=T^2\times S^1$ to $S^1\times S^1$ first and then $\mathcal{S}$ shrinks $\sigma=S^1\times S^1$ to $S^1$. Also notice that Wilson operator carries $\delta \left( \int_{\sigma}{B^i} \right)$, where $\sigma$ is a submanifold of the world volume $\tilde{\omega}$. Since when we shrink the world volume $\tilde{\omega}$ to world sheet $S^1\times S^1$, $\sigma$ is still a submanifold of this world sheet, we do not need to shrink $\sigma$ in $\delta \left( \int_{\sigma}{B^i} \right)$. Thus we have
\begin{align}
	\langle \mathcal{S} \left( \mathsf{M^T}_{001} \right) \rangle =&\langle \underset{\tilde{\omega} \rightarrow T^2}{\lim}\mathsf{M^T}_{001} \rangle \nonumber
	\\
	=&\underset{\tilde{\omega} \rightarrow T^2}{\lim}\frac{1}{\mathcal{Z}}\int{\mathcal{D} \left[ ABC\right] \exp \left( iS \right)}\times 4\exp \left[ i\int_{\tilde{\omega}}{C^3+\frac{1}{2}\frac{2\pi q}{N_3}\left( d^{-1}B^1B^2+d^{-1}B^2B^1 \right)} \right] \nonumber
	\\
	&\times\delta \left( \int_{\sigma}{B^1} \right) \delta \left( \int_{\sigma}{B^2} \right) \,.
\end{align}
By further expanding delta functionals, we have
\begin{align}
	\langle \mathcal{S} \left( \mathsf{M^T}_{001} \right) \rangle 
	=&\langle \mathsf{1}\oplus \mathsf{L}_{100}\oplus \mathsf{L}_{010}\oplus \mathsf{L}_{110} \rangle \,.
\end{align}
If we regard $\mathsf{1}$ as a trivial loop, shrinking $\mathsf{M^T}_{001}$ once gives us a superposition of a vacuum and  all three pure loops in table~\ref{tab_excitation}. We also see that the normalization factor $4$ ensures that the shrinking coefficients are integers. 


Since the above shrinking process leads to nontrivial loops, we can continue to shrink the loops to particles (the symbol $\mathcal{S}^2=\mathcal{S}\mathcal{S}$):
\begin{align}
	\langle \mathcal{S} ^2\left( \mathsf{M^T}_{001} \right) \rangle =&\langle \mathcal{S} \left( \mathsf{1}\oplus \mathsf{L}_{100}\oplus \mathsf{L}_{010}\oplus \mathsf{L}_{110} \right) \rangle =\langle \underset{\sigma \rightarrow S^1}{\lim}\left( \mathsf{1}\oplus \mathsf{L}_{100}\oplus \mathsf{L}_{010}\oplus \mathsf{L}_{110} \right) \rangle \nonumber
	\\
	=&\underset{\sigma \rightarrow S^1}{\lim}\frac{1}{\mathcal{Z}}\int{\mathcal{D} \left[ABC  \right] \exp \left( iS \right)}\nonumber
	\\
	&\times \left[ 1+\exp \left( i\int_{\sigma}{B^1} \right) +\exp \left( i\int_{\sigma}{B^2} \right) +\exp \left( i\int_{\sigma}{B^1+B^2} \right) \right] \nonumber
	\\
	=&\frac{1}{\mathcal{Z}}\int{\mathcal{D} \left[ ABC \right] \exp \left( iS \right) \times \left[ 1+\exp \left( i0 \right) +\exp \left( i0 \right) +\exp \left( i0 \right) \right]}\nonumber
	\\
	=&\frac{1}{\mathcal{Z}}\int{\mathcal{D} \left[ ABC \right] \exp \left( iS \right) \times 4\cdot \mathsf{1}}=\langle 4\cdot \mathsf{1} \rangle \,.
\end{align}
In summary,  after two steps of shrinking operation, we finally shrink a $\mathbb{Z}_{N_3}$ membrane in the shape of $T^2$ to $4$ vacua, which is    different from shrinking an $S^2$ membrane.

\subsection{Hierarchical shrinking process (ii): shrinking an $\mathsf{M^{ST}}$-excitation}

 For an $\mathsf{M^{ST}}$-excitation, we define that shrinking operator $\mathcal{S}$ shrinks $\omega=S^2\times S^1$ and $\tilde{\omega}=T^2\times S^1$ to $S^1$ and $S^1\times S^1$ respectively first. Then we can continue to shrink all $S^1\times S^1$ world sheets to $S^1$. Thus we obtain 
\begin{align}
	\langle \mathcal{S} \left( \mathsf{M^{ST}} \right) \rangle =&\langle \underset{\tilde{\omega}\rightarrow T^2,\omega \rightarrow S^1}{\lim}\mathsf{M^{ST}} \rangle \nonumber
	\\
	=&\underset{\tilde{\omega}\rightarrow T^2,\omega \rightarrow S^1}{\lim}\frac{1}{\mathcal{Z}}\int{\mathcal{D} \left[ ABC \right] \exp \left( iS \right)} \nonumber
	\\
	&\times 8\exp \left[ i\int_{\tilde{\omega}}{C^3+\frac{1}{2}\frac{2\pi q}{N_3}\left( d^{-1}B^1B^2+d^{-1}B^2B^1 \right)} \right] \nonumber
	\\
	&\times \exp \left[ i\int_{\omega}{C^3+\frac{1}{2}\frac{2\pi q}{N_3}\left( d^{-1}B^1B^2+d^{-1}B^2B^1 \right)} \right] \delta \left( \int_{\sigma}{B^1} \right) \delta \left( \int_{\sigma}{B^2} \right) \nonumber
	\\
	=&\frac{1}{\mathcal{Z}}\int{\mathcal{D} \left[ ABC\right] \exp \left( iS \right) 8\exp \left( i0 \right)}\exp \left( i0 \right) \nonumber
	\\
	&\times \frac{1}{4}\left[ 1+\exp \left( i\int_{\sigma}{B^1} \right) \right] \left[ 1+\exp \left( i\int_{\sigma}{B^2} \right) \right] \nonumber
	\\
	=&\frac{1}{\mathcal{Z}}\int{\mathcal{D} \left[ ABC \right] \exp \left( iS \right)} \nonumber
	\\
	&\times 2\cdot\left[ 1+\exp \left( i\int_{\sigma}{B^1} \right) +\exp \left( i\int_{\sigma}{B^2} \right) +\exp \left( i\int_{\sigma}{B^1+B^2} \right) \right] \nonumber
	\\
	=&\langle 2\cdot\left(\mathsf{1}\oplus \mathsf{L}_{100}\oplus \mathsf{L}_{010}\oplus \mathsf{L}_{110}\right) \rangle 
\end{align}
If we regard $\mathsf{1}$ as a trivial loop, then shrinking $\mathsf{M^{ST}}$ once will give a $2$ times superposition of vacua and all pure loops.  Next, we continue to shrink all loops to particles:  
\begin{align}
	\langle \mathcal{S} ^2\left( \mathsf{M^{ST}} \right) \rangle =&\langle \mathcal{S} \left[ 2\cdot \left( \mathsf{1}\oplus \mathsf{L}_{100}\oplus \mathsf{L}_{010}\oplus \mathsf{L}_{110} \right) \right]  \rangle =\langle \underset{\sigma \rightarrow S^1}{\lim}\left[ 2\cdot \left( \mathsf{1}\oplus \mathsf{L}_{100}\oplus \mathsf{L}_{010}\oplus \mathsf{L}_{110} \right) \right]  \rangle \nonumber
	\\
	=&\underset{\sigma \rightarrow S^1}{\lim}\frac{1}{\mathcal{Z}}\int{\mathcal{D} \left[ ABC \right] \exp \left( iS \right)} \nonumber
	\\
	&\times 2\left[ 1+\exp \left( i\int_{\sigma}{B^1} \right) +\exp \left( i\int_{\sigma}{B^2} \right) +\exp \left( i\int_{\sigma}{B^1+B^2} \right) \right] \nonumber
	\\
	=&\frac{1}{\mathcal{Z}}\int{\mathcal{D} \left[ ABC\right] \exp \left( iS \right) \times 2\left[ 1+\exp \left( i0 \right) +\exp \left( i0 \right) +\exp \left( i0 \right) \right]} \nonumber
	\\
	=&\frac{1}{\mathcal{Z}}\int{\mathcal{D} \left[ ABC \right] \exp \left( iS \right) \times 8\cdot \mathsf{1}} =\langle 8\cdot \mathsf{1} \rangle 
\end{align}
Finally, we get $8$ vacua.

If we consider shrinking an $\mathsf{M^S}_{001}$ and an $\mathsf{M^T}_{001}$ before fusing them, we have 
\begin{align}
	\mathcal{S} \left( \mathsf{M^S}_{001} \right)\otimes\mathcal{S} \left( \mathsf{M^T}_{001} \right)=&\left(2\cdot\mathsf{1}\right)\otimes\left( \mathsf{1}\oplus \mathsf{L}_{100}\oplus \mathsf{L}_{010}\oplus \mathsf{L}_{110} \right)\nonumber
	\\
	=&2\cdot\left(\mathsf{1}\oplus \mathsf{L}_{100}\oplus \mathsf{L}_{010}\oplus \mathsf{L}_{110}\right) =\mathcal{S} \left( \mathsf{M^{ST}} \right)\,.
\end{align}
Similarly, we have 
\begin{align}
	\mathcal{S}^2 \left( \mathsf{M^S}_{001} \right)\otimes\mathcal{S}^2 \left( \mathsf{M^T}_{001} \right)=&\left(2\cdot\mathsf{1}\right)\otimes\left(4\cdot\mathsf{1}\right)=8\cdot\mathsf{1}=\mathcal{S}^2 \left( \mathsf{M^{ST}} \right)
\end{align}
This shows that our definition for shrinking an $\mathsf{M^{ST}}$-excitation respect the fusion rules.

\subsection{Shrinking table}
All shrinking rules for the $BBA$ twisted term are shown in table~\ref{table_shrinking}. We can see that all pure loops are shrunk to vacuum, and pure $S^2$ membrane is shrunk to $2$ vacua. Pure $T^2$ membrane can be ultimately shrunk to $4$ vacua, $\mathsf{M^{ST}}$ ultimately shrunk to $8$ vacua. Thus all pure spatially extended topological excitations eventually shrunk to a vacuum or several vacua. We can also verify that quantum dimensions are preserved in the shrinking process. 
\begin{table}
	\caption{\label{table_shrinking}Shrinking table  for the $5$D twisted $BF$ theory with the $BBA$ twisted term. }
	\centering
	\scriptsize
	\begin{tabular*}{\textwidth}{@{\extracolsep{\fill}}|c|c|c|c|c|c|c|c|c|c|c|c|c|}
		\hline
		excitation & $\mathsf{L}_{100}$ & $\mathsf{L}_{010}$ & $\mathsf{L}_{110}$ & $\mathsf{L}_{100}^{001}$ & $\mathsf{L}_{010}^{001}$ & $\mathsf{L}_{110}^{001}$ & $\mathsf{M^S}_{001}$ & $\mathsf{M^S}_{001}^{001,}$ & $\mathsf{M^T}_{001}$ & $\mathsf{M^T}_{001}^{001,}$ & $\mathsf{M^{ST}}$ & $\mathsf{M^{ST}}^{001,}$ \\
		\hline
		$\mathcal{S} \left( \text{excitation}\right) $ & $\mathsf{1}$ & $\mathsf{1}$ & $\mathsf{1}$ & $\mathsf{P}_{001}$ & $\mathsf{P}_{001}$ & $\mathsf{P}_{001}$ & $2\cdot\mathsf{1}$ & $2\cdot \mathsf{P}_{001}$ & \makecell{$\mathsf{1}$ \\ $\oplus \mathsf{L}_{100}$ \\ $\oplus \mathsf{L}_{010}$ \\ $\oplus \mathsf{L}_{110}$} & \makecell{$\mathsf{P}_{001}$ \\ $\oplus \mathsf{L}_{100}^{001}$ \\ $\oplus \mathsf{L}_{010}^{001}$ \\ $\oplus \mathsf{L}_{110}^{001}$} & \makecell{$2\cdot\mathsf{1}$ \\ $\oplus2\cdot \mathsf{L}_{100}$ \\ $\oplus2\cdot \mathsf{L}_{010}$ \\ $\oplus2\cdot \mathsf{L}_{110}$} & \makecell{$2\cdot \mathsf{P}_{001}$ \\ $\oplus2\cdot \mathsf{L}_{100}^{001}$ \\ $\oplus2\cdot \mathsf{L}_{010}^{001}$ \\ $\oplus2\cdot \mathsf{L}_{110}^{001}$} \\
		\hline
		$\mathcal{S} ^2\left( \text{excitation}\right) $ & $\mathsf{1}$ & $\mathsf{1}$ & $\mathsf{1}$ & $\mathsf{P}_{001}$ & $\mathsf{P}_{001}$ & $\mathsf{P}_{001}$ & $2\cdot\mathsf{1}$ & $2\cdot \mathsf{P}_{001}$ & $4\cdot\mathsf{1}$ & $4\cdot \mathsf{P}_{001}$ & $8\cdot\mathsf{1}$ & $8\cdot \mathsf{P}_{001}$ \\
		\hline
	\end{tabular*}
\end{table}
\section{Anomaly-free algebraic structure in fusion rules and shrinking rules}\label{section_algebraic}

It is fundamentally important to study the algebraic structure encoded in topological data that characterize topological orders. Especially, all data, e.g., fusion coefficients and shrinking coefficients, should be consistent to each other such that they can together define an anomaly-free topological order. Otherwise, we expect that there should be a quantum anomaly of some kind that obstructs the existence of such topological order. In this section, we study the consistency relations between fusion and shrinking coefficients with the help of the results of fusion and shrinking rules obtained through the path-integral approach in the previous sections. 


Using table~\ref{table_fusionBBA} and table~\ref{table_shrinking}, we can verify that shrinking rules and hierarchical shrinking rules all respect the fusion rules:
\begin{gather}
	\mathcal{S} \left( \mathsf{a}\otimes \mathsf{b} \right) =\mathcal{S} \left( \mathsf{a} \right) \otimes \mathcal{S} \left( \mathsf{b} \right) \,,
	\mathcal{S}^2 \left( \mathsf{a}\otimes \mathsf{b} \right) =\mathcal{S}^2 \left( \mathsf{a} \right) \otimes \mathcal{S}^2 \left( \mathsf{b} \right)\,.
	\label{eq19}
\end{gather}

Eq.~(\ref{eq19}) also implies that there are some algebraic relationships between the fusion coefficient and shrinking coefficient. Using eq.~(\ref{eq_general_fusion}) and eq.~(\ref{eq_shrinking_pathintegral}), we first calculate $\mathcal{S} \left( \mathsf{a}\otimes \mathsf{b} \right)$:
\begin{align}
	\mathcal{S} \left( \mathsf{a}\otimes \mathsf{b} \right) =&\mathcal{S} \left( \oplus _{\mathsf{i}}N_{\mathsf{i}}^{\mathsf{ab}}\mathsf{i} \right) =\oplus _{\mathsf{i}}N_{\mathsf{i}}^{\mathsf{ab}}\mathcal{S} \left( \mathsf{i} \right) =\oplus _{\mathsf{i}}N_{\mathsf{i}}^{\mathsf{ab}}\left( \oplus _{\mathsf{c}}\mathrm{S}_{\mathsf{c}}^{\mathsf{i}}\mathsf{c} \right) 
	=\oplus _{\mathsf{c}}\left( \sum_{\mathsf{i}}{\mathrm{S}_{\mathsf{c}}^{\mathsf{i}}N_{\mathsf{i}}^{\mathsf{ab}}} \right) \mathsf{c}	\,.
	\label{eq15}
\end{align}
Keep shrinking, we will get $\mathcal{S} ^2\left( \mathsf{a}\otimes \mathsf{b} \right)$:
\begin{align}
	\mathcal{S} ^2\left( \mathsf{a}\otimes \mathsf{b} \right) =&\mathcal{S} \left[ \oplus _{\mathsf{j}}\left( \sum_{\mathsf{i}}{\mathrm{S}_{\mathsf{j}}^{\mathsf{i}}N_{\mathsf{i}}^{\mathsf{ab}}} \right) \mathsf{j} \right] =\oplus _{\mathsf{j}}\left[ \left( \sum_{\mathsf{i}}{\mathrm{S}_{\mathsf{j}}^{\mathsf{i}}N_{\mathsf{i}}^{\mathsf{ab}}} \right) \mathcal{S} \left( \mathsf{j} \right) \right] 	=\oplus _{\mathsf{j}}\left[ \left( \sum_{\mathsf{i}}{\mathrm{S}_{\mathsf{j}}^{\mathsf{i}}N_{\mathsf{i}}^{\mathsf{ab}}} \right) \left( \oplus _{\mathsf{c}}\mathrm{S}_{\mathsf{c}}^{\mathsf{j}}\mathsf{c} \right) \right] \nonumber\\
	=&\oplus _{\mathsf{c}}\left( \sum_{\mathsf{ij}}{\mathrm{S}_{\mathsf{c}}^{\mathsf{j}}\mathrm{S}_{\mathsf{j}}^{\mathsf{i}}N_{\mathsf{i}}^{a\mathsf{b}}} \right) \mathsf{c}\,.
	\label{eq16}
\end{align}
Now we calculate $\mathcal{S} \left( \mathsf{a} \right) \otimes \mathcal{S} \left( \mathsf{b} \right)$:
\begin{align}
	\mathcal{S} \left( \mathsf{a} \right) \otimes \mathcal{S} \left( \mathsf{b} \right) =&\left( \oplus _{\mathsf{i}}\mathrm{S}_{\mathsf{i}}^{\mathsf{a}}\mathsf{i} \right) \otimes \left( \oplus _{\mathsf{j}}\mathrm{S}_{\mathsf{j}}^{\mathsf{b}}\mathsf{j} \right) =\oplus _{\mathsf{i}}\oplus _{\mathsf{j}}\left[ \mathrm{S}_{\mathsf{i}}^{\mathsf{a}}\mathrm{S}_{\mathsf{j}}^{\mathsf{b}}\left( \mathsf{i}\otimes \mathsf{j} \right) \right] =\oplus _{\mathsf{i}}\oplus _{\mathsf{j}}\left[ \mathrm{S}_{\mathsf{i}}^{\mathsf{a}}\mathrm{S}_{\mathsf{j}}^{\mathsf{b}}\left( \oplus _{\mathsf{c}}N_{\mathsf{c}}^{\mathsf{ij}}\mathsf{c} \right) \right] 	\nonumber\\
	=&\oplus _{\mathsf{c}}\left( \sum_{\mathsf{ij}}{\mathrm{S}_{\mathsf{i}}^{\mathsf{a}}\mathrm{S}_{\mathsf{j}}^{\mathsf{b}}N_{\mathsf{c}}^{\mathsf{ij}}} \right) \mathsf{c}\,.
	\label{eq17}
\end{align}
Similarly, we obtain $\mathcal{S} ^2\left( \mathsf{a} \right) \otimes \mathcal{S} ^2\left( \mathsf{a} \right)$:
\begin{align}
	\mathcal{S} ^2\left( \mathsf{a} \right) \otimes \mathcal{S} ^2\left( \mathsf{b} \right) =&\mathcal{S} \left( \oplus _{\mathsf{i}}\mathrm{S}_{\mathsf{i}}^{\mathsf{a}}\mathsf{i} \right) \otimes \mathcal{S} \left( \oplus _{\mathsf{j}}\mathrm{S}_{\mathsf{j}}^{\mathsf{b}}\mathsf{j} \right) =\left[ \oplus _{\mathsf{k}}\left( \sum_{\mathsf{i}}{\mathrm{S}_{\mathsf{i}}^{\mathsf{a}}\mathrm{S}_{\mathsf{k}}^{\mathsf{i}}} \right) \mathsf{k} \right] \otimes \left[ \oplus _{\mathsf{l}}\left( \sum_{\mathsf{j}}{\mathrm{S}_{\mathsf{j}}^{\mathsf{b}}\mathrm{S}_{\mathsf{l}}^{\mathsf{j}}} \right) \mathsf{l} \right] \nonumber
	\\
	=&\oplus _{\mathsf{k}}\oplus _{\mathsf{l}}\left[ \left( \sum_{\mathsf{ij}}{\mathrm{S}_{\mathsf{i}}^{\mathsf{a}}\mathrm{S}_{\mathsf{k}}^{\mathsf{i}}\mathrm{S}_{\mathsf{j}}^{\mathsf{b}}\mathrm{S}_{\mathsf{l}}^{\mathsf{j}}} \right) \left( \mathsf{k}\otimes \mathsf{l} \right) \right] =\oplus _{\mathsf{k}}\oplus _{\mathsf{l}}\left[ \left( \sum_{\mathsf{ij}}{\mathrm{S}_{\mathsf{i}}^{\mathsf{a}}\mathrm{S}_{\mathsf{k}}^{\mathsf{i}}\mathrm{S}_{\mathsf{j}}^{\mathsf{b}}\mathrm{S}_{\mathsf{l}}^{\mathsf{j}}} \right) \left( \oplus _{\mathsf{c}}N_{\mathsf{c}}^{\mathsf{kl}}\mathsf{c} \right) \right] \nonumber
	\\
	=&\oplus _{\mathsf{c}}\left( \sum_{\mathsf{ijkl}}{\mathrm{S}_{\mathsf{i}}^{\mathsf{a}}\mathrm{S}_{\mathsf{k}}^{\mathsf{i}}\mathrm{S}_{\mathsf{j}}^{\mathsf{b}}\mathrm{S}_{\mathsf{l}}^{\mathsf{j}}}N_{\mathsf{c}}^{\mathsf{kl}} \right) \mathsf{c}\,.
	\label{eq18}
\end{align}

Comparing eq.~(\ref{eq15}) and eq.~(\ref{eq17}), eq.~(\ref{eq16}) and eq.~(\ref{eq18}), we obtain
\begin{gather}
	\sum_{\mathsf{i}}{\mathrm{S}_{\mathsf{c}}^{\mathsf{i}}N_{\mathsf{i}}^{\mathsf{ab}}}=\sum_{\mathsf{ij}}{\mathrm{S}_{\mathsf{i}}^{\mathsf{a}}\mathrm{S}_{\mathsf{j}}^{\mathsf{b}}N_{\mathsf{c}}^{\mathsf{ij}}}\,,\label{eq20-}
	\\ \sum_{\mathsf{ij}}{\mathrm{S}_{\mathsf{c}}^{\mathsf{j}}\mathrm{S}_{\mathsf{j}}^{\mathsf{i}}N_{\mathsf{i}}^{a\mathsf{b}}}=\sum_{\mathsf{ijkl}}{\mathrm{S}_{\mathsf{i}}^{\mathsf{a}}\mathrm{S}_{\mathsf{k}}^{\mathsf{i}}\mathrm{S}_{\mathsf{j}}^{\mathsf{b}}\mathrm{S}_{\mathsf{l}}^{\mathsf{j}}}N_{\mathsf{c}}^{\mathsf{kl}}\,.
	\label{eq20}
\end{gather}
Eq.~(\ref{eq20-}) describes the first shrinking process and $\mathsf{c}$ can be either loops or particles. Eq.~(\ref{eq20}) describes the second shrinking process and $\mathsf{c}$ must be particles. 

Notice that eq.~(\ref{eq20-}) and eq.~(\ref{eq20}) are not independent, eq.~(\ref{eq20}) can also be derived from eq.~(\ref{eq20-}) directly by acting a shrinking operator on both sides of eq.~(\ref{eq20-}). Only eq.~(\ref{eq20-}) is the fundamental relationship between fusion and shrinking coefficients and eq.~(\ref{eq20}) can be seen as the descendant of eq.~(\ref{eq20-}). In a topological order higher than $5$D, we may derive other descendants describing hierarchical shrinking processes from eq.~(\ref{eq20-}) by similarly acting more shrinking operators.




\section{General discussions on all twisted terms in 5D\label{section_general_discuss_twisted5D}}
In the above analysis, we have only considered the twisted $BF$ theory with the  $BBA$ twisted term, but there are many other legitimate twisted terms in $5$D topological orders, e.g., $AAAB$, $AAAAA$, $AAAdA$, $AAC$, $AdAdA$, $AdAB$, and $AAdB$. If gauge-invariant combinations among them are taken into account, there are even more twisted $BF$ theories as shown in ref.~\cite{Zhang:2021ycl}.  In principle, we can construct equivalence classes of  Wilson operators, and compute fusion and shrinking rules explicitly. For example, we find $AAAB$ also simultaneously supports non-Abelian fusion rules, non-Abelian shrinking rules and hierarchical shrinking rules (see appendix~\ref{ap4}). As shown in appendix~\ref{ap2}, we find that a simple method allows us to determine the general properties of fusion and shrinking rules from gauge transformations and some simple Wilson operators. We summarize these properties in table~\ref{table_comparison_5dtwisted} and discuss this table below. 
\begin{table}
	\caption{\label{table_comparison_5dtwisted}Comparision of fusion and shrinking rules in different twisted terms.}
	\centering
	\begin{tabular*}{\textwidth}{@{\extracolsep{\fill}}|c|c|c|c|}
		\hline
		twisted term & fusion rules & shrinking rules  & hierarchical shrinking rules \\
		\hline
		$BBA$ & non-Abelian & non-Abelian & yes \\
		\hline
		$AAAB$ & non-Abelian & non-Abelian & yes \\
		\hline
		$AAAAA$ & non-Abelian & non-Abelian & no \\
		\hline
		$AAAdA$ & non-Abelian & non-Abelian & no \\
		\hline
		$AAC$ & non-Abelian & non-Abelian & no \\
		\hline
		$AdAdA$ & Abelian & Abelian & no \\
		\hline
		$AdAB$ & Abelian & Abelian & no \\
		\hline
		$AAdB$ & Abelian & Abelian & no \\
		\hline
	\end{tabular*}
\end{table}
\begin{itemize}
\item Both $BBA$ and $AAAB$ twisted terms simultaneously support non-Abelian fusion rules, non-Abelian shrinking rules, and hierarchical shrinking rules. According to the discussion in appendix~\ref{ap2}, delta functionals carried by Wilson operators determine the type of fusion and shrinking rules. We find that in $BBA$ and $AAAB$, Wilson operators for membranes carry $\delta \left( \int_{\sigma}{B^i} \right) $, which will lead to hierarchical shrinking rules. The presence of hierarchical shrinking rules indicates that some fusion and shrinking rules are non-Abelian.

\item $AAAAA$, $AAAdA$, and $AAC$ support non-Abelian fusion rules and non-Abelian shrinking rules but they do not support hierarchical shrinking rules. As studied in ref.~\cite{Zhang:2021ycl}, the twisted $BF$ theories with these twisted terms contain only type-I $BF$ term, i.e., $CdA$. Since there are no $B^i$ fields, none of membranes carry $\delta \left( \int_{\sigma}{B^i} \right) $ and thus they do not have hierarchical shrinking. A pure $T^2$ membrane can only be shrunk to a trivial loop (vacuum). Although $\delta \left( \int_{\sigma}{B^i} \right) $ is absent, some Wilson operators for membranes carry $\delta \left( \int_{\gamma}{A^i} \right) $, which can lead to non-Abelian fusion and shrinking rules.

\item $AdAdA$, $AdAB$, and $AAdB$ twisted terms only support Abelian fusion and Abelian shrinking rules. Some Wilson operators may contain nontrivial terms besides $A^i$ or $B^i$ or $C^i$ to compensate for the nontrivial shift in gauge transformation, but they do not carry delta functionals that potentially lead to non-Abelian fusion rules and non-Abelian shrinking rules.
\end{itemize}



The presence of hierarchical shrinking leads to an interesting phenomenon: $S^2$ and $T^2$ membranes have different quantum dimensions. In a theory without hierarchical shrinking, they only differ by shape. Remember we request the fusion coefficients and the shrinking coefficients are all integers, this affects the normalization factors of Wilson operators. In the above analysis, we have seen that an $S^2$ membrane and a $T^2$ membrane carry the same gauge charges and have different normalization factors as long as the $T^2$ membrane can shrink to nontrivial loops. Thus they have different quantum dimensions. This is absent in the theory without hierarchical shrinking, since the $S^2$ and $T^2$ membranes carry the same gauge charges and have identical normalization factors. To be more specific, each gauge charge of $C^i$ field contributes an extra $2$ to the normalization factor of a $T^2$ membrane. This can be clearly viewed from $AAAB$ twisted term shown in appendix~\ref{ap4}. From eq.~(\ref{eq31}), we can see when the membrane carries $n$ units of gauge charges of $C^i$ field, then we have $\tilde{N}_{abc}=2^nN_{abc}$. This result also applies to the $BBA$ twisted term since there are only $\mathbb{Z}_{N_3}$ membranes and we do find that $\mathsf{M^T}_{001}$ has normalization factor $4$ while $\mathsf{M^S}_{001}$ has $2$. 



\section{Summary and outlook}\label{section_summary}

In this paper, as a series of work on $5$D topological orders, we   concentrated on the fusion and shrinking rules of $5$D topological orders by means of the path-integral formalism of TQFTs proposed in ref.~\cite{Zhang:2021ycl}. More concretely, we constructed Wilson operators for all topologically distinct excitations by carefully investigating equivalence classes among all possible gauge-invariant Wilson operators.  
Based on Wilson operators, we may derive fusion rules, quantum dimensions, and shrinking rules. We found exotic non-Abelian fusion and non-Abelian shrinking, and especially identified exotic hierarchical structure in shrinking processes.  The underlying algebraic structure of the fusion coefficients and shrinking coefficients was also found in terms of several exact algebraic equations about the coefficients, which sets an anomaly-free condition on $5$D topological orders.  Below we present several future directions motivated by this paper.  (i) By adopting the strategy in ref.~\cite{Ning2018prb}, we may explore symmetry-enrichment in $5$D topological orders. While the idea of ``mixed $3$-loop statistics'' plays an important role in charactering and classifying SET phases in ref.~\cite{Ning2018prb}, it will be interesting to generalize such mixed type of braiding statistics to probe SET order in higher dimensions, especially when higher-form symmetry~\cite{mcgreevy2023} is considered and higher-form response gauge fields are involved.
(ii) We have only considered topological excitations in a simple shape such as particles, simple unknotted loops, $S^2$ membranes,   $T^2$ membranes, and composite of the two membranes. It is important to explore more complicated spatially extended excitations, e.g., $T^2$ membranes with higher genus in order to see if there are more exotic topological data that should be involved for completeness. 
(iii) Finally, it will be interesting to explore fermionic topological orders by putting TQFTs on a spacetime manifold with spin structure.
  

\acknowledgements
This work was supported by NSFC Grant  No.~12074438, and Guangdong Basic and Applied Basic Research Foundation under Grant No.~2020B1515120100.

\appendix
\section{Nontrivial terms in Wilson operators\label{ap1}}
When we construct the Wilson operators, we have to make sure they are gauge invariant. If the gauge field transforms with nontrivial shifts, we will have to add some nontrivial terms to compensate for these nontrivial shifts in Wilson operators. For example in $BBA$ twisted term, $C^3$ transforms with nontrivial shifts
\begin{align}
	C^3\rightarrow C^3+dT^3+\frac{2\pi q}{N_3}\left( -V^1B^2-B^1V^2-\frac{1}{2}V^1dV^2-\frac{1}{2}V^2dV^1 \right) 
\end{align}
thus we construct the Wilson operator for $\mathsf{M^S}_{001}$ as 
\begin{align}
	\mathsf{M^S}_{001}=2\exp \left[ i\int_{\omega}{C^3+\frac{1}{2}\frac{2\pi q}{N_3}\left( d^{-1}B^1B^2+d^{-1}B^2B^1 \right)} \right] \delta \left( \int_{\sigma}{B^1} \right) \delta \left( \int_{\sigma}{B^2} \right) 
\end{align}
where $\frac{1}{2}\frac{2\pi q}{N_3}\left( d^{-1}B^1B^2+d^{-1}B^2B^1 \right)$ is introduce to compensate for the nontrivial shifts.

When we consider $\langle \mathsf{M^S}_{001}\otimes \mathsf{M^S}_{001} \rangle $, we will encounter 
\begin{align}
	&\exp \left[ i2\int_{\omega}{\frac{1}{2}\frac{2\pi}{N_3}\frac{pN_1N_2N_3}{\left( 2\pi \right) ^2N_{123}}\left( d^{-1}\hat{B}^1\hat{B}^2+d^{-1}\hat{B}^2\hat{B}^1 \right)} \right] \nonumber
	\\
	=&\exp \left[ i2\times\frac{1}{2}\frac{2\pi}{N_3}\frac{pN_1N_2N_3}{\left( 2\pi \right) ^2N_{123}}\left( \frac{2\pi k_1}{N_1}\frac{2\pi m_2}{N_2}+\frac{2\pi m_1}{N_1}\frac{2\pi k_2}{N_2} \right) \right] \nonumber
	\\
	=&\exp \left[ i\pi p\left( k_1m_2+m_1k_2 \right) \right] 
\end{align}
At the first sight, $\exp \left[ i\pi p\left( k_1m_2+m_1k_2 \right) \right]$ can be $+1$ or $-1$ since $k_1,k_2,m_1,m_2$ are all integers. However we can find that $k_1m_2$ is actually equal to $m_1k_2$, because $d^{-1}B^1B^2$ is equal to $d^{-1}B^2B^1$ up to a boundary term. Notice that the delta functionals enforce $\int_{\sigma}{B^1}$ and $\int_{\sigma}{B^2}$ are $0$ mod $2\pi$, thus $B^1$ and $B^2$ are exact. We can find two gauge fields $F_1$ and $F_2$ satisfy that $dF_1=B^1$ and $dF_2=B^2$, thus we have
\begin{align}
	d^{-1}B^2B^1=F_2dF_1=dF_2F_1-d\left( F_2F_1 \right) =F_1dF_2=d^{-1}B^1B^2
\end{align}
We drop the boundary term because it will not contribute to the integral. Now we can see that $\frac{2\pi k_1}{N_1}\frac{2\pi m_2}{N_2}=\frac{2\pi m_1}{N_1}\frac{2\pi k_2}{N_2}$ and thus $k_1m_2=m_1k_2$. Finally, we obtain 
\begin{align}
	\exp \left[ i\pi p\left( k_1m_2+m_1k_2 \right) \right] =1
\end{align}

Similarly, in $AAAB$ twisted term, we construct the gauge invariant Wilson operator for $\mathsf{M^S}_{100}$ as 
\begin{align}
	\mathsf{M^S}_{100}=&4\exp \left[ i\int_{\omega}{C^1-\frac{1}{3}\frac{2\pi q}{N_1}\left( B^4A^2d^{-1}A^3-B^4A^3d^{-1}A^2-A^2A^3d^{-1}B^4 \right)} \right] \nonumber
	\\
	&\times \delta \left( \int_{\gamma}{A^2} \right) \delta \left( \int_{\gamma}{A^3} \right) \delta \left( \int_{\sigma}{B^4} \right) 
\end{align}
Due to delta functionals, we can see that $B^4$, $A^2$ and $A^3$ are all exact and we can find $dF_4=B^4$, $dF_2=A^2$ and $dF_3=A^3$. Thus
\begin{align}
	B^4A^2d^{-1}A^3=&dF_4dF_2F_3=d\left( F_4dF_2 \right) F_3=d\left( F_4dF_2F_3 \right) -F_4dF_2dF_3 \nonumber
	\\
	=&-dF_2dF_3F_4=-A^2A^3d^{-1}B^4 \nonumber
	\\
	B^4A^3d^{-1}A^2=&dF_4dF_3F_2=d\left( F_4dF_3 \right) F_2=d\left( F_4dF_3F_2 \right) -F_4dF_3dF_2 \nonumber
	\\
	=&dF_2dF_3F_4=A^2A^3d^{-1}B^4
\end{align}
We obtain that $B^4A^2d^{-1}A^3=-B^4A^3d^{-1}A^2=-A^2A^3d^{-1}B^4$ after dropping some boundary terms. When we consider fusing two $\mathsf{M^S}_{100}$, we will encounter 
\begin{align}
	&\exp \left[ i2\int_{\omega}{\frac{1}{3}\frac{2\pi}{N_1}\frac{pN_1N_2N_3N_4}{\left( 2\pi \right) ^3N_{1234}}\left( \hat{B}^4\hat{A}^2d^{-1}\hat{A}^3-\hat{B}^4\hat{A}^3d^{-1}\hat{A}^2-\hat{A}^2\hat{A}^3d^{-1}\hat{B}^4 \right)} \right] \nonumber
	\\
	=&\exp \left[ i2\times \frac{1}{3}\frac{2\pi}{N_1}\frac{pN_1N_2N_3N_4}{\left( 2\pi \right) ^3N_{1234}} \right. \nonumber
	\\
	&\left. \times \left( \frac{2\pi m_1}{N_4}\frac{2\pi m_2}{N_2}\frac{2\pi k_3}{N_3}-\frac{2\pi m_4}{N_4}\frac{2\pi m_3}{N_3}\frac{2\pi k_2}{N_2}-\frac{2\pi m_2}{N_2}\frac{2\pi m_3}{N_3}\frac{2\pi k_4}{N_4} \right) \right] \nonumber
	\\
	=&\exp \left[ \frac{i2\pi p}{3}\left( m_1m_2k_3-m_4m_3k_2-m_2m_3k_4 \right) \right] \nonumber
	\\
	=&\exp \left[ \frac{i2\pi p}{3}\left( 3m_1m_2k_3\right) \right]=1
\end{align}
For other excitations, we have similar calculations. 

In general, when calculating fusion rules, we have $\langle \exp \left[ i2\int_{X}{f^i+\left( \cdots \right)} \right] \rangle =1$ in our theory, where $X$ is the world-line(sheet or volume) and $f^i$ is the gauge field. 
\section{Fusion and shrinking rules for the $AAAB$ twisted term\label{ap4}}
We have already seen that $BBA$ twisted term can present non-Abelian fusion and shrinking rules and hierarchical shrinking rules. $AAAB$ twisted term can also present these properties. In this appendix, we will study fusion and shrinking rules for the $AAAB$ term. Compared to the $BBA$ twisted term, pure spatially extended topological excitations are not necessarily shrunk to vacua eventually in the $AAAB$ term. Besides, we can clearly see that for $T^2$ membranes, each gauge charge of $C^i$ field contributes an extra $2$ factor to the quantum dimension compared to $S^2$ membranes.

When the gauge group is $G={\left( \mathbb{Z} _{2} \right)}^4$, we find that there are $100$ nonequivalent excitations, which makes fusion and shrinking rules complicated and hard to present all of them in a table. So we will use several formulas to describe them. One can derive the fusion and shrinking rules for any excitation following our formulas. Since we deal with the $AAAB$ term in a similar way to the $BBA$ term, we will only show our results. 
\subsection{TQFT action and gauge transformations for $AAAB$ twisted term}
For the $AAAB$ twisted term, the index of all $A$ and $B$ fields must be different. Thus the simplest case is take the gauge group $G={\left( \mathbb{Z} _{2} \right)}^4$. Then we can write down the TQFT action:
\begin{align}
	S=\int{\sum_{i=1}^3{\frac{N_i}{2\pi}C^idA^i}+}\frac{N_4}{2\pi}\tilde{B}^4dB^4+qA^1A^2A^3B^4
\end{align}
where $q=\frac{pN_1N_2N_3N_4}{\left( 2\pi \right) ^3N_{1234}}$, $p\in \mathbb{Z} _{N_{1234}}$, $N_{1,2,3,4}=N_{1234}=2$. The gauge transformations are:
\begin{align}
	&A^i\rightarrow A^i+d\chi ^i,B^4\rightarrow B^4+dV^4
	\label{eq21}
	\\
	&C^i\rightarrow C^i+dT^i+\frac{2\pi q}{N_i}\sum_{j,k}{\epsilon ^{ijk}}\left[ \left( A^j\chi ^k-\frac{1}{2}\chi ^jd\chi ^k \right) \left( B^4+dV^4 \right) -\frac{1}{2}A^jA^kV^4 \right] 
	\\
	&\tilde{B}^4\rightarrow \tilde{B}^4+d\tilde{V}^4+\frac{2\pi q}{N_4}\sum_{i,j,k}{\epsilon ^{4ijk}}\left( -\frac{1}{2}\chi ^iA^jA^k+\frac{1}{2}A^i\chi ^jd\chi ^k-\frac{1}{6}\chi ^id\chi ^jd\chi ^k \right) 
\end{align}
where $\epsilon$ is the antisymmetric tensor and $i,j,k=1,2,3$.
\subsection{Nonequivalent Wilson operators for excitations}
Now we can construct Wilson operators for different excitations. We will illustrate it in $3$ different cases: Abelian excitations(particles and loops), membranes, and $\mathsf{M^{ST}}$-excitations.

\textbf{\textit{particles and loops.}}
All particles and loops are Abelian excitations because we can see that in eq.~(\ref{eq21}), $A^i$ and $B^4$ transform normally (without nontrivial shifts). Thus we can easily write down their Wilson operators in a general way:
\begin{align}
	\mathsf{L}_{000d}^{abc0}=\exp \left( id\int_{\sigma}{B^4}+ia\int_{\gamma}{A^1}+ib\int_{\gamma}{A^2}+ic\int_{\gamma}{A^3} \right)
	\label{eq7}
\end{align}
where $a,b,c,d=0,1$ and they take value independently. Notice that when $d=0$, $\mathsf{L}_{000d}^{abc0}=\mathsf{P}_{abc0}$, when $a=b=c=d=0$, $\mathsf{L}_{000d}^{abc0}=\mathsf{1}$. For simplicity, we use $\mathsf{P}_{abc}$ to denote $\mathsf{P}_{abc0}$ in the following discussion. From eq.~(\ref{eq7}), we can find that there are $2^4=16$ nonequivalent excitations. 

\textbf{\textit{$S^2$ and $T^2$ membranes.}}
All membranes are non-Abelian excitations. We give a general formula for all nonequivalent $S^2$ membranes first, they can be written as 
\begin{align}
	\mathsf{M^S}_{abc}^{a^{\prime}b^{\prime}c^{\prime},}=&N_{abc}\exp \left[ ia^{\prime}\int_{\gamma}{A^1}+ib^{\prime}\int_{\gamma}{A^2}+ic^{\prime}\int_{\gamma}{A^3} \right. \nonumber
	\\
	&+ia\int_{\omega}{C^1-\frac{1}{3}\frac{2\pi q}{N_1}\left( B^4A^2d^{-1}A^3-B^4A^3d^{-1}A^2-A^2A^3d^{-1}B^4 \right)} \nonumber
	\\
	&+ib\int_{\omega}{C^2+\frac{1}{3}\frac{2\pi q}{N_2}\left( B^4A^1d^{-1}A^3-B^4A^3d^{-1}A^1-A^1A^3d^{-1}B^4 \right)} \nonumber
	\\
	&+\left. ic\int_{\omega}{C^3-\frac{1}{3}\frac{2\pi q}{N_3}\left( B^4A^1d^{-1}A^2-B^4A^2d^{-1}A^1-A^1A^2d^{-1}B^4 \right)} \right]  \nonumber
	\\
	&\times \delta \left[ \left( ab \right) \int_{\gamma}{A^2-A^1} \right] \delta \left[ \left( ac \right) \int_{\gamma}{A^3-A^1} \right] \delta \left[ \left( bc \right) \int_{\gamma}{A^3-A^2} \right] \nonumber
	\\
	&\times \delta \left[ \left( 1-a \right) \int_{\gamma}{A^1} \right] \delta \left[ \left( 1-b \right) \int_{\gamma}{A^2} \right] \delta \left[ \left( 1-c \right) \int_{\gamma}{A^3} \right] \delta \left( \int_{\sigma}{B^4} \right) 	
	\label{eq8}
\end{align}
For simplicity, we use $\mathsf{M^S}_{abc}^{a^{\prime}b^{\prime}c^{\prime},}$ to denote $\mathsf{M^S}_{abc0}^{a^{\prime}b^{\prime}c^{\prime}0,}$. Although eq.~(\ref{eq8}) only describes $S^2$ membranes, $T^2$ membranes have a similar expression. We just need to replace all $\mathsf{M^S}$, $\omega$ and $N_{abc}$ with $\mathsf{M^T}$, $\tilde{\omega}$ and $\tilde{N}_{abc}$ respectively in eq.~(\ref{eq8}), then we get the correct formula for $T^2$ membranes. In eq.~(\ref{eq8}), $N_{abc}$ is normalization factor, $a,b,c=0,1$ and at least one of $a,b,c$ is nonzero. So there are $7$ possible combinations:
\begin{gather}
	\left\{ \left( abc \right) \right\} =\left\{ \left( 100 \right) ,\left( 010 \right) ,\left( 001 \right) ,\left( 110 \right) ,\left( 101 \right) ,\left( 011 \right) ,\left( 111 \right) \right\}
	\label{eq31-}  
	\\
	N_{abc}=\left\{4,4,4,8,8,8,16\right\},\qquad\tilde{N}_{abc}=\left\{8,8,8,32,32,32,128\right\}
	\label{eq31}
\end{gather}
For a specific $\left( abc \right)$, we need to choose the corresponding normalization factor shown in eq.~(\ref{eq31}). Given a $\left( abc \right)$, the possible values of $\left( a^{\prime}b^{\prime}c^{\prime} \right)$ are also restricted. When $\left( abc \right)$ takes value in the order of eq.~(\ref{eq31-}), correspondingly we can choose
\begin{align}
	\left\{ \left( a^{\prime}b^{\prime}c^{\prime} \right) \right\} =\left\{ \left( a^{\prime}00 \right) ,\left( 0b^{\prime}0 \right) ,\left( 00c^{\prime} \right) ,\left( a^{\prime}00 \right) ,\left( a^{\prime}00 \right) ,\left( 0b^{\prime}0 \right) ,\left( a^{\prime}00 \right) \right\}  
	\label{eq9}
\end{align}
where $a^{\prime},b^{\prime},c^{\prime}=0,1$. We can see that once given a specific $\left( abc \right)$, two of $a^{\prime},b^{\prime},c^{\prime}$ will be forced to zero and hence there are only two possible $\left( a^{\prime}b^{\prime}c^{\prime} \right)$. For example, when we choose $\left( abc \right)=\left( 011 \right)$, then $\left( a^{\prime}b^{\prime}c^{\prime} \right)=\left( 010 \right)$ or $\left( 000 \right)$. Eq.~(\ref{eq9}) is not the only choice, other choices can also correctly represent these nonequivalent excitations, but these choices are equivalent to eq.~(\ref{eq9}). In the following discussion, we will fix our choice in eq.~(\ref{eq9}).

Eq.~(\ref{eq8}) contains $7\times2=14$ nonequivalent $S^2$ membranes since there are $7$ different $\left( abc \right)$ and each $\left( abc \right)$ correspond to $2$ nonequivalent $\left( a^{\prime}b^{\prime}c^{\prime} \right)$. For $T^2$ membranes we have a similar discussion, just replace all $\mathsf{M^S}$, $\omega$ and $N_{abc}$ with $\mathsf{M^T}$, $\tilde{\omega}$ and $\tilde{N}_{abc}$ respectively in eq.~(\ref{eq8}), we obtain $14$ nonequivalent $T^2$ membranes.

\textbf{\textit{$\mathsf{M^{ST}}$-excitations.}}
Similar to $BBA$ twisted term, $\mathsf{M^{ST}}$-excitations in $AAAB$ twisted term are also created by fusing $S^2$ membranes and  $T^2$ membranes. When we fuse an $\mathsf{M^S}_{abc}^{a^{\prime}b^{\prime}c^{\prime},}$ and an $\mathsf{M^T}_{\tilde{a}\tilde{b}\tilde{c}}^{\tilde{a}^{\prime}\tilde{b}^{\prime}\tilde{c}^{\prime},}$, we label the result as $\mathsf{M^{ST}}_{abc,\tilde{a}\tilde{b}\tilde{c}}^{\left( a^{\prime}+\tilde{a}^{\prime} \right) \left( b^{\prime}+\tilde{b}^{\prime} \right) \left( c^{\prime}+\tilde{c}^{\prime} \right),}$. Any addition that appears in subscript and superscript is addition mod $2$, which means $1+1=0$ and $0+1=1$. Notice that $\mathsf{M^S}_{abc}^{a^{\prime}b^{\prime}c^{\prime},}\otimes\mathsf{M^T}_{\tilde{a}\tilde{b}\tilde{c}}^{\tilde{a}^{\prime}\tilde{b}^{\prime}\tilde{c}^{\prime},}$ is different from $\mathsf{M^T}_{abc}^{a^{\prime}b^{\prime}c^{\prime},}\otimes \mathsf{M^S}_{\tilde{a}\tilde{b}\tilde{c}}^{\tilde{a}^{\prime}\tilde{b}^{\prime}\tilde{c}^{\prime},}$, so $\mathsf{M^{ST}}_{abc,\tilde{a}\tilde{b}\tilde{c}}^{\left( a^{\prime}+\tilde{a}^{\prime} \right) \left( b^{\prime}+\tilde{b}^{\prime} \right) \left( c^{\prime}+\tilde{c}^{\prime} \right) ,}$ is not equal to $\mathsf{M^{ST}}_{\tilde{a}\tilde{b}\tilde{c},abc}^{\left( a^{\prime}+\tilde{a}^{\prime} \right) \left( b^{\prime}+\tilde{b}^{\prime} \right) \left( c^{\prime}+\tilde{c}^{\prime} \right),}$ generally. We find that $\mathsf{M^{ST}}_{abc,\tilde{a}\tilde{b}\tilde{c}}^{\left( a^{\prime}+\tilde{a}^{\prime} \right) \left( b^{\prime}+\tilde{b}^{\prime} \right) \left( c^{\prime}+\tilde{c}^{\prime} \right),}$ with different $\left(abca^{\prime}b^{\prime}c^{\prime}\tilde{a}\tilde{b}\tilde{c}\tilde{a}^{\prime}\tilde{b}^{\prime}\tilde{c}^{\prime}\right)$ are not necessarily nonequivalent in our calculation. These $\mathsf{M^{ST}}$-excitations can be classified into $56$ equivalence classes and we are going to give two formulas to describe all nonequivalent $\mathsf{M^{ST}}$-excitations. They can be represented by:
\begin{align}
	\mathsf{M^{ST}}_{abc,abc}^{a^{\prime\prime}b^{\prime\prime}c^{\prime\prime},}=&N_{abc}\tilde{N}_{abc}\exp \left[ ia^{\prime\prime}\int_{\gamma}{A^1}+ib^{\prime\prime}\int_{\gamma}{A^2}+ic^{\prime\prime}\int_{\gamma}{A^3} \right. \nonumber 
	\\
	&+ia\int_{\omega}{C^1-\frac{1}{3}\frac{2\pi q}{N_1}\left( B^4A^2d^{-1}A^3-B^4A^3d^{-1}A^2-A^2A^3d^{-1}B^4 \right)} \nonumber
	\\
	&+ib\int_{\omega}{C^2+\frac{1}{3}\frac{2\pi q}{N_2}\left( B^4A^1d^{-1}A^3-B^4A^3d^{-1}A^1-A^1A^3d^{-1}B^4 \right)} \nonumber
	\\
	&+ic\int_{\omega}{C^3-\frac{1}{3}\frac{2\pi q}{N_3}\left( B^4A^1d^{-1}A^2-B^4A^2d^{-1}A^1-A^1A^2d^{-1}B^4 \right)} \nonumber
	\\
	&+ia\int_{\tilde{\omega}}{C^1-\frac{1}{3}\frac{2\pi q}{N_1}\left( B^4A^2d^{-1}A^3-B^4A^3d^{-1}A^2-A^2A^3d^{-1}B^4 \right)} \nonumber
	\\
	&+ib\int_{\tilde{\omega}}{C^2+\frac{1}{3}\frac{2\pi q}{N_2}\left( B^4A^1d^{-1}A^3-B^4A^3d^{-1}A^1-A^1A^3d^{-1}B^4 \right)} \nonumber
	\\
	&+\left. ic\int_{\tilde{\omega}}{C^3-\frac{1}{3}\frac{2\pi q}{N_3}\left( B^4A^1d^{-1}A^2-B^4A^2d^{-1}A^1-A^1A^2d^{-1}B^4 \right)} \right] \nonumber
	\\
	&\times \delta \left[ \left( ab \right) \int_{\gamma}{A^2-A^1} \right] \delta \left[ \left( ac \right) \int_{\gamma}{A^3-A^1} \right] \delta \left[ \left( bc \right) \int_{\gamma}{A^3-A^2} \right] \nonumber
	\\
	&\times \delta \left[ \left( 1-a \right) \int_{\gamma}{A^1} \right] \delta \left[ \left( 1-b \right) \int_{\gamma}{A^2} \right] \delta \left[ \left( 1-c \right) \int_{\gamma}{A^3} \right] \delta \left( \int_{\sigma}{B^4} \right) 
\end{align}
where $\left(abc\right)$ is equal to $\left(\tilde{a}\tilde{b}\tilde{c}\right)$, i.e., $a=\tilde{a},b=\tilde{b},c=\tilde{c}$. Possible values of $\left( a^{\prime\prime}b^{\prime\prime}c^{\prime\prime} \right)$ still follows eq.~(\ref{eq9}), so two of $a^{\prime\prime}$, $b^{\prime\prime}$ and $c^{\prime\prime}$ will be force to zero once $\left( abc \right)$ is given, the last one is determined by 
\begin{align}
	a^{\prime\prime}&=a^{\prime}+\tilde{a}^{\prime} \quad \text{mod} \quad 2 \nonumber 
	\\
	b^{\prime\prime}&=b^{\prime}+\tilde{b}^{\prime} \quad \text{mod} \quad 2 \nonumber 
	\\
	c^{\prime\prime}&=c^{\prime}+\tilde{c}^{\prime} \quad \text{mod} \quad 2 
\end{align}

When $\left(abc\right)$ is not equal to $\left(\tilde{a}\tilde{b}\tilde{c}\right)$, then we obtain
\begin{align}
	\mathsf{M^{ST}}_{abc,\tilde{a}\tilde{b}\tilde{c}}=&N_{abc}\tilde{N}_{\tilde{a}\tilde{b}\tilde{c}}\exp \left[ ia\int_{\omega}{C^1-\frac{1}{3}\frac{2\pi q}{N_1}\left( B^4A^2d^{-1}A^3-B^4A^3d^{-1}A^2-A^2A^3d^{-1}B^4 \right)} \right. \nonumber
	\\
	&+ib\int_{\omega}{C^2+\frac{1}{3}\frac{2\pi q}{N_2}\left( B^4A^1d^{-1}A^3-B^4A^3d^{-1}A^1-A^1A^3d^{-1}B^4 \right)} \nonumber
	\\
	&+ic\int_{\omega}{C^3-\frac{1}{3}\frac{2\pi q}{N_3}\left( B^4A^1d^{-1}A^2-B^4A^2d^{-1}A^1-A^1A^2d^{-1}B^4 \right)} \nonumber
	\\
	&+i\tilde{a}\int_{\tilde{\omega}}{C^1-\frac{1}{3}\frac{2\pi q}{N_1}\left( B^4A^2d^{-1}A^3-B^4A^3d^{-1}A^2-A^2A^3d^{-1}B^4 \right)} \nonumber
	\\
	&+i\tilde{b}\int_{\tilde{\omega}}{C^2+\frac{1}{3}\frac{2\pi q}{N_2}\left( B^4A^1d^{-1}A^3-B^4A^3d^{-1}A^1-A^1A^3d^{-1}B^4 \right)} \nonumber
	\\
	&+\left. i\tilde{c}\int_{\tilde{\omega}}{C^3-\frac{1}{3}\frac{2\pi q}{N_3}\left( B^4A^1d^{-1}A^2-B^4A^2d^{-1}A^1-A^1A^2d^{-1}B^4 \right)} \right] \nonumber
	\\
	&\times \delta \left( \int_{\gamma}{A^1} \right) \delta \left( \int_{\gamma}{A^2} \right) \delta \left( \int_{\gamma}{A^3} \right) \delta \left( \int_{\sigma}{B^4} \right) 
\end{align}
superscript is omitted since it is $\left(000,\right)$. 

We conclude that $\mathsf{M^{ST}}_{abc,abc}^{\left( a^{\prime}+\tilde{a}^{\prime} \right) \left( b^{\prime}+\tilde{b}^{\prime} \right) \left( c^{\prime}+\tilde{c}^{\prime} \right),}$ is equivalent to $\mathsf{M^{ST}}_{abc,abc}^{a^{\prime\prime}b^{\prime\prime}c^{\prime\prime},}$ when $\left(abc\right)=\left(\tilde{a}\tilde{b}\tilde{c}\right)$ and $\mathsf{M^{ST}}_{abc,\tilde{a}\tilde{b}\tilde{c}}^{\left( a^{\prime}+\tilde{a}^{\prime} \right) \left( b^{\prime}+\tilde{b}^{\prime} \right) \left( c^{\prime}+\tilde{c}^{\prime} \right),}$ is equivalent to $\mathsf{M^{ST}}_{abc,\tilde{a}\tilde{b}\tilde{c}}$ when $\left(abc\right)\ne\left(\tilde{a}\tilde{b}\tilde{c}\right)$. Now we obtain all equivalence classes. There are $16$ Abelian excitations and $14+14+56=84$ non-Abelian excitations.
\subsection{Fusion rules for the $AAAB$ twisted term}
The discussion above shows that any possible excitation can be represented by one of the following expressions:
\begin{align}
	\mathsf{L}_{000d}^{abc0} \qquad 	\mathsf{M^S}_{abc}^{a^{\prime}b^{\prime}c^{\prime},} \qquad \mathsf{M^T}_{\tilde{a}\tilde{b}\tilde{c}}^{\tilde{a}^{\prime}\tilde{b}^{\prime}\tilde{c}^{\prime},} \qquad \mathsf{M^{ST}}_{abc,abc}^{a^{\prime\prime}b^{\prime\prime}c^{\prime\prime},} \qquad \mathsf{M^{ST}}_{abc,\tilde{a}\tilde{b}\tilde{c}}
\end{align}
Thus if we want to derive all fusion rules for all nonequivalent excitations, we can try to obtain the fusion rules for these $5$ expressions. We can formally express the fusion table as table~\ref{table5}. Once we can explicitly write down the formulas of every element in table~\ref{table5}, we obtain the complete fusion rules.
\begin{table}
	\caption{\label{table5}Fusion rules for $AAAB$ twisted term. The fusion table is symmetric, so we only have to write down the upper triangular part of this table. To get all information about fusion rules, we have to know what exactly these elements are.}
	\centering
	\begin{tabular*}{\textwidth}{@{\extracolsep{\fill}}|c|c|c|c|c|c|}
		\hline
		& $\mathsf{L}_{000d}^{abc0}$ & $\mathsf{M^S}_{abc}^{a^{\prime}b^{\prime}c^{\prime},}$ & $\mathsf{M^T}_{\tilde{a}\tilde{b}\tilde{c}}^{\tilde{a}^{\prime}\tilde{b}^{\prime}\tilde{c}^{\prime},}$ & $\mathsf{M^{ST}}_{abc,abc}^{a^{\prime\prime}b^{\prime\prime}c^{\prime\prime},}$ & $\mathsf{M^{ST}}_{abc,\tilde{a}\tilde{b}\tilde{c}}$ \\
		\hline
		$\mathsf{L}_{000h}^{efg0}$ & $\mathsf{L}\otimes \mathsf{L}$ & $\mathsf{M^S}\otimes \mathsf{L}$ & $\mathsf{M^T}\otimes \mathsf{L}$ & $\mathsf{M^{ST}}_1\otimes \mathsf{L}$ & $\mathsf{M^{ST}}_2\otimes \mathsf{L}$ \\
		\hline
		$\mathsf{M^S}_{efg}^{e^{\prime}f^{\prime}g^{\prime},}$ && $\mathsf{M^S}\otimes \mathsf{M^S}$ & $\mathsf{M^T}\otimes \mathsf{M^S}$ & $\mathsf{M^{ST}}_1\otimes \mathsf{M^S}$ & $\mathsf{M^{ST}}_2\otimes \mathsf{M^S}$ \\
		\hline
		$\mathsf{M^T}_{\tilde{e}\tilde{f}\tilde{g}}^{\tilde{e}^{\prime}\tilde{f}^{\prime}\tilde{g}^{\prime},}$ &&& $\mathsf{M^T}\otimes \mathsf{M^T}$ & $\mathsf{M^{ST}}_1\otimes \mathsf{M^T}$ & $\mathsf{M^{ST}}_2\otimes \mathsf{M^T}$ \\
		\hline
		$\mathsf{M^{ST}}_{efg,efg}^{e^{\prime\prime}f^{\prime\prime}g^{\prime\prime},}$ &&&& $\mathsf{M^{ST}}_1\otimes \mathsf{M^{ST}}_1$ & $\mathsf{M^{ST}}_2\otimes \mathsf{M^{ST}}_1$ \\
		\hline
		$\mathsf{M^{ST}}_{efg,\tilde{e}\tilde{f}\tilde{g}}$ &&&&& $\mathsf{M^{ST}}_2\otimes \mathsf{M^{ST}}_2$ \\
		\hline
	\end{tabular*}
\end{table}

There are $15$ elements in table~\ref{table5}, we will show how to derive the fusion rules between them in $3$ different cases: fusion concerning Abelian excitations(particles and loops), fusion concerning two membranes and fusion concerning $\mathsf{M^{ST}}$-excitations.

\textbf{\textit{Fusion rules for Abelian excitations.}}
We can directly derive $5$ explicit formulas concerning Abelian fusion, they correspond to the $5$ elements in the first row of table~\ref{table5}:
\begin{align}
	&\mathsf{L}_{000d}^{abc0}\otimes \mathsf{L}_{000h}^{efg0}=\mathsf{L}_{000\left( d+h \right)}^{\left( a+e \right) \left( b+f \right) \left( c+g \right) 0}
	\\
	&\mathsf{M^S}_{abc}^{a^{\prime}b^{\prime}c^{\prime},}\otimes \mathsf{L}_{000h}^{efg0}=\mathsf{M^S}_{abc}^{F_{abc}\left[ \left( a^{\prime}+e \right) \left( b^{\prime}+f \right) \left( c^{\prime}+g \right) \right] ,}
	\\
	&\mathsf{M^T}_{\tilde{a}\tilde{b}\tilde{c}}^{\tilde{a}^{\prime}\tilde{b}^{\prime}\tilde{c}^{\prime},}\otimes \mathsf{L}_{000h}^{efg0}=\mathsf{M^T}_{\tilde{a}\tilde{b}\tilde{c}}^{F_{\tilde{a}\tilde{b}\tilde{c}}\left[ \left( \tilde{a}^{\prime}+e \right) \left( \tilde{b}^{\prime}+f \right) \left( \tilde{c}^{\prime}+g \right) \right] ,}
	\\
	&\mathsf{M^{ST}}_{abc,abc}^{a^{\prime\prime}b^{\prime\prime}c^{\prime\prime},}\otimes \mathsf{L}_{000h}^{efg0}=\mathsf{M^{ST}}_{abc,abc}^{F_{abc}\left[ \left( a^{\prime\prime}+e \right) \left( b^{\prime\prime}+f \right) \left( c^{\prime\prime}+g \right) \right] ,}
	\\
	&\mathsf{M^{ST}}_{abc,\tilde{a}\tilde{b}\tilde{c}}\otimes \mathsf{L}_{000h}^{efg0}=\mathsf{M^{ST}}_{abc,\tilde{a}\tilde{b}\tilde{c}}
\end{align}
where any addition is addition mod $2$. $F_{abc}$ is a function depends on the value of $\left( abc\right)$. For
\begin{align}
	\left\{ \left( abc \right) \right\} =\left\{ \left( 100 \right), \left( 010 \right), \left( 001 \right), \left( 110 \right), \left( 101 \right), \left( 011 \right), \left( 111 \right) \right\} 
\end{align}
the corresponding $F_{abc}\left[ \left( a_1+b_1 \right) \left( a_2+b_2 \right) \left( a_3+b_3 \right) \right]$ is given by
\begin{align}
	&F_{abc}\left[ \left( a_1+b_1 \right) \left( a_2+b_2 \right) \left( a_3+b_3 \right) \right] \nonumber
	\\
	=&\left\{ \left(\left( a_1+b_1 \right)00\right),\left(0\left( a_2+b_2 \right)0\right),\left(00\left( a_3+b_3 \right)\right),\left(\left( a_1+b_1 \right)00\right),\left(\left( a_1+b_1 \right)00\right),\right.\nonumber\\
	&\left.\left(0\left( a_2+b_2 \right)0\right),\left(\left( a_1+b_1 \right)00\right) \right\}
\end{align}
Now all fusion rules concerning Abelian excitations are obtained. 

\textbf{\textit{Fusion rules for two membranes.}}
Fusion rules for two membranes are much more complicated and it is hard to find a general formula. We have to calculate these fusion rules carefully and finally, we derive a fusion table for two $S^2$ membranes shown in table~\ref{table6}. 
\begin{table}
	\caption{\label{table6}Fusion rules for $\mathsf{M^S}_{abc}^{a^{\prime}b^{\prime}c^{\prime},}\otimes \mathsf{M^S}_{efg}^{e^{\prime}f^{\prime}g^{\prime},}$. This table is also symmetric, so we only have to write down the upper triangular part of this table. Repeated indices imply summation. $\oplus _{\left( abc \right) d}$ means that $a,b,c$ are not independent to each other, we constrain $\left\{\left( abc \right)\right\}=\left\{\left( 000 \right),\left( 110 \right),\left( 101 \right),\left( 011 \right)\right\}$. To be more specifically, $32\cdot\oplus _{\left( abc \right) d}\mathsf{L}_{000d}^{\left(a+a^{\prime}+\tilde{a}\right)bc0}=32\cdot \mathsf{P}_{\left(a^{\prime}+\tilde{a}\right)00}\oplus 32\cdot \mathsf{L}_{0001}^{\left(a^{\prime}+\tilde{a}\right)000}\oplus 32\cdot \mathsf{L}_{0001}^{\left(1+a^{\prime}+\tilde{a}\right)100}\oplus 32\cdot \mathsf{L}_{0001}^{\left(1+a^{\prime}+\tilde{a}\right)010}\oplus 32\cdot \mathsf{L}_{0001}^{\left(a^{\prime}+\tilde{a}\right)110}\oplus 32\cdot \mathsf{P}_{\left(1+a^{\prime}+\tilde{a}\right)10}\oplus 32\cdot \mathsf{P}_{\left(1+a^{\prime}+\tilde{a}\right)01}\oplus 32\cdot \mathsf{P}_{\left(a^{\prime}+\tilde{a}\right)11} $.}
	\centering
	\begin{tabular*}{\textwidth}{@{\extracolsep{\fill}}|c|c|c|c|c|c|c|c|}
		\hline
		& $\mathsf{M^S}_{100}^{a^{\prime}00,}$ & $\mathsf{M^S}_{010}^{0b^{\prime}0,}$ & $\mathsf{M^S}_{001}^{00c^{\prime},}$ & $\mathsf{M^S}_{110}^{a^{\prime}00,}$ & $\mathsf{M^S}_{101}^{a^{\prime}00,}$ & $\mathsf{M^S}_{011}^{0b^{\prime}0,}$ & $\mathsf{M^S}_{111}^{a^{\prime}00,}$ \\
		\hline
		$\mathsf{M^S}_{100}^{\tilde{a}00,}$ & \makecell{$2\cdot\oplus _{bcd}$\\$\mathsf{L}_{000d}^{\left(a^{\prime}+\tilde{a}\right)bc0}$} & \makecell{$\mathsf{M^S}_{110}^{100,}$\\$\oplus \mathsf{M^S}_{110}$} & \makecell{$\mathsf{M^S}_{101}^{100,}$\\$\oplus \mathsf{M^S}_{101}$} & \makecell{$4\cdot \mathsf{M^S}_{010}^{010,}$\\$\oplus4\cdot \mathsf{M^S}_{010}$} & \makecell{$4\cdot \mathsf{M^S}_{001}^{001,}$\\$\oplus4\cdot \mathsf{M^S}_{001}$} & \makecell{$\mathsf{M^S}_{111}^{100,}$\\$\oplus \mathsf{M^S}_{111}$} & \makecell{$4\cdot \mathsf{M^S}_{011}^{010,}$\\$\oplus4\cdot \mathsf{M^S}_{011}$} \\
		\hline
		$\mathsf{M^S}_{010}^{0\tilde{b}0,}$ &  & \makecell{$2\cdot\oplus _{acd}$\\  $\mathsf{L}_{000d}^{a\left(b^{\prime}+\tilde{b}\right)c0}$} & \makecell{$\mathsf{M^S}_{011}^{010,}$\\$\oplus \mathsf{M^S}_{011}$} & \makecell{$4\cdot \mathsf{M^S}_{100}^{100,}$\\$\oplus4\cdot \mathsf{M^S}_{100}$} & \makecell{$\mathsf{M^S}_{111}^{100,}$\\$\oplus \mathsf{M^S}_{111}$} & \makecell{$4\cdot \mathsf{M^S}_{001}^{001,}$\\$\oplus4\cdot \mathsf{M^S}_{001}$} & \makecell{$4\cdot \mathsf{M^S}_{101}^{100,}$\\$\oplus4\cdot \mathsf{M^S}_{101}$} \\
		\hline
		$\mathsf{M^S}_{001}^{00\tilde{c},}$ &  &  & \makecell{$2\cdot\oplus _{abd}$\\$\mathsf{L}_{000d}^{ab\left(c^{\prime}+\tilde{c}\right)0}$} & \makecell{$\mathsf{M^S}_{111}^{100,}$\\$\oplus \mathsf{M^S}_{111}$} & \makecell{$4\cdot \mathsf{M^S}_{100}^{100,}$\\$\oplus4\cdot \mathsf{M^S}_{100}$} & \makecell{$4\cdot \mathsf{M^S}_{010}^{010,}$\\$\oplus4\cdot \mathsf{M^S}_{010}$} & \makecell{$4\cdot \mathsf{M^S}_{110}^{100,}$\\$\oplus4\cdot \mathsf{M^S}_{110}$} \\
		\hline
		$\mathsf{M^S}_{110}^{\tilde{a}00,}$ &  &  &  & \makecell{$8\cdot\oplus _{acd}$\\$\mathsf{L}_{000d}^{\left(a+a^{\prime}+\tilde{a}\right)ac0}$} & \makecell{$4\cdot \mathsf{M^S}_{011}^{010,}$\\$\oplus4\cdot \mathsf{M^S}_{011}$} & \makecell{$4\cdot \mathsf{M^S}_{101}^{100,}$\\$\oplus4\cdot \mathsf{M^S}_{101}$} & \makecell{$16\cdot \mathsf{M^S}_{001}^{001,}$\\$\oplus16\cdot \mathsf{M^S}_{001}$} \\
		\hline
		$\mathsf{M^S}_{101}^{\tilde{a}00,}$ &  &  &  &  & \makecell{$8\cdot\oplus _{abd}$\\$\mathsf{L}_{000d}^{\left(a+a^{\prime}+\tilde{a}\right)ba0}$} & \makecell{$4\cdot \mathsf{M^S}_{110}^{100,}$\\$\oplus4\cdot \mathsf{M^S}_{110}$} & \makecell{$16\cdot \mathsf{M^S}_{010}^{010,}$\\$\oplus16\cdot \mathsf{M^S}_{010}$} \\
		\hline
		$\mathsf{M^S}_{011}^{0\tilde{b}0,}$ &  &  &  &  &  & \makecell{$8\cdot\oplus _{abd}$\\$\mathsf{L}_{000d}^{a\left(b+b^{\prime}+\tilde{b}\right)b0}$} & \makecell{$16\cdot \mathsf{M^S}_{100}^{100,}$\\$\oplus16\cdot \mathsf{M^S}_{100}$} \\
		\hline
		$\mathsf{M^S}_{111}^{\tilde{a}00,}$ &  &  &  &  &  &  & \makecell{$32\cdot\oplus _{\left( abc \right) d}$\\$\mathsf{L}_{000d}^{\left(a+a^{\prime}+\tilde{a}\right)bc0}$} \\
		\hline
	\end{tabular*}
\end{table}

For two $T^2$ membranes, we have a similar table. First, we replace all $\mathsf{M^S}$ with $\mathsf{M^T}$ in table~\ref{table6}. Notice that fusion coefficients larger than $1$ in table~\ref{table6} can be $2,4,8,16,32$, we replace them with $8,16,128,256,2048$ respectively, then we have a fusion table for $\mathsf{M^T}_{\tilde{a}\tilde{b}\tilde{c}}^{\tilde{a}^{\prime}\tilde{b}^{\prime}\tilde{c}^{\prime},}\otimes \mathsf{M^T}_{\tilde{e}\tilde{f}\tilde{g}}^{\tilde{e}^{\prime}\tilde{f}^{\prime}\tilde{g}^{\prime},}$. We do this replacement because the coefficients of Wilson operators for $T^2$ membranes are larger than $S^2$ membranes so that all shrinking coefficients are integers. For example, we have 
\begin{align}
	\mathsf{M^S}_{110}^{a^{\prime}00,}\otimes \mathsf{M^S}_{110}^{\tilde{a}00,}=8\cdot \oplus _{acd}\mathsf{L}_{000d}^{\left( a+a^{\prime}+\tilde{a} \right) ac0},\quad \mathsf{M^S}_{111}^{a^{\prime}00,}\otimes \mathsf{M^S}_{011}^{0\tilde{b}0,}=16\cdot \mathsf{M^S}_{100}^{100,}\oplus 16\cdot \mathsf{M^S}_{100} \nonumber
\end{align}
in table~\ref{table6}. Do the replacement we obtain 
\begin{align}
	\mathsf{M^T}_{101}^{a^{\prime}00,}\otimes \mathsf{M^T}_{110}^{\tilde{a}00,}=16\cdot \mathsf{M^T}_{011}^{010,}\oplus 16\cdot \mathsf{M^T}_{011},\quad \mathsf{M^T}_{111}^{a^{\prime}00,}\otimes \mathsf{M^T}_{011}^{0\tilde{b}0,}=256\cdot \mathsf{M^T}_{100}^{100,}\oplus 256\cdot \mathsf{M^T}_{100} \nonumber
\end{align}

As for fusing an $S^2$ membrane and a $T^2$ membrane, we have:
\begin{align}
	\mathsf{M^T}_{\tilde{a}\tilde{b}\tilde{c}}^{\tilde{a}^{\prime}\tilde{b}^{\prime}\tilde{c}^{\prime},}\otimes \mathsf{M^S}_{efg}^{e^{\prime}f^{\prime}g^{\prime},}=\mathsf{M^{ST}}_{efg,\tilde{a}\tilde{b}\tilde{c}}^{\left( e^{\prime}+\tilde{a}^{\prime} \right) \left( f^{\prime}+\tilde{b}^{\prime} \right) \left( g^{\prime}+\tilde{c}^{\prime} \right)}=\begin{cases}
		\mathsf{M^{ST}}_{efg,efg}^{\tilde{e}^{\prime \prime}\tilde{f}^{\prime \prime}\tilde{g}^{\prime \prime}}&\quad\text{for}\quad efg=\tilde{a}\tilde{b}\tilde{c}\\
		\mathsf{M^{ST}}_{efg,\tilde{a}\tilde{b}\tilde{c}}&\quad\text{for}\quad efg\ne \tilde{a}\tilde{b}\tilde{c}
	\end{cases}
\end{align}
Now we obtain three more elements in table~\ref{table5}.

\textbf{\textit{Fusion rules for $\mathsf{M^{ST}}$-excitations.}}
At the beginning of section~\ref{s3}, we illustrate that fusion rules in our theory should satisfy commutative and associative law, i.e., fusion rules satisfy eq.~(\ref{eq_fusion_algebra_rule}). We will use this powerful condition to calculate fusion rules for $\mathsf{M^{ST}}$-excitations. 

Notice that $\mathsf{M^{ST}}_{abc,abc}^{a^{\prime\prime}b^{\prime\prime}c^{\prime\prime},}$ and $\mathsf{M^{ST}}_{abc,\tilde{a}\tilde{b}\tilde{c}}$ are equivalent to:
\begin{align}
	\mathsf{M^{ST}}_{abc,abc}^{a^{\prime\prime}b^{\prime\prime}c^{\prime\prime},}=&\mathsf{M^S}_{abc}^{a^{\prime\prime}b^{\prime\prime}c^{\prime\prime},}\otimes \mathsf{M^T}_{abc}
	\\
	\mathsf{M^{ST}}_{abc,\tilde{a}\tilde{b}\tilde{c}}=&\mathsf{M^S}_{abc}\otimes \mathsf{M^T}_{\tilde{a}\tilde{b}\tilde{c}} \qquad \left(abc\ne\tilde{a}\tilde{b}\tilde{c}\right)
\end{align}
Thus we have 
\begin{align}
	\mathsf{M^{ST}}_{abc,abc}^{a^{\prime\prime}b^{\prime\prime}c^{\prime\prime},}\otimes \mathsf{M^S}_{efg}^{e^{\prime}f^{\prime}g^{\prime},}=\mathsf{M^S}_{abc}^{a^{\prime\prime}b^{\prime\prime}c^{\prime\prime},}\otimes \mathsf{M^T}_{abc}\otimes \mathsf{M^S}_{efg}^{e^{\prime}f^{\prime}g^{\prime},}=\mathsf{M^S}_{abc}^{a^{\prime\prime}b^{\prime\prime}c^{\prime\prime},}\otimes \mathsf{M^S}_{efg}^{e^{\prime}f^{\prime}g^{\prime},}\otimes \mathsf{M^T}_{abc}
\end{align}
Using table~\ref{table6} we can always write:
\begin{align}
	\mathsf{M^S}_{abc}^{a^{\prime\prime}b^{\prime\prime}c^{\prime\prime},}\otimes \mathsf{M^S}_{efg}^{e^{\prime}f^{\prime}g^{\prime},}=\oplus_{l}\mathsf{E}_{l}
\end{align}
where $\mathsf{E}_{l}$ represent excitations appear in the fusion results in table~\ref{table6}, coefficient is absorbed in $\mathsf{E}_{l}$. Then we obtain 
\begin{align}
	\mathsf{M^{ST}}_{abc,abc}^{a^{\prime\prime}b^{\prime\prime}c^{\prime\prime},}\otimes \mathsf{M^S}_{efg}^{e^{\prime}f^{\prime}g^{\prime},}=\oplus_{l}\left(\mathsf{E}_{l}\otimes \mathsf{M^T}_{abc}\right)
\end{align}
Since we have already known the fusion rules for $\mathsf{M^T}_{\tilde{a}\tilde{b}\tilde{c}}^{\tilde{a}^{\prime}\tilde{b}^{\prime}\tilde{c}^{\prime},}\otimes \mathsf{L}_{000h}^{efg0}$ and $\mathsf{M^T}_{\tilde{a}\tilde{b}\tilde{c}}^{\tilde{a}^{\prime}\tilde{b}^{\prime}\tilde{c}^{\prime},}\otimes \mathsf{M^S}_{efg}^{e^{\prime}f^{\prime}g^{\prime},}$, also $\mathsf{E}_{l}$ eventually takes the form of $\mathsf{L}_{000d}^{abc0}$ or $\mathsf{M^S}_{abc}^{a^{\prime}b^{\prime}c^{\prime},}$, we can thus obtain the fusion results of $\mathsf{E}_{l}\otimes \mathsf{M^T}_{abc}$, denoted as $\oplus_{l^{\prime}}\mathsf{\epsilon}_{ll^{\prime}}$, finally we obtain 
\begin{align}
	\mathsf{M^{ST}}_{abc,abc}^{a^{\prime\prime}b^{\prime\prime}c^{\prime\prime},}\otimes \mathsf{M^S}_{efg}^{e^{\prime}f^{\prime}g^{\prime},}=\oplus_{l}\left(\oplus_{l^{\prime}}\mathsf{\epsilon}_{ll^{\prime}}\right)=\oplus_{ll^{\prime}}\mathsf{\epsilon}_{ll^{\prime}}
\end{align}
For other fusion rules concerning $\mathsf{M^{ST}}$-excitations, such as $\mathsf{M^{ST}}_{abc,\tilde{a}\tilde{b}\tilde{c}}\otimes \mathsf{M^S}_{efg}^{e^{\prime}f^{\prime}g^{\prime},}$, we have similar discussions. So in principle, we can obtain the rest of $7$ elements in table~\ref{table5}.

For example, consider $\mathsf{M^{ST}}_{110,110}^{100,}\otimes \mathsf{M^S}_{011}^{010,}$, following the method in this section we have 
\begin{align}
	\mathsf{M^{ST}}_{110,110}^{100,}\otimes \mathsf{M^S}_{011}^{010,}=&\mathsf{M^S}_{110}^{100,}\otimes \mathsf{M^T}_{110}\otimes \mathsf{M^S}_{011}^{010,}=\mathsf{M^S}_{110}^{100,}\otimes \mathsf{M^S}_{011}^{010,}\otimes \mathsf{M^T}_{110} \nonumber
	\\
	=&4\cdot \left( \mathsf{M^S}_{101}^{100,}\oplus \mathsf{M^S}_{101} \right) \otimes \mathsf{M^T}_{110}=4\cdot \left( \mathsf{M^{ST}}_{101,110}\oplus \mathsf{M^{ST}}_{101,110} \right) \nonumber
	\\
	=&8\cdot \mathsf{M^{ST}}_{101,110}
\end{align}

Now for any two given excitations, we can obtain their fusion rules. 
\subsection{Shrinking rules for the $AAAB$ twisted term}
We have already defined shrinking operator $\mathcal{S}$ for different excitations in section~\ref{s4}, we continue to use this definition. 

\textbf{\textit{Shrinking an $S^2$ membrane.}}
Shrinking rules for $S^2$ membranes can be written as 
\begin{align}
	\mathcal{S} \left( \mathsf{M^S}_{abc}^{a^{\prime}b^{\prime}c^{\prime},} \right) =&N_{abc}\left\{ a^{\prime}\mathsf{P}_{a^{\prime}00}\otimes b^{\prime}\mathsf{P}_{0b^{\prime}0}\otimes c^{\prime}\mathsf{P}_{00c^{\prime}} \right. \nonumber
	\\
	&\otimes \left[ \frac{ab}{2}\left( \mathsf{1}\oplus \mathsf{P}_{110} \right) \right] \otimes \left[ \frac{ac}{2}\left( \mathsf{1}\oplus \mathsf{P}_{101} \right) \right] \otimes \left[ \frac{bc}{2}\left( \mathsf{1}\oplus \mathsf{P}_{011} \right) \right] \nonumber
	\\
	&\left. \otimes \left[ \frac{\left( 1-a \right)}{2}\left( \mathsf{1}\oplus \mathsf{P}_{100} \right) \right] \otimes \left[ \frac{\left( 1-b \right)}{2}\left( \mathsf{1}\oplus \mathsf{P}_{010} \right) \right] \otimes \left[ \frac{\left( 1-c \right)}{2}\left( \mathsf{1}\oplus \mathsf{P}_{001} \right) \right] \right\} 
\end{align}
In the above equation, if the coefficient of the excitation is zero, we just omit it. For example terms like $\otimes \left[ \frac{0}{2}\left( 1\oplus \mathsf{P}_{abc} \right) \right] $ are omitted. If we consider shrinking an $\mathsf{M^S}_{101}^{100,}$, we have 
\begin{align}
	\mathcal{S} \left( {\mathsf{M}^{\mathsf{S}}}_{101}^{100,} \right) =&8\left\{ \mathsf{P}_{100}\otimes \left[ \frac{1}{2}\left( 1\oplus \mathsf{P}_{101} \right) \right] \otimes \left[ \frac{1}{2}\left( 1\oplus \mathsf{P}_{010} \right) \right] \right\} \nonumber
	\\
	=&2\cdot \left( \mathsf{P}_{100}\oplus \mathsf{P}_{001}\oplus \mathsf{P}_{110}\oplus \mathsf{P}_{011} \right) 
\end{align}

Through our calculation, we find that shrinking results of all $S^2$ membranes contain $4$ different kinds of particles. This is different from the $BBA$ twisted term where all pure membranes are shrunk to vacua.

\textbf{\textit{Shrinking a $T^2$ membrane.}}
Shrinking rules for $T^2$ membranes can be written as 
\begin{align}
	\mathcal{S} \left( \mathsf{M^T}_{abc}^{a^{\prime}b^{\prime}c^{\prime},} \right) =&\frac{\tilde{N}_{abc}}{2}\left\{ a^{\prime}\mathsf{P}_{a^{\prime}00}\otimes b^{\prime}\mathsf{P}_{0b^{\prime}0}\otimes c^{\prime}\mathsf{P}_{00c^{\prime}}\otimes \left( \mathsf{1}\oplus \mathsf{L}_{0001} \right) \right. \nonumber
	\\
	&\otimes \left[ \frac{ab}{2}\left( \mathsf{1}\oplus \mathsf{P}_{110} \right) \right] \otimes \left[ \frac{ac}{2}\left( \mathsf{1}\oplus \mathsf{P}_{101} \right) \right] \otimes \left[ \frac{bc}{2}\left( \mathsf{1}\oplus \mathsf{P}_{011} \right) \right] \nonumber
	\\
	&\otimes \left[ \frac{\left( 1-a \right)}{2}\left( \mathsf{1}\oplus \mathsf{P}_{100} \right) \right] \otimes \left[ \frac{\left( 1-b \right)}{2}\left( \mathsf{1}\oplus \mathsf{P}_{010} \right) \right] \nonumber
	\\
	&\left.\otimes \left[ \frac{\left( 1-c \right)}{2}\left( \mathsf{1}\oplus \mathsf{P}_{001} \right) \right] \right\} \nonumber
	\\
	\mathcal{S}^{2} \left( \mathsf{M^T}_{abc}^{a^{\prime}b^{\prime}c^{\prime},} \right) =&\frac{\tilde{N}_{abc}}{N_{abc}}\mathcal{S} \left( \mathsf{M^S}_{abc}^{a^{\prime}b^{\prime}c^{\prime},} \right)
	\label{eq30}
\end{align}
If a $T^2$ and an $S^2$ membrane carry the same gauge charges and decorations, they can eventually be shrunk to the same superposition of particles up to a global coefficient $\frac{\tilde{N}_{abc}}{N_{abc}}$. When they only carry one unit of gauge charge, this global coefficient is $2^1=2$, carry two units is $2^2=4$, and carry three units is $2^3=8$. Every unit of gauge charge contributes an extra $2$ to the normalization factor $\tilde{N}_{abc}$ in the front of Wilson operators for $T^2$ membranes.

\textbf{\textit{Shrinking an $\mathsf{M^{ST}}$-excitation.}}
When $\left(abc\right)=\left(\tilde{a}\tilde{b}\tilde{c}\right)$, shrinking rules for $\mathsf{M^{ST}}$-excitations can be written as 
\begin{align}
	\mathcal{S} \left( \mathsf{M^{ST}}_{abc,abc}^{a^{\prime}b^{\prime}c^{\prime},} \right) =&N_{abc}\mathcal{S} \left( \mathsf{M^T}_{abc}^{a^{\prime}b^{\prime}c^{\prime},} \right) \nonumber
	\\
	\mathcal{S} ^2\left( \mathsf{M^{ST}}_{abc,abc}^{a^{\prime}b^{\prime}c^{\prime},} \right) =&\tilde{N}_{abc}\mathcal{S} \left( \mathsf{M^S}_{abc}^{a^{\prime}b^{\prime}c^{\prime},} \right) 
\end{align}
Compare to the shrinking rules eq.~(\ref{eq30}) for $T^2$ membranes, we see that the shrinking rules for $\mathsf{M^{ST}}$-excitations with $\left(abc\right)=\left(\tilde{a}\tilde{b}\tilde{c}\right)$ are actually $N_{abc}$ times eq.~(\ref{eq30}). 

When $\left(abc\right)\ne\left(\tilde{a}\tilde{b}\tilde{c}\right)$, we have
\begin{align}
	\mathcal{S} \left( \mathsf{M^{ST}}_{abc,\tilde{a}\tilde{b}\tilde{c}} \right) =&\frac{1}{16}N_{abc}\tilde{N}_{\tilde{a}\tilde{b}\tilde{c}}\left[ \left( \mathsf{1}\oplus \mathsf{L}_{0001} \right) \otimes \left( \mathsf{1}\oplus \mathsf{P}_{100} \right) \otimes \left( \mathsf{1}\oplus \mathsf{P}_{010} \right) \otimes \left( \mathsf{1}\oplus \mathsf{P}_{001} \right) \right] \nonumber
	\\
	=&\frac{1}{16}N_{abc}\tilde{N}_{\tilde{a}\tilde{b}\tilde{c}}\left( \oplus _{efgh}\mathsf{L}_{000h}^{efg0} \right) \nonumber
	\\
	=&\frac{1}{16}N_{abc}\tilde{N}_{\tilde{a}\tilde{b}\tilde{c}}\left( \mathsf{Ab} \right)  
\end{align}
where we use $\mathsf{Ab}$ to denote $\left( \oplus _{efgh}\mathsf{L}_{000h}^{efg0} \right)$, which is the superposition of all Abelian excitations. Keep shrinking we obtain 
\begin{align}
	\mathcal{S} ^2\left( \mathsf{M^{ST}}_{abc,\tilde{a}\tilde{b}\tilde{c}} \right) =&\frac{1}{8}N_{abc}\tilde{N}_{\tilde{a}\tilde{b}\tilde{c}}\left( \oplus _{efg}\mathsf{P}_{efg} \right)
\end{align}
which is the superposition of all particles.

Now we have obtained all shrinking rules for $AAAB$ twisted term. One can verify that shrinking rules respect fusion rules, i.e., satisfy eq.~(\ref{eq19}):
\begin{align}
	\mathcal{S} \left( a\otimes b \right) =\mathcal{S} \left( a \right) \otimes \mathcal{S} \left( b \right), \qquad \mathcal{S}^2 \left( a\otimes b \right) =\mathcal{S}^2 \left( a \right) \otimes \mathcal{S}^2 \left( b \right)
\end{align}

\section{General method to determine the type of fusion and shrinking rules\label{ap2}}
Given a specific twisted term, we do not have to derive all fusion and shrinking rules explicitly to find out whether the fusion and shrinking rules are Abelian or non-Abelian. We just need to derive the gauge transformations and a few Wilson operators, then we can determine the type of fusion and shrinking rules.

Generally, Wilson operators for topological excitations in $BF$ theory can be written in these two forms 
\begin{align}
	\mathsf{E}_a=&N_a\prod_i{\exp \left( i\int_{X _i}{F^i+f^i} \right)} \label{eq1ap2}
	\\
	\mathsf{E}_b=&N_b\left[ \prod_j{\exp \left( i\int_{X _j}{F^j+f^j} \right)} \right] \left[ \prod_k{\delta \left( \int_{X _k}{F^k} \right)} \right]  \label{eq2ap2}
\end{align}
where $N_a$ and $N_b$ are normalization factors, $X _i$ and $X _j$ are spacetime trajectories of excitations, $F^i$ and $F^j$ are gauge fields, $f^i$ and $f^j$ are nontrivial terms. If $F^i$ transforms without nontrivial shifts, $f^i$ will vanish. 

We conclude that all excitations in the form of eq.~(\ref{eq1ap2}) are Abelian excitation because we can always find that fusing two $\mathsf{E}_a$-form excitations gives another excitation that still has $\mathsf{E}_a$-form, fusing a $\mathsf{E}_a$-form excitations and a $\mathsf{E}_b$-form excitations gives another excitation that has $\mathsf{E}_b$-form. Since all fusion channels for $\mathsf{E}_a$ are single, this is an Abelian excitation.

For all excitations in the form of eq.~(\ref{eq2ap2}), we conclude that they are non-Abelian. Notice that $\delta \left( \int_{X _k}{F^k} \right) $ will be expanded as $\frac{1}{2}\left[ 1+\exp \left( \frac{i2\pi m_k}{2} \right) \right] $ thus 
\begin{align}
	\mathsf{E}_b\otimes \mathsf{E}_b=N\left[ \otimes _k\left( \mathsf{1}\oplus \mathsf{E}_k \right) \right] 
\end{align}
This is sufficient to show that $\mathsf{E}_b$ in the form of eq.~(\ref{eq2ap2}) is non-Abelian.

Now we consider shrinking excitations. When we apply the shrinking operator to the excitation, $\prod_j{\exp \left( i\int_{X _j}{F^j+f^j} \right)}$ will vanish, only some delta functionals may survive and these delta functionals can be expanded, causing other nontrivial excitations. Thus, only excitations that carry delta functionals may have non-Abelian shrinking rules. If we want to have hierarchical shrinking rules, which can only happen in $T^2$ membranes, there must be delta functionals in the form of $\delta \left( \int_{\sigma}{B^i} \right) $. When we shrink the $T^2$ membranes, $\delta \left( \int_{\sigma}{B^i} \right) $ will not vanish and we can expand it as $\frac{1}{2}\left[ 1+\exp \left( \frac{i2\pi m_i}{2} \right) \right] $, causing $\left( \mathsf{1}\oplus \mathsf{L}_i \right)$ to appear in the shrinking results.

From the above discussion, we can see that delta functionals determine whether there exist non-Abelian fusion rules and non-Abelian shrinking rules. We can list the basic Wilson operators whose corresponding gauge field transformations with nontrivial shifts to see if there are delta functionals.

For $AAC$ twisted term, TQFT action is 
\begin{align}
	S=\int{\sum_{i=1}^3{\frac{N_i}{2\pi}C^idA^i+qA^1A^2C^3}}
\end{align} 
Gauge transformations are 
\begin{align}
	&A^{1,2}\rightarrow A^{1,2}+d\chi ^{1,2},\qquad C^3\rightarrow C^3+dT^3 \nonumber
	\\
	&C^1\rightarrow C^1+dT^1-\frac{2\pi q}{N_1}\left( \chi ^2C^3-A^2T^3+\chi ^2dT^3 \right) \nonumber
	\\
	&C^2\rightarrow C^2+dT^2+\frac{2\pi q}{N_2}\left( \chi ^1C^3-A^1T^3+\chi ^1dT^3 \right) \nonumber
	\\
	&A^3\rightarrow A^3+d\chi ^3-\frac{2\pi q}{N_3}\left[ \left( \chi ^1A^2+\frac{1}{2}\chi ^1d\chi ^2 \right) - \left( \chi ^2A^1+\frac{1}{2}\chi ^2d\chi ^1 \right)\right]  
\end{align}
Basic nontrivial Wilson operators are
\begin{align}
	&W_1=N_1\exp \left[ i\int_{\omega}{C^1+\frac{1}{2}\frac{2\pi q}{N_1}\left( d^{-1}A^2C^3-d^{-1}C^3A^2 \right)} \right] \delta \left( \int_{\gamma}{A^2} \right) \delta \left( \int_{\omega}{C^3} \right) \nonumber
	\\
	&W_2=N_2\exp \left[ i\int_{\omega}{C^2+\frac{1}{2}\frac{2\pi q}{N_2}\left( d^{-1}C^3A^1-d^{-1}A^1C^3 \right)} \right] \delta \left( \int_{\gamma}{A^1} \right) \delta \left( \int_{\omega}{C^3} \right) \nonumber
	\\
	&W_3=N_3\exp \left[ i\int_{\omega}{C^3+\frac{1}{2}\frac{2\pi q}{N_3}\left( d^{-1}A^1A^2-d^{-1}A^2A^1 \right)} \right] \delta \left( \int_{\gamma}{A^1} \right) \delta \left( \int_{\gamma}{A^2} \right) 
\end{align}
Other Wilson operators that have nontrivial terms and delta functionals can always be constructed from the above three basic nontrivial Wilson operators. We can see that there exist $\delta \left( \int_{\gamma}{A^1}\right)$, $\delta \left( \int_{\gamma}{A^2}\right)$ and $\delta \left( \int_{\omega}{C^3} \right)$, but $\delta \left( \int_{\sigma}{B^i} \right)$ is absents. Thus we can conclude that $AAC$ twisted terms have non-Abelian fusion rules and non-Abelian shrinking rules, but do not have hierarchical shrinking rules.

Similarly, for $AAAAA$ twisted term, TQFT action is
\begin{align}
	S=\int{\sum_{i=1}^5{\frac{N_i}{2\pi}C^idA^i+qA^1A^2A^3A^4A^5}}
\end{align}
Gauge transformations are 
\begin{align}
	A^i\rightarrow& A^i+d\chi ^i \nonumber
	\\
	C^i\rightarrow& C^i+dT^i-\frac{2\pi q}{N_i}\sum_{k<l<m}{\epsilon ^{ijklm}}\chi ^jA^kA^lA^m-\frac{2\pi q}{N_i}\sum_{j<k,l<m}{\epsilon ^{ijklm}}\chi ^jd\chi ^kA^lA^m \nonumber
	\\
	&-\frac{2\pi q}{N_i}\sum_{j<k<l}{\sum_{m=1}^5{\epsilon ^{ijklm}}}\chi ^jd\chi ^kd\chi ^lA^m-\frac{2\pi q}{N_i}\sum_{j<k<l<m}{\epsilon ^{ijklm}}\chi ^jd\chi ^kd\chi ^ld\chi ^m 
\end{align}
Basic nontrivial Wilson operators are
\begin{align}
	W_i=N_i\exp \left[ i\int_{\omega}{C^1+\frac{1}{4}\frac{2\pi q}{N_i}\sum_{k<l<m}{\sum_{j=1}^5{\epsilon ^{ijklm}}}d^{-1}A^jA^kA^lA^m} \right] \prod_{n=1}^5{\delta \left( \int_{\gamma}{A^n} \right)}
\end{align}
Thus $AAAAA$ twisted terms have non-Abelian fusion rules and non-Abelian shrinking rules, but do not have hierarchical shrinking rules.

For $AAAdA$ twisted term, TQFT action is
\begin{align}
	S=\int{\sum_{i=1}^3{\frac{N_i}{2\pi}}C^idA^i}+qA^1A^2A^3dA^1
\end{align}
Gauge transformations are
\begin{align}
	&A^1\rightarrow A^1+d\chi ^1,C^1\rightarrow C^1+dT^1-\frac{2\pi q}{N_1}\left( A^1d\chi ^2A^3+A^1A^2d\chi ^3+A^1d\chi ^2d\chi ^3 \right) ,\nonumber
	\\
	&A^2\rightarrow A^2+d\chi ^2,C^2\rightarrow C^2+dT^2+\frac{2\pi q}{N_2}\left( d\chi ^1A^3A^1+d\chi ^1d\chi ^3A^1 \right) ,\nonumber
	\\
	&A^3\rightarrow A^3+d\chi ^3,C^3\rightarrow C^3+dT^3-\frac{2\pi q}{N_3}\left( d\chi ^1A^2A^1+d\chi ^1d\chi ^2A^1 \right) 
\end{align}
Basic nontrivial Wilson operators are
\begin{align}
	W_1=&N_1\exp \left[ i\int_{\omega}{C^1+\frac{1}{2}\frac{2\pi q}{N_1}\left( d^{-1}A^2A^3dA^1-d^{-1}A^3A^2dA^1 \right)} \right] \delta \left( \int_{\gamma}{A^2} \right) \delta \left( \int_{\gamma}{A^3} \right) \nonumber
	\\
	W_2=&N_2\exp \left[ i\int_{\omega}{C^2+\frac{1}{2}\frac{2\pi q}{N_2}\left( d^{-1}A^1A^3dA^1-d^{-1}A^3A^1dA^1 \right)} \right] \delta \left( \int_{\gamma}{A^1} \right) \delta \left( \int_{\gamma}{A^3} \right) \nonumber
	\\
	W_3=&N_3\exp \left[ i\int_{\omega}{C^3+\frac{1}{2}\frac{2\pi q}{N_3}\left( d^{-1}A^1A^2dA^1-d^{-1}A^2A^1dA^1 \right)} \right] \delta \left( \int_{\gamma}{A^1} \right) \delta \left( \int_{\gamma}{A^2} \right) 
\end{align}
Thus $AAAdA$ twisted terms have non-Abelian fusion rules and non-Abelian shrinking rules, but do not have hierarchical shrinking rules.

For $AdAdA$ twisted term, TQFT action is
\begin{align}
	S=\int{\sum_{i=1}^3{\frac{N_i}{2\pi}C^idA^i+qA^1dA^2dA^3}}
\end{align}
Gauge transformations are 
\begin{align}
	A^i\rightarrow A^i+d\chi ^i,\qquad C^i\rightarrow C^i+dT^i
\end{align}
All fields transform without nontrivial shifts and thus all Wilson operators do not have nontrivial terms and delta functionals. Thus $AdAdA$ twisted term only has Abelian fusion rules and Abelian shrinking rules.

For $AdAB$ twisted term, TQFT action is
\begin{align}
	S=\int{\frac{N_1}{2\pi}C^1dA^1+\frac{N_2}{2\pi}C^2dA^2+\frac{N_2}{2\pi}\tilde{B}^3dB^3+qA^1dA^2B^3}
\end{align}
Gauge transformations are 
\begin{align}
	&A^1\rightarrow A^1+d\chi ^1,\qquad C^1\rightarrow C^1+dT^1-\frac{2\pi q}{N_1}dV^3A^2 \nonumber
	\\
	&A^2\rightarrow A^2+d\chi ^2,\qquad C^2\rightarrow C^2+dT^2 \nonumber
	\\
	&B^3\rightarrow B^3+dV^3,\qquad \tilde{B}^3\rightarrow \tilde{B}^3+d\tilde{V}^3-\frac{2\pi q}{N_3}d\chi ^1A^2
\end{align}
Basic nontrivial Wilson operators are
\begin{align}
	W_1=&N_1\exp \left( i\int_{\omega}{C^1+\frac{2\pi q}{N_i}}\int_{\Xi}{B^3dA^2} \right) \nonumber
	\\
	W_3=&N_3\exp \left( i\int_{\sigma}{B^3-\frac{2\pi q}{N_3}}\int_{\Omega}{A^1dA^2} \right) 
\end{align}
Although some Wilson operators have nontrivial terms, delta functionals are still absent. These nontrivial terms cannot change fusion and shrinking rules. Thus $AdAB$ twisted term only has Abelian fusion rules and Abelian shrinking rules.

For $AAdB$ twisted term, notice that 
\begin{align}
	d\left( A^1A^2B^3 \right) =A^1A^2dB^3+A^2dA^1B^3-A^1dA^2B^3
\end{align}
This means $AAdB$ twisted term can always be written as two $AdAB$ twisted terms up to a boundary term. Thus $AAdB$ twisted term only has Abelian fusion rules and Abelian shrinking rules.

% figures should be put into the text as floats.
% Use the graphics or graphicx packages (distributed with LaTeX2e)
% and the \includegraphics macro defined in those packages.
% See the LaTeX Graphics Companion by Michel Goosens, Sebastian Rahtz,
% and Frank Mittelbach for instance.
%
% Here is an example of the general form of a figure:
% Fill in the caption in the braces of the \caption{} command. Put the label
% that you will use with \ref{} command in the braces of the \label{} command.
% Use the figure* environment if the figure should span across the
% entire page. There is no need to do explicit centering.

% \begin{figure}
% \includegraphics{}%
% \caption{\label{}}
% \end{figure}

% Surround figure environment with turnpage environment for landscape
% figure
% \begin{turnpage}
% \begin{figure}
% \includegraphics{}%
% \caption{\label{}}
% \end{figure}
% \end{turnpage}

% tables should appear as floats within the text
%
% Here is an example of the general form of a table:
% Fill in the caption in the braces of the \caption{} command. Put the label
% that you will use with \ref{} command in the braces of the \label{} command.
% Insert the column specifiers (l, r, c, d, etc.) in the empty braces of the
% \begin{tabular}{} command.
% The ruledtabular enviroment adds doubled rules to table and sets a
% reasonable default table settings.
% Use the table* environment to get a full-width table in two-column
% Add \usepackage{longtable} and the longtable (or longtable*}
% environment for nicely formatted long tables. Or use the the [H]
% placement option to break a long table (with less control than 
% in longtable).
% \begin{table}%[H] add [H] placement to break table across pages
% \caption{\label{}}
% \begin{ruledtabular}
% \begin{tabular}{}
% Lines of table here ending with \\
% \end{tabular}
% \end{ruledtabular}
% \end{table}

% Surround table environment with turnpage environment for landscape
% table
% \begin{turnpage}
% \begin{table}
% \caption{\label{}}
% \begin{ruledtabular}
% \begin{tabular}{}
% \end{tabular}
% \end{ruledtabular}
% \end{table}
% \end{turnpage}

% Specify following sections are appendices. Use \appendix* if there
% only one appendix.
%\appendix
%\section{}

% If you have acknowledgments, this puts in the proper section head.
%\begin{acknowledgments}
% put your acknowledgments here.
%\end{acknowledgments}

% Create the reference section using BibTeX:
\bibliographystyle{JHEP}
%\bibliography{braid}
 
\begin{thebibliography}{10}

\bibitem{wenZootopoRMP}
X.-G.~Wen, \emph{Colloquium: Zoo of quantum-topological phases of matter},
  \href{https://doi.org/10.1103/RevModPhys.89.041004}{\emph{Rev. Mod. Phys.}
  {\bfseries 89} (2017) 041004}.

\bibitem{string1}
M.A.~Levin and X.-G.~Wen, \emph{String-net condensation: A physical mechanism
  for topological phases},
  \href{https://doi.org/10.1103/PhysRevB.71.045110}{\emph{Phys. Rev. B}
  {\bfseries 71} (2005) 045110}.

\bibitem{dwitten}
R.~Dijkgraaf and E.~Witten, \emph{Topological gauge theories and group
  cohomology}, \href{https://doi.org/10.1007/BF02096988}{\emph{Commun. Math.
  Phys.} {\bfseries 129} (1990) 393}.

\bibitem{Kitaev2003faulttolerant}
A.Y.~{Kitaev}, \emph{{Fault-tolerant quantum computation by anyons}},
  \href{https://doi.org/10.1016/S0003-4916(02)00018-0}{\emph{Annals of Physics}
  {\bfseries 303} (2003) 2}
  [\href{https://arxiv.org/abs/quant-ph/9707021}{{\ttfamily
  quant-ph/9707021}}].

\bibitem{CGW2010LRE}
X.~Chen, Z.-C.~Gu and X.-G.~Wen, \emph{Local unitary transformation, long-range
  quantum entanglement, wave function renormalization, and topological order},
  \href{https://doi.org/10.1103/PhysRevB.82.155138}{\emph{Phys. Rev. B}
  {\bfseries 82} (2010) 155138}.

\bibitem{Li_Haldane}
H.~Li and F.D.M.~Haldane, \emph{Entanglement spectrum as a generalization of
  entanglement entropy: Identification of topological order in non-abelian
  fractional quantum hall effect states},
  \href{https://doi.org/10.1103/PhysRevLett.101.010504}{\emph{Phys. Rev. Lett.}
  {\bfseries 101} (2008) 010504}.

\bibitem{KitaevTopoEE}
A.~Kitaev and J.~Preskill, \emph{Topological entanglement entropy},
  \href{https://doi.org/10.1103/PhysRevLett.96.110404}{\emph{Phys. Rev. Lett.}
  {\bfseries 96} (2006) 110404}.

\bibitem{Levin_Wen_TEE}
M.~Levin and X.-G.~Wen, \emph{Detecting topological order in a ground state
  wave function},
  \href{https://doi.org/10.1103/PhysRevLett.96.110405}{\emph{Phys. Rev. Lett.}
  {\bfseries 96} (2006) 110405}.

\bibitem{Kitaev2006}
A.~Kitaev, \emph{Anyons in an exactly solved model and beyond},
  \href{https://doi.org/10.1016/j.aop.2005.10.005}{\emph{Ann. Phys.} {\bfseries
  321} (2006) 2}.

\bibitem{Kong2020}
L.~Kong, T.~Lan, X.G.~Wen, Z.H.~Zhang and H.~Zheng, \emph{{Algebraic higher
  symmetry and categorical symmetry - a holographic and entanglement view of
  symmetry}},  may, 2020.
\newblock 10.1103/physrevresearch.2.043086.

\bibitem{Witten1989}
E.~Witten, \emph{Quantum field theory and the jones polynomial},
  \href{https://doi.org/10.1007/BF01217730}{\emph{Commun. Math. Phys.}
  {\bfseries 121} (1989) 351}.

\bibitem{Turaev2016}
V.G.~Turaev, \emph{Quantum Invariants of Knots and 3-Manifolds}, De Gruyter,
  Berlin, Boston (2016),
  \href{https://doi.org/doi:10.1515/9783110435221}{doi:10.1515/9783110435221}.

\bibitem{sarma_08_TQC}
C.~Nayak, S.H.~Simon, A.~Stern, M.~Freedman and S.~Das~Sarma, \emph{Non-abelian
  anyons and topological quantum computation},
  \href{https://doi.org/10.1103/RevModPhys.80.1083}{\emph{Rev. Mod. Phys.}
  {\bfseries 80} (2008) 1083}.

\bibitem{moore1989classical}
G.~Moore and N.~Seiberg, \emph{Classical and quantum conformal field theory},
  {\emph{Communications in Mathematical Physics} {\bfseries 123} (1989) 177}.

\bibitem{Wen_grav_prd2013}
X.-G.~Wen, \emph{Classifying gauge anomalies through symmetry-protected trivial
  orders and classifying gravitational anomalies through topological orders},
  \href{https://doi.org/10.1103/PhysRevD.88.045013}{\emph{Phys. Rev. D}
  {\bfseries 88} (2013) 045013}.

\bibitem{string8}
L.~Kong and X.-G.~Wen, \emph{Braided fusion categories, gravitational
  anomalies, and the mathematical framework for topological orders in any
  dimensions},  2014.

\bibitem{Wen99}
X.-G.~Wen, \emph{Projective construction of non-abelian quantum hall liquids},
  \href{https://doi.org/10.1103/PhysRevB.60.8827}{\emph{Phys. Rev. B}
  {\bfseries 60} (1999) 8827}.

\bibitem{Barkeshli2010}
M.~Barkeshli and X.-G.~Wen, \emph{Effective field theory and projective
  construction for ${Z}_{k}$ parafermion fractional quantum hall states},
  \href{https://doi.org/10.1103/PhysRevB.81.155302}{\emph{Phys. Rev. B}
  {\bfseries 81} (2010) 155302}.

\bibitem{barkeshli_wen}
M.~Barkeshli and X.-G.~Wen,
  \emph{$u(1)\ifmmode\times\else\texttimes\fi{}u(1)\ensuremath{\rtimes}{Z}_{2}$
  chern-simons theory and ${Z}_{4}$ parafermion fractional quantum hall
  states}, \href{https://doi.org/10.1103/PhysRevB.81.045323}{\emph{Phys. Rev.
  B} {\bfseries 81} (2010) 045323}.

\bibitem{composite_Jain_2}
J.K.~Jain, \emph{Composite fermions}, Cambridge University Press (2007).

\bibitem{composite_Fradkin}
A.~Lopez and E.~Fradkin, \emph{Fractional quantum hall effect and chern-simons
  gauge theories}, \href{https://doi.org/10.1103/PhysRevB.44.5246}{\emph{Phys.
  Rev. B} {\bfseries 44} (1991) 5246}.

\bibitem{PhysRevLett.62.82}
S.C.~Zhang, T.H.~Hansson and S.~Kivelson, \emph{Effective-field-theory model
  for the fractional quantum hall effect},
  \href{https://doi.org/10.1103/PhysRevLett.62.82}{\emph{Phys. Rev. Lett.}
  {\bfseries 62} (1989) 82}.

\bibitem{PhysRevB.93.155121}
Y.-M.~Lu and A.~Vishwanath, \emph{Classification and properties of
  symmetry-enriched topological phases: Chern-simons approach with applications
  to ${Z}_{2}$ spin liquids},
  \href{https://doi.org/10.1103/PhysRevB.93.155121}{\emph{Phys. Rev. B}
  {\bfseries 93} (2016) 155121}.

\bibitem{PhysRevB.87.195103}
L.-Y.~Hung and Y.~Wan, \emph{$k$ matrix construction of symmetry-enriched
  phases of matter},
  \href{https://doi.org/10.1103/PhysRevB.87.195103}{\emph{Phys. Rev. B}
  {\bfseries 87} (2013) 195103}.

\bibitem{PhysRevB.86.125119}
Y.-M.~Lu and A.~Vishwanath, \emph{Theory and classification of interacting
  integer topological phases in two dimensions: A chern-simons approach},
  \href{https://doi.org/10.1103/PhysRevB.86.125119}{\emph{Phys. Rev. B}
  {\bfseries 86} (2012) 125119}.

\bibitem{YW12}
P.~Ye and X.-G.~Wen, \emph{Projective construction of two-dimensional
  symmetry-protected topological phases with u(1), so(3), or su(2) symmetries},
  \href{https://doi.org/10.1103/PhysRevB.87.195128}{\emph{Phys. Rev. B}
  {\bfseries 87} (2013) 195128}.

\bibitem{PhysRevB.93.115136}
Z.-C.~Gu, J.C.~Wang and X.-G.~Wen, \emph{Multikink topological terms and
  charge-binding domain-wall condensation induced symmetry-protected
  topological states: Beyond chern-simons/bf field theories},
  \href{https://doi.org/10.1103/PhysRevB.93.115136}{\emph{Phys. Rev. B}
  {\bfseries 93} (2016) 115136}.

\bibitem{Ye14b}
Z.-X.~Liu, J.-W.~Mei, P.~Ye and X.-G.~Wen,
  \emph{$u(1)\ifmmode\times\else\texttimes\fi{}u(1)$ symmetry-protected
  topological order in gutzwiller wave functions},
  \href{https://doi.org/10.1103/PhysRevB.90.235146}{\emph{Phys. Rev. B}
  {\bfseries 90} (2014) 235146}.

\bibitem{PhysRevLett.112.141602}
M.~Cheng and Z.-C.~Gu, \emph{Topological response theory of abelian
  symmetry-protected topological phases in two dimensions},
  \href{https://doi.org/10.1103/PhysRevLett.112.141602}{\emph{Phys. Rev. Lett.}
  {\bfseries 112} (2014) 141602}.

\bibitem{Hung_Wen_gauge}
L.-Y.~Hung and X.-G.~Wen, \emph{Quantized topological terms in weak-coupling
  gauge theories with a global symmetry and their connection to
  symmetry-enriched topological phases},
  \href{https://doi.org/10.1103/PhysRevB.87.165107}{\emph{Phys. Rev. B}
  {\bfseries 87} (2013) 165107}.

\bibitem{wu84}
Y.-S.~Wu, \emph{General theory for quantum statistics in two dimensions},
  \href{https://doi.org/10.1103/PhysRevLett.52.2103}{\emph{Phys. Rev. Lett.}
  {\bfseries 52} (1984) 2103}.

\bibitem{Horowitz1990}
G.T.~Horowitz and M.~Srednicki, \emph{A quantum field theoretic description of
  linking numbers and their generalization},
  \href{https://doi.org/10.1007/BF02099875}{\emph{Communications in
  Mathematical Physics} {\bfseries 130} (1990) 83}.

\bibitem{hansson2004superconductors}
T.H.~Hansson, V.~Oganesyan and S.L.~Sondhi, \emph{Superconductors are
  topologically ordered},
  \href{https://doi.org/http://dx.doi.org/10.1016/j.aop.2004.05.006}{\emph{Annals
  of Physics} {\bfseries 313} (2004) 497}.

\bibitem{YW13a}
P.~Ye and X.-G.~Wen, \emph{Constructing symmetric topological phases of bosons
  in three dimensions via fermionic projective construction and dyon
  condensation}, \href{https://doi.org/10.1103/PhysRevB.89.045127}{\emph{Phys.
  Rev. B} {\bfseries 89} (2014) 045127}.

\bibitem{ye16a}
P.~Ye, T.L.~Hughes, J.~Maciejko and E.~Fradkin, \emph{Composite particle theory
  of three-dimensional gapped fermionic phases: Fractional topological
  insulators and charge-loop excitation symmetry},
  \href{https://doi.org/10.1103/PhysRevB.94.115104}{\emph{Phys. Rev. B}
  {\bfseries 94} (2016) 115104}.

\bibitem{Moy_Fradkin2023}
B.~Moy, H.~Goldman, R.~Sohal and E.~Fradkin, \emph{{Theory of oblique
  topological insulators}},
  \href{https://doi.org/10.21468/SciPostPhys.14.2.023}{\emph{SciPost Phys.}
  {\bfseries 14} (2023) 023}.

\bibitem{2016arXiv161209298P}
P.~Putrov, J.~Wang and S.-T.~Yau, \emph{Braiding statistics and link invariants
  of bosonic/fermionic topological quantum matter in 2+1 and 3+1 dimensions},
  \href{https://doi.org/https://doi.org/10.1016/j.aop.2017.06.019}{\emph{Annals
  of Physics} {\bfseries 384} (2017) 254 }.

\bibitem{PhysRevB.99.235137}
Q.-R.~Wang, M.~Cheng, C.~Wang and Z.-C.~Gu, \emph{Topological quantum field
  theory for abelian topological phases and loop braiding statistics in
  $(3+1)$-dimensions},
  \href{https://doi.org/10.1103/PhysRevB.99.235137}{\emph{Phys. Rev. B}
  {\bfseries 99} (2019) 235137}.

\bibitem{YeGu2015}
P.~Ye and Z.-C.~Gu, \emph{Topological quantum field theory of three-dimensional
  bosonic abelian-symmetry-protected topological phases},
  \href{https://doi.org/10.1103/PhysRevB.93.205157}{\emph{Phys. Rev. B}
  {\bfseries 93} (2016) 205157}.

\bibitem{ypdw}
X.~Wen, H.~He, A.~Tiwari, Y.~Zheng and P.~Ye, \emph{Entanglement entropy for
  (3+1)-dimensional topological order with excitations},
  \href{https://doi.org/10.1103/PhysRevB.97.085147}{\emph{Phys. Rev. B}
  {\bfseries 97} (2018) 085147}.

\bibitem{yp18prl}
A.P.O.~Chan, P.~Ye and S.~Ryu, \emph{Braiding with borromean rings in
  ($3+1$)-dimensional spacetime},
  \href{https://doi.org/10.1103/PhysRevLett.121.061601}{\emph{Phys. Rev. Lett.}
  {\bfseries 121} (2018) 061601}.

\bibitem{abeffect}
Y.~Aharonov and D.~Bohm, \emph{Significance of electromagnetic potentials in
  the quantum theory},
  \href{https://doi.org/10.1103/PhysRev.115.485}{\emph{Phys. Rev.} {\bfseries
  115} (1959) 485}.

\bibitem{PRESKILL199050}
J.~Preskill and L.M.~Krauss, \emph{Local discrete symmetry and
  quantum-mechanical hair},
  \href{https://doi.org/https://doi.org/10.1016/0550-3213(90)90262-C}{\emph{Nuclear
  Physics B} {\bfseries 341} (1990) 50 }.

\bibitem{PhysRevLett.62.1071}
M.G.~Alford and F.~Wilczek, \emph{Aharonov-bohm interaction of cosmic strings
  with matter}, \href{https://doi.org/10.1103/PhysRevLett.62.1071}{\emph{Phys.
  Rev. Lett.} {\bfseries 62} (1989) 1071}.

\bibitem{PhysRevLett.62.1221}
L.M.~Krauss and F.~Wilczek, \emph{Discrete gauge symmetry in continuum
  theories}, \href{https://doi.org/10.1103/PhysRevLett.62.1221}{\emph{Phys.
  Rev. Lett.} {\bfseries 62} (1989) 1221}.

\bibitem{ALFORD1992251}
M.G.~Alford, K.-M.~Lee, J.~March-Russell and J.~Preskill, \emph{Quantum field
  theory of non-abelian strings and vortices},
  \href{https://doi.org/https://doi.org/10.1016/0550-3213(92)90468-Q}{\emph{Nuclear
  Physics B} {\bfseries 384} (1992) 251 }.

\bibitem{wang_levin1}
C.~Wang and M.~Levin, \emph{Braiding statistics of loop excitations in three
  dimensions},
  \href{https://doi.org/10.1103/PhysRevLett.113.080403}{\emph{Phys. Rev. Lett.}
  {\bfseries 113} (2014) 080403}.

\bibitem{kauffman2006formal}
L.H.~Kauffman, \emph{Formal knot theory}, Courier Corporation (2006).

\bibitem{mellor2003geometric}
B.~Mellor and P.~Melvin, \emph{A geometric interpretation of milnor's triple
  linking numbers}, \href{https://doi.org/10.2140/agt.2003.3.557}{\emph{Algebr.
  Geom. Topol.} {\bfseries 3} (2003) 557}.

\bibitem{zhang2021compatible}
Z.-F.~Zhang and P.~Ye, \emph{Compatible braidings with hopf links, multiloop,
  and borromean rings in $(3+1)$-dimensional spacetime},
  \href{https://doi.org/10.1103/PhysRevResearch.3.023132}{\emph{Phys. Rev.
  Research} {\bfseries 3} (2021) 023132}.

\bibitem{PhysRevB.107.165117}
Z.-F.~Zhang, Q.-R.~Wang and P.~Ye, \emph{Non-abelian fusion, shrinking, and
  quantum dimensions of abelian gauge fluxes},
  \href{https://doi.org/10.1103/PhysRevB.107.165117}{\emph{Phys. Rev. B}
  {\bfseries 107} (2023) 165117}.

\bibitem{ye16_set}
S.-Q.~Ning, Z.-X.~Liu and P.~Ye, \emph{Symmetry enrichment in three-dimensional
  topological phases},
  \href{https://doi.org/10.1103/PhysRevB.94.245120}{\emph{Phys. Rev. B}
  {\bfseries 94} (2016) 245120}.

\bibitem{2016arXiv161008645Y}
P.~Ye, \emph{Three-dimensional anomalous twisted gauge theories with global
  symmetry: Implications for quantum spin liquids},
  \href{https://doi.org/10.1103/PhysRevB.97.125127}{\emph{Phys. Rev. B}
  {\bfseries 97} (2018) 125127}.

\bibitem{Ye:2017aa}
P.~Ye, M.~Cheng and E.~Fradkin, \emph{Fractional $s$-duality, classification of
  fractional topological insulators, and surface topological order},
  \href{https://doi.org/10.1103/PhysRevB.96.085125}{\emph{Phys. Rev. B}
  {\bfseries 96} (2017) 085125}.

\bibitem{Ning2018prb}
S.-Q.~Ning, Z.-X.~Liu and P.~Ye, \emph{Fractionalizing global symmetry on
  looplike topological excitations},
  \href{https://doi.org/10.1103/PhysRevB.105.205137}{\emph{Phys. Rev. B}
  {\bfseries 105} (2022) 205137}.

\bibitem{bti2}
P.~Ye and Z.-C.~Gu, \emph{Vortex-line condensation in three dimensions: A
  physical mechanism for bosonic topological insulators},
  \href{https://doi.org/10.1103/PhysRevX.5.021029}{\emph{Phys. Rev. X}
  {\bfseries 5} (2015) 021029}.

\bibitem{PhysRevLett.114.031601}
J.C.~Wang, Z.-C.~Gu and X.-G.~Wen, \emph{Field-theory representation of
  gauge-gravity symmetry-protected topological invariants, group cohomology,
  and beyond},
  \href{https://doi.org/10.1103/PhysRevLett.114.031601}{\emph{Phys. Rev. Lett.}
  {\bfseries 114} (2015) 031601}.

\bibitem{PhysRevB.99.205120}
B.~Han, H.~Wang and P.~Ye, \emph{Generalized wen-zee terms},
  \href{https://doi.org/10.1103/PhysRevB.99.205120}{\emph{Phys. Rev. B}
  {\bfseries 99} (2019) 205120}.

\bibitem{bti6}
P.~Ye and J.~Wang, \emph{Symmetry-protected topological phases with charge and
  spin symmetries: Response theory and dynamical gauge theory in two and three
  dimensions}, \href{https://doi.org/10.1103/PhysRevB.88.235109}{\emph{Phys.
  Rev. B} {\bfseries 88} (2013) 235109}.

\bibitem{Kapustin:2014gua}
A.~Kapustin and N.~Seiberg, \emph{{Coupling a QFT to a TQFT and Duality}},
  \href{https://doi.org/10.1007/JHEP04(2014)001}{\emph{JHEP} {\bfseries 04}
  (2014) 001} [\href{https://arxiv.org/abs/1401.0740}{{\ttfamily 1401.0740}}].

\bibitem{ye2023unp}
Z.-F.~Zhang, Q.-R.~Wang and P.~Ye{\emph{Unpublished} (2023) }.

\bibitem{lapa17}
M.F.~Lapa, C.-M.~Jian, P.~Ye and T.L.~Hughes, \emph{Topological electromagnetic
  responses of bosonic quantum hall, topological insulator, and chiral
  semimetal phases in all dimensions},
  \href{https://doi.org/10.1103/PhysRevB.95.035149}{\emph{Phys. Rev. B}
  {\bfseries 95} (2017) 035149}.

\bibitem{Zhang:2021ycl}
Z.-F.~Zhang and P.~Ye, \emph{{Topological orders, braiding statistics, and
  mixture of two types of twisted BF theories in five dimensions}},
  \href{https://doi.org/10.1007/JHEP04(2022)138}{\emph{JHEP} {\bfseries 04}
  (2022) 138} [\href{https://arxiv.org/abs/2104.07067}{{\ttfamily
  2104.07067}}].

\bibitem{mcgreevy2023}
J.~McGreevy, \emph{Generalized symmetries in condensed matter},
  \href{https://doi.org/10.1146/annurev-conmatphys-040721-021029}{\emph{Annual
  Review of Condensed Matter Physics} {\bfseries 14} (2023) 57}
  [\href{https://arxiv.org/abs/https://doi.org/10.1146/annurev-conmatphys-040721-021029}{{\ttfamily
  https://doi.org/10.1146/annurev-conmatphys-040721-021029}}].

\end{thebibliography}


\end{document}
%
% ****** End of file apstemplate.tex ******

